% Generated by `python -m scripts.reporting.test_cases` on 2026-11-18 (commit e743616).
% One table per test case, in the layout of the report's Table 7.1.
% Needs \usepackage{array,longtable} in the preamble.
\newcolumntype{L}[1]{>{\raggedright\arraybackslash}p{#1}}
\newcolumntype{C}[1]{>{\centering\arraybackslash}p{#1}}

\begin{longtable}{|L{0.21\linewidth}|L{0.27\linewidth}|L{0.17\linewidth}|L{0.19\linewidth}|}
\caption{Test case TC-01: Late aggregation settings}\label{tab:tc-01} \\ \hline
\multicolumn{4}{|C{\dimexpr 0.84\linewidth+6\tabcolsep+3\arrayrulewidth\relax}|}{\textbf{Late aggregation settings}} \\ \hline
\multicolumn{4}{|C{\dimexpr 0.84\linewidth+6\tabcolsep+3\arrayrulewidth\relax}|}{\textbf{src/\allowbreak{}experts/\allowbreak{}levels.py, scripts/\allowbreak{}train\_\allowbreak{}models.py, src/\allowbreak{}experts/\allowbreak{}expert.py; tests/\allowbreak{}test\_\allowbreak{}aggregation\_\allowbreak{}modes.py (TestLateSettings)}} \\ \hline
\textbf{Test Case ID:} & TC-01 & \textbf{QA Test Engineer:} & Abi \\ \hline
\textbf{Test case Version:} & 1 & \textbf{Reviewed By:} & - \\ \hline
\textbf{Test Date:} & 2026-11-18 & \textbf{Use Case Reference(s):} & Design \S{}8.1, \S{}8.2; config engagement.aggregate \\ \hline
\textbf{Revision History:} & \multicolumn{3}{L{\dimexpr 0.63\linewidth+4\tabcolsep+2\arrayrulewidth\relax}|}{v1 2026-11-15 (db92029)} \\ \hline
\textbf{Objective:} & \multicolumn{3}{L{\dimexpr 0.63\linewidth+4\tabcolsep+2\arrayrulewidth\relax}|}{Check the decision threshold for each aggregation mode, the mapping from window-mean feature names to per-student names, and that the hand-weighted baseline reads per-student columns in late mode.} \\ \hline
\textbf{Product/Ver/ Module:} & \multicolumn{3}{L{\dimexpr 0.63\linewidth+4\tabcolsep+2\arrayrulewidth\relax}|}{FYP-1 attention pipeline, commit e743616. Engagement stage: settings that differ between early (one row per window) and late (one row per student per window) aggregation.} \\ \hline
\textbf{Environment:} & \multicolumn{3}{L{\dimexpr 0.63\linewidth+4\tabcolsep+2\arrayrulewidth\relax}|}{Windows 11 (10.0.26200); conda environment fyp-local: Python 3.13.18, NumPy 2.5.4, SciPy 1.18.1, scikit-learn 1.9.1, pandas 3.1.1, OpenCV 4.11.0, PyTorch 2.9.0+cu126. CPU only; no GPU, dataset or model weights needed unless stated. Run from the repository root: python -m unittest discover -s tests -t .} \\ \hline
\textbf{Assumptions:} & \multicolumn{3}{L{\dimexpr 0.63\linewidth+4\tabcolsep+2\arrayrulewidth\relax}|}{The hand-weighted weights come from configs/\allowbreak{}deliverable3.json, which names window-level keys such as vec.kg.stillness.mean.} \\ \hline
\textbf{Pre-Requisite:} & \multicolumn{3}{L{\dimexpr 0.63\linewidth+4\tabcolsep+2\arrayrulewidth\relax}|}{configs/\allowbreak{}deliverable3.json must be present and loadable.} \\ \hline
\textbf{Step No.} & \textbf{Execution description} & \multicolumn{2}{L{\dimexpr 0.36\linewidth+2\tabcolsep+1\arrayrulewidth\relax}|}{\textbf{Procedure result}} \\ \hline
1 & \texttt{test\_\allowbreak{}late\_\allowbreak{}scores\_\allowbreak{}are\_\allowbreak{}cut\_\allowbreak{}at\_\allowbreak{}the\_\allowbreak{}label\_\allowbreak{}share}: Call decision\_\allowbreak{}threshold with mixture\_\allowbreak{}tau 0.5 and tau\_\allowbreak{}share 0.6 for 'late' and for 'early'. & \multicolumn{2}{L{\dimexpr 0.36\linewidth+2\tabcolsep+1\arrayrulewidth\relax}|}{Late returns 0.6 (tau\_\allowbreak{}share). Early returns 0.5 (mixture\_\allowbreak{}tau). (Passed)} \\ \hline
2 & \texttt{test\_\allowbreak{}window\_\allowbreak{}mean\_\allowbreak{}key\_\allowbreak{}maps\_\allowbreak{}to\_\allowbreak{}the\_\allowbreak{}student\_\allowbreak{}feature}: Call student\_\allowbreak{}key on 'vec.kg.stillness.mean', 'vec.kg.stillness.std' and 'vec.n\_\allowbreak{}tracks'. & \multicolumn{2}{L{\dimexpr 0.36\linewidth+2\tabcolsep+1\arrayrulewidth\relax}|}{Returns 'trk.kg.stillness', 'vec.kg.stillness.std' and 'vec.n\_\allowbreak{}tracks'. Only .mean keys are mapped. (Passed)} \\ \hline
3 & \texttt{test\_\allowbreak{}hand\_\allowbreak{}weighted\_\allowbreak{}baseline\_\allowbreak{}reads\_\allowbreak{}students\_\allowbreak{}in\_\allowbreak{}late\_\allowbreak{}mode}: Load configs/\allowbreak{}deliverable3.json, set engagement to \{aggregate: 'late', late\_\allowbreak{}context: false\}, and build a hand\_\allowbreak{}weighted expert for experiment '02' with columns ['trk.kg.stillness', 'trk.scb.write']. & \multicolumn{2}{L{\dimexpr 0.36\linewidth+2\tabcolsep+1\arrayrulewidth\relax}|}{The expert's selected column indices cover both columns, so the window-level weights were mapped to the per-student features. (Passed)} \\ \hline
\multicolumn{4}{|L{\dimexpr 0.84\linewidth+6\tabcolsep+3\arrayrulewidth\relax}|}{\textbf{Comments:} -} \\ \hline
\multicolumn{4}{|C{\dimexpr 0.84\linewidth+6\tabcolsep+3\arrayrulewidth\relax}|}{\textbf{\fbox{Passed} \quad Failed \quad Not Executed}} \\ \hline
\end{longtable}

\begin{longtable}{|L{0.21\linewidth}|L{0.27\linewidth}|L{0.17\linewidth}|L{0.19\linewidth}|}
\caption{Test case TC-02: Late inference with no visible student}\label{tab:tc-02} \\ \hline
\multicolumn{4}{|C{\dimexpr 0.84\linewidth+6\tabcolsep+3\arrayrulewidth\relax}|}{\textbf{Late inference with no visible student}} \\ \hline
\multicolumn{4}{|C{\dimexpr 0.84\linewidth+6\tabcolsep+3\arrayrulewidth\relax}|}{\textbf{src/\allowbreak{}experts/\allowbreak{}levels.py; tests/\allowbreak{}test\_\allowbreak{}aggregation\_\allowbreak{}modes.py (TestNobodyVisible)}} \\ \hline
\textbf{Test Case ID:} & TC-02 & \textbf{QA Test Engineer:} & Abi \\ \hline
\textbf{Test case Version:} & 1 & \textbf{Reviewed By:} & - \\ \hline
\textbf{Test Date:} & 2026-11-18 & \textbf{Use Case Reference(s):} & Design \S{}8 \\ \hline
\textbf{Revision History:} & \multicolumn{3}{L{\dimexpr 0.63\linewidth+4\tabcolsep+2\arrayrulewidth\relax}|}{v1 2026-11-15 (db92029)} \\ \hline
\textbf{Objective:} & \multicolumn{3}{L{\dimexpr 0.63\linewidth+4\tabcolsep+2\arrayrulewidth\relax}|}{Check that late aggregation handles a session with no visible tracks: no rows are built, the model is not called, and every window gets the fill value.} \\ \hline
\textbf{Product/Ver/ Module:} & \multicolumn{3}{L{\dimexpr 0.63\linewidth+4\tabcolsep+2\arrayrulewidth\relax}|}{FYP-1 attention pipeline, commit e743616. Engagement stage at inference, late aggregation.} \\ \hline
\textbf{Environment:} & \multicolumn{3}{L{\dimexpr 0.63\linewidth+4\tabcolsep+2\arrayrulewidth\relax}|}{Windows 11 (10.0.26200); conda environment fyp-local: Python 3.13.18, NumPy 2.5.4, SciPy 1.18.1, scikit-learn 1.9.1, pandas 3.1.1, OpenCV 4.11.0, PyTorch 2.9.0+cu126. CPU only; no GPU, dataset or model weights needed unless stated. Run from the repository root: python -m unittest discover -s tests -t .} \\ \hline
\textbf{Assumptions:} & \multicolumn{3}{L{\dimexpr 0.63\linewidth+4\tabcolsep+2\arrayrulewidth\relax}|}{Two windows whose only feature vec.a.mean is NaN, and an empty track table with columns window and tid (what to\_\allowbreak{}track\_\allowbreak{}frame returns with no tracks).} \\ \hline
\textbf{Pre-Requisite:} & \multicolumn{3}{L{\dimexpr 0.63\linewidth+4\tabcolsep+2\arrayrulewidth\relax}|}{None beyond the Python environment.} \\ \hline
\textbf{Step No.} & \textbf{Execution description} & \multicolumn{2}{L{\dimexpr 0.36\linewidth+2\tabcolsep+1\arrayrulewidth\relax}|}{\textbf{Procedure result}} \\ \hline
1 & \texttt{test\_\allowbreak{}late\_\allowbreak{}inference\_\allowbreak{}without\_\allowbreak{}any\_\allowbreak{}student\_\allowbreak{}falls\_\allowbreak{}back\_\allowbreak{}to\_\allowbreak{}fill}: Call late\_\allowbreak{}rows with context off and fixed columns ['trk.a', 'vec.a.mean'], then predict\_\allowbreak{}windows for windows [0, 1] with a stub model whose predict\_\allowbreak{}proba raises an AssertionError, gate 0.5 and fill 0.3. & \multicolumn{2}{L{\dimexpr 0.36\linewidth+2\tabcolsep+1\arrayrulewidth\relax}|}{late\_\allowbreak{}rows returns X with shape (0, 2). predict\_\allowbreak{}windows returns [0.3, 0.3] without calling the model. (Passed)} \\ \hline
\multicolumn{4}{|L{\dimexpr 0.84\linewidth+6\tabcolsep+3\arrayrulewidth\relax}|}{\textbf{Comments:} -} \\ \hline
\multicolumn{4}{|C{\dimexpr 0.84\linewidth+6\tabcolsep+3\arrayrulewidth\relax}|}{\textbf{\fbox{Passed} \quad Failed \quad Not Executed}} \\ \hline
\end{longtable}

\begin{longtable}{|L{0.21\linewidth}|L{0.27\linewidth}|L{0.17\linewidth}|L{0.19\linewidth}|}
\caption{Test case TC-03: Late row construction and window scores}\label{tab:tc-03} \\ \hline
\multicolumn{4}{|C{\dimexpr 0.84\linewidth+6\tabcolsep+3\arrayrulewidth\relax}|}{\textbf{Late row construction and window scores}} \\ \hline
\multicolumn{4}{|C{\dimexpr 0.84\linewidth+6\tabcolsep+3\arrayrulewidth\relax}|}{\textbf{src/\allowbreak{}experts/\allowbreak{}levels.py; tests/\allowbreak{}test\_\allowbreak{}aggregation\_\allowbreak{}modes.py (TestRows)}} \\ \hline
\textbf{Test Case ID:} & TC-03 & \textbf{QA Test Engineer:} & Abi \\ \hline
\textbf{Test case Version:} & 1 & \textbf{Reviewed By:} & - \\ \hline
\textbf{Test Date:} & 2026-11-18 & \textbf{Use Case Reference(s):} & Design \S{}8 \\ \hline
\textbf{Revision History:} & \multicolumn{3}{L{\dimexpr 0.63\linewidth+4\tabcolsep+2\arrayrulewidth\relax}|}{v1 2026-11-15 (db92029)} \\ \hline
\textbf{Objective:} & \multicolumn{3}{L{\dimexpr 0.63\linewidth+4\tabcolsep+2\arrayrulewidth\relax}|}{Check that late rows carry each track's features plus optional window context, that fixed column lists fill missing columns with NaN, and that window scores average rows and fill empty windows.} \\ \hline
\textbf{Product/Ver/ Module:} & \multicolumn{3}{L{\dimexpr 0.63\linewidth+4\tabcolsep+2\arrayrulewidth\relax}|}{FYP-1 attention pipeline, commit e743616. Engagement stage: builds expert inputs per (window, track) and turns row probabilities into window scores.} \\ \hline
\textbf{Environment:} & \multicolumn{3}{L{\dimexpr 0.63\linewidth+4\tabcolsep+2\arrayrulewidth\relax}|}{Windows 11 (10.0.26200); conda environment fyp-local: Python 3.13.18, NumPy 2.5.4, SciPy 1.18.1, scikit-learn 1.9.1, pandas 3.1.1, OpenCV 4.11.0, PyTorch 2.9.0+cu126. CPU only; no GPU, dataset or model weights needed unless stated. Run from the repository root: python -m unittest discover -s tests -t .} \\ \hline
\textbf{Assumptions:} & \multicolumn{3}{L{\dimexpr 0.63\linewidth+4\tabcolsep+2\arrayrulewidth\relax}|}{Three windows with vec.a.mean [1, 2, 3] and vec.n\_\allowbreak{}tracks [2, 0, 1]; tracks in windows [0, 0, 2] with tids [7, 8, 7] and trk.a [0.5, 1.5, 3.0].} \\ \hline
\textbf{Pre-Requisite:} & \multicolumn{3}{L{\dimexpr 0.63\linewidth+4\tabcolsep+2\arrayrulewidth\relax}|}{None beyond the Python environment.} \\ \hline
\textbf{Step No.} & \textbf{Execution description} & \multicolumn{2}{L{\dimexpr 0.36\linewidth+2\tabcolsep+1\arrayrulewidth\relax}|}{\textbf{Procedure result}} \\ \hline
1 & \texttt{test\_\allowbreak{}late\_\allowbreak{}rows\_\allowbreak{}carry\_\allowbreak{}window\_\allowbreak{}context}: Call late\_\allowbreak{}rows with context=True. & \multicolumn{2}{L{\dimexpr 0.36\linewidth+2\tabcolsep+1\arrayrulewidth\relax}|}{Columns are ['trk.a', 'vec.a.mean', 'vec.n\_\allowbreak{}tracks'], window indices are [0, 0, 2], and X is [[0.5, 1, 2], [1.5, 1, 2], [3, 3, 1]]. (Passed)} \\ \hline
2 & \texttt{test\_\allowbreak{}late\_\allowbreak{}rows\_\allowbreak{}without\_\allowbreak{}context\_\allowbreak{}and\_\allowbreak{}fixed\_\allowbreak{}columns}: Call late\_\allowbreak{}rows with context=False, then again with fixed columns ['trk.a', 'trk.missing']. & \multicolumn{2}{L{\dimexpr 0.36\linewidth+2\tabcolsep+1\arrayrulewidth\relax}|}{The first call has columns ['trk.a']. In the second call the 'trk.missing' column is all NaN. (Passed)} \\ \hline
3 & \texttt{test\_\allowbreak{}window\_\allowbreak{}scores\_\allowbreak{}mean\_\allowbreak{}rows\_\allowbreak{}and\_\allowbreak{}fill\_\allowbreak{}empty\_\allowbreak{}windows}: Call window\_\allowbreak{}scores with row probabilities [0.2, 0.6, 0.9], row windows [0, 0, 2], 3 windows and fill 0.4. & \multicolumn{2}{L{\dimexpr 0.36\linewidth+2\tabcolsep+1\arrayrulewidth\relax}|}{Returns [0.4, 0.4, 0.9]: window 0 is the mean of 0.2 and 0.6, window 1 has no rows and gets the fill. (Passed)} \\ \hline
\multicolumn{4}{|L{\dimexpr 0.84\linewidth+6\tabcolsep+3\arrayrulewidth\relax}|}{\textbf{Comments:} -} \\ \hline
\multicolumn{4}{|C{\dimexpr 0.84\linewidth+6\tabcolsep+3\arrayrulewidth\relax}|}{\textbf{\fbox{Passed} \quad Failed \quad Not Executed}} \\ \hline
\end{longtable}

\begin{longtable}{|L{0.21\linewidth}|L{0.27\linewidth}|L{0.17\linewidth}|L{0.19\linewidth}|}
\caption{Test case TC-04: Early aggregation equivalence}\label{tab:tc-04} \\ \hline
\multicolumn{4}{|C{\dimexpr 0.84\linewidth+6\tabcolsep+3\arrayrulewidth\relax}|}{\textbf{Early aggregation equivalence}} \\ \hline
\multicolumn{4}{|C{\dimexpr 0.84\linewidth+6\tabcolsep+3\arrayrulewidth\relax}|}{\textbf{src/\allowbreak{}experts/\allowbreak{}levels.py, src/\allowbreak{}experts/\allowbreak{}mixture.py, src/\allowbreak{}experts/\allowbreak{}expert.py; tests/\allowbreak{}test\_\allowbreak{}aggregation\_\allowbreak{}modes.py (TestEarlyEqualsDirectMixture)}} \\ \hline
\textbf{Test Case ID:} & TC-04 & \textbf{QA Test Engineer:} & Abi \\ \hline
\textbf{Test case Version:} & 1 & \textbf{Reviewed By:} & - \\ \hline
\textbf{Test Date:} & 2026-11-18 & \textbf{Use Case Reference(s):} & Design \S{}6, \S{}8 \\ \hline
\textbf{Revision History:} & \multicolumn{3}{L{\dimexpr 0.63\linewidth+4\tabcolsep+2\arrayrulewidth\relax}|}{v1 2026-11-15 (db92029)} \\ \hline
\textbf{Objective:} & \multicolumn{3}{L{\dimexpr 0.63\linewidth+4\tabcolsep+2\arrayrulewidth\relax}|}{Check that the row-based fitting path in early mode gives the same predictions as fitting the window-level mixture directly.} \\ \hline
\textbf{Product/Ver/ Module:} & \multicolumn{3}{L{\dimexpr 0.63\linewidth+4\tabcolsep+2\arrayrulewidth\relax}|}{FYP-1 attention pipeline, commit e743616. Engagement stage, early aggregation (one row per window).} \\ \hline
\textbf{Environment:} & \multicolumn{3}{L{\dimexpr 0.63\linewidth+4\tabcolsep+2\arrayrulewidth\relax}|}{Windows 11 (10.0.26200); conda environment fyp-local: Python 3.13.18, NumPy 2.5.4, SciPy 1.18.1, scikit-learn 1.9.1, pandas 3.1.1, OpenCV 4.11.0, PyTorch 2.9.0+cu126. CPU only; no GPU, dataset or model weights needed unless stated. Run from the repository root: python -m unittest discover -s tests -t .} \\ \hline
\textbf{Assumptions:} & \multicolumn{3}{L{\dimexpr 0.63\linewidth+4\tabcolsep+2\arrayrulewidth\relax}|}{60 synthetic windows (seed 0) with features vec.x and vec.z, experiments alternating '01' and '02', and label y = (vec.x + 0.3 * noise > 0). Per-experiment soft mixture of logistic regression with seed 0. Gate probabilities [0.3, 0.7]. Windows 0-39 train, 40-59 test.} \\ \hline
\textbf{Pre-Requisite:} & \multicolumn{3}{L{\dimexpr 0.63\linewidth+4\tabcolsep+2\arrayrulewidth\relax}|}{None beyond the Python environment.} \\ \hline
\textbf{Step No.} & \textbf{Execution description} & \multicolumn{2}{L{\dimexpr 0.36\linewidth+2\tabcolsep+1\arrayrulewidth\relax}|}{\textbf{Procedure result}} \\ \hline
1 & \texttt{test\_\allowbreak{}early\_\allowbreak{}rows\_\allowbreak{}reproduce\_\allowbreak{}the\_\allowbreak{}window\_\allowbreak{}level\_\allowbreak{}mixture}: Build rows with early\_\allowbreak{}rows and call fit\_\allowbreak{}predict. Separately fit a new mixture on the training windows and predict the test windows. & \multicolumn{2}{L{\dimexpr 0.36\linewidth+2\tabcolsep+1\arrayrulewidth\relax}|}{The two prediction arrays are equal (assert\_\allowbreak{}allclose). (Passed)} \\ \hline
\multicolumn{4}{|L{\dimexpr 0.84\linewidth+6\tabcolsep+3\arrayrulewidth\relax}|}{\textbf{Comments:} -} \\ \hline
\multicolumn{4}{|C{\dimexpr 0.84\linewidth+6\tabcolsep+3\arrayrulewidth\relax}|}{\textbf{\fbox{Passed} \quad Failed \quad Not Executed}} \\ \hline
\end{longtable}

\begin{longtable}{|L{0.21\linewidth}|L{0.27\linewidth}|L{0.17\linewidth}|L{0.19\linewidth}|}
\caption{Test case TC-05: Per-student labels for late aggregation}\label{tab:tc-05} \\ \hline
\multicolumn{4}{|C{\dimexpr 0.84\linewidth+6\tabcolsep+3\arrayrulewidth\relax}|}{\textbf{Per-student labels for late aggregation}} \\ \hline
\multicolumn{4}{|C{\dimexpr 0.84\linewidth+6\tabcolsep+3\arrayrulewidth\relax}|}{\textbf{src/\allowbreak{}labels/\allowbreak{}window\_\allowbreak{}labels.py, src/\allowbreak{}labels/\allowbreak{}attention.py; tests/\allowbreak{}test\_\allowbreak{}aggregation\_\allowbreak{}modes.py (TestStudentLabels)}} \\ \hline
\textbf{Test Case ID:} & TC-05 & \textbf{QA Test Engineer:} & Abi \\ \hline
\textbf{Test case Version:} & 1 & \textbf{Reviewed By:} & - \\ \hline
\textbf{Test Date:} & 2026-11-18 & \textbf{Use Case Reference(s):} & Design \S{}8.1 \\ \hline
\textbf{Revision History:} & \multicolumn{3}{L{\dimexpr 0.63\linewidth+4\tabcolsep+2\arrayrulewidth\relax}|}{v1 2026-11-15 (db92029)} \\ \hline
\textbf{Objective:} & \multicolumn{3}{L{\dimexpr 0.63\linewidth+4\tabcolsep+2\arrayrulewidth\relax}|}{Check that per-student labels use the same rule as the window label, so their mean equals the window's share engaged.} \\ \hline
\textbf{Product/Ver/ Module:} & \multicolumn{3}{L{\dimexpr 0.63\linewidth+4\tabcolsep+2\arrayrulewidth\relax}|}{FYP-1 attention pipeline, commit e743616. Label side of late aggregation: per-student training targets.} \\ \hline
\textbf{Environment:} & \multicolumn{3}{L{\dimexpr 0.63\linewidth+4\tabcolsep+2\arrayrulewidth\relax}|}{Windows 11 (10.0.26200); conda environment fyp-local: Python 3.13.18, NumPy 2.5.4, SciPy 1.18.1, scikit-learn 1.9.1, pandas 3.1.1, OpenCV 4.11.0, PyTorch 2.9.0+cu126. CPU only; no GPU, dataset or model weights needed unless stated. Run from the repository root: python -m unittest discover -s tests -t .} \\ \hline
\textbf{Assumptions:} & \multicolumn{3}{L{\dimexpr 0.63\linewidth+4\tabcolsep+2\arrayrulewidth\relax}|}{Three synthetic students, each with one labeler giving a constant rating from 10:00:00 (4, 2 and 5).} \\ \hline
\textbf{Pre-Requisite:} & \multicolumn{3}{L{\dimexpr 0.63\linewidth+4\tabcolsep+2\arrayrulewidth\relax}|}{None beyond the Python environment.} \\ \hline
\textbf{Step No.} & \textbf{Execution description} & \multicolumn{2}{L{\dimexpr 0.36\linewidth+2\tabcolsep+1\arrayrulewidth\relax}|}{\textbf{Procedure result}} \\ \hline
1 & \texttt{test\_\allowbreak{}student\_\allowbreak{}labels\_\allowbreak{}match\_\allowbreak{}the\_\allowbreak{}window\_\allowbreak{}share}: With WindowLabeler(tau\_\allowbreak{}att=4, tau\_\allowbreak{}share=0.5, min\_\allowbreak{}labeled\_\allowbreak{}students=1), call student\_\allowbreak{}labels and label on the 8 s window from 10:00:00. & \multicolumn{2}{L{\dimexpr 0.36\linewidth+2\tabcolsep+1\arrayrulewidth\relax}|}{student\_\allowbreak{}labels returns \{s1: 1.0, s2: 0.0, s3: 1.0\}. The window's share\_\allowbreak{}engaged equals their mean (2/\allowbreak{}3). (Passed)} \\ \hline
\multicolumn{4}{|L{\dimexpr 0.84\linewidth+6\tabcolsep+3\arrayrulewidth\relax}|}{\textbf{Comments:} -} \\ \hline
\multicolumn{4}{|C{\dimexpr 0.84\linewidth+6\tabcolsep+3\arrayrulewidth\relax}|}{\textbf{\fbox{Passed} \quad Failed \quad Not Executed}} \\ \hline
\end{longtable}

\begin{longtable}{|L{0.21\linewidth}|L{0.27\linewidth}|L{0.17\linewidth}|L{0.19\linewidth}|}
\caption{Test case TC-06: Feature pass per-track rows}\label{tab:tc-06} \\ \hline
\multicolumn{4}{|C{\dimexpr 0.84\linewidth+6\tabcolsep+3\arrayrulewidth\relax}|}{\textbf{Feature pass per-track rows}} \\ \hline
\multicolumn{4}{|C{\dimexpr 0.84\linewidth+6\tabcolsep+3\arrayrulewidth\relax}|}{\textbf{src/\allowbreak{}core/\allowbreak{}pipeline.py, src/\allowbreak{}core/\allowbreak{}types.py; tests/\allowbreak{}test\_\allowbreak{}aggregation\_\allowbreak{}modes.py (TestFeaturePassRows)}} \\ \hline
\textbf{Test Case ID:} & TC-06 & \textbf{QA Test Engineer:} & Abi \\ \hline
\textbf{Test case Version:} & 1 & \textbf{Reviewed By:} & - \\ \hline
\textbf{Test Date:} & 2026-11-18 & \textbf{Use Case Reference(s):} & Design \S{}4, \S{}6 \\ \hline
\textbf{Revision History:} & \multicolumn{3}{L{\dimexpr 0.63\linewidth+4\tabcolsep+2\arrayrulewidth\relax}|}{v1 2026-11-15 (db92029)} \\ \hline
\textbf{Objective:} & \multicolumn{3}{L{\dimexpr 0.63\linewidth+4\tabcolsep+2\arrayrulewidth\relax}|}{Check that the feature pass emits one per-track row for each visible track in each window, consistent with the window's track count, and that the track frame has per-track feature columns.} \\ \hline
\textbf{Product/Ver/ Module:} & \multicolumn{3}{L{\dimexpr 0.63\linewidth+4\tabcolsep+2\arrayrulewidth\relax}|}{FYP-1 attention pipeline, commit e743616. Feature pass of the pipeline (after tracking, before the engagement experts).} \\ \hline
\textbf{Environment:} & \multicolumn{3}{L{\dimexpr 0.63\linewidth+4\tabcolsep+2\arrayrulewidth\relax}|}{Windows 11 (10.0.26200); conda environment fyp-local: Python 3.13.18, NumPy 2.5.4, SciPy 1.18.1, scikit-learn 1.9.1, pandas 3.1.1, OpenCV 4.11.0, PyTorch 2.9.0+cu126. CPU only; no GPU, dataset or model weights needed unless stated. Run from the repository root: python -m unittest discover -s tests -t .} \\ \hline
\textbf{Assumptions:} & \multicolumn{3}{L{\dimexpr 0.63\linewidth+4\tabcolsep+2\arrayrulewidth\relax}|}{A synthetic SessionTrack with 100 frames at 5 fps (0 to 19.8 s), frame size 1280 x 720, and three fixed boxes. Tracks 0 and 1 are visible for all 20 s; track 2 only for the first 8 s (40 frames). Scene rows have motion 0. Settings come from configs/\allowbreak{}deliverable3.json.} \\ \hline
\textbf{Pre-Requisite:} & \multicolumn{3}{L{\dimexpr 0.63\linewidth+4\tabcolsep+2\arrayrulewidth\relax}|}{configs/\allowbreak{}deliverable3.json must be present. No video, detector or GPU is needed: the detector is built only on first use.} \\ \hline
\textbf{Step No.} & \textbf{Execution description} & \multicolumn{2}{L{\dimexpr 0.36\linewidth+2\tabcolsep+1\arrayrulewidth\relax}|}{\textbf{Procedure result}} \\ \hline
1 & \texttt{test\_\allowbreak{}one\_\allowbreak{}row\_\allowbreak{}per\_\allowbreak{}visible\_\allowbreak{}track\_\allowbreak{}per\_\allowbreak{}window}: Run Pipeline(cfg).features on the synthetic session, count track\_\allowbreak{}rows per window, and convert with to\_\allowbreak{}track\_\allowbreak{}frame. & \multicolumn{2}{L{\dimexpr 0.36\linewidth+2\tabcolsep+1\arrayrulewidth\relax}|}{The row count per window equals each window's n\_\allowbreak{}tracks in the Vec. The track frame has a 'trk.kg.speed' column and its tids are \{0, 1, 2\}. (Passed)} \\ \hline
\multicolumn{4}{|L{\dimexpr 0.84\linewidth+6\tabcolsep+3\arrayrulewidth\relax}|}{\textbf{Comments:} -} \\ \hline
\multicolumn{4}{|C{\dimexpr 0.84\linewidth+6\tabcolsep+3\arrayrulewidth\relax}|}{\textbf{\fbox{Passed} \quad Failed \quad Not Executed}} \\ \hline
\end{longtable}

\begin{longtable}{|L{0.21\linewidth}|L{0.27\linewidth}|L{0.17\linewidth}|L{0.19\linewidth}|}
\caption{Test case TC-07: Session-level data augmentation}\label{tab:tc-07} \\ \hline
\multicolumn{4}{|C{\dimexpr 0.84\linewidth+6\tabcolsep+3\arrayrulewidth\relax}|}{\textbf{Session-level data augmentation}} \\ \hline
\multicolumn{4}{|C{\dimexpr 0.84\linewidth+6\tabcolsep+3\arrayrulewidth\relax}|}{\textbf{src/\allowbreak{}ingestion/\allowbreak{}augment.py, src/\allowbreak{}ingestion/\allowbreak{}source.py, scripts/\allowbreak{}train\_\allowbreak{}models.py, src/\allowbreak{}spu/\allowbreak{}splits.py; tests/\allowbreak{}test\_\allowbreak{}augment.py (TestAugmentation)}} \\ \hline
\textbf{Test Case ID:} & TC-07 & \textbf{QA Test Engineer:} & Abi \\ \hline
\textbf{Test case Version:} & 3 & \textbf{Reviewed By:} & - \\ \hline
\textbf{Test Date:} & 2026-11-18 & \textbf{Use Case Reference(s):} & - \\ \hline
\textbf{Revision History:} & \multicolumn{3}{L{\dimexpr 0.63\linewidth+4\tabcolsep+2\arrayrulewidth\relax}|}{v1 2026-11-15 (a497de6); v2 2026-11-15 (be0a468); v3 2026-11-16 (d3ab703)} \\ \hline
\textbf{Objective:} & \multicolumn{3}{L{\dimexpr 0.63\linewidth+4\tabcolsep+2\arrayrulewidth\relax}|}{Check that augmentation parameters are fixed per (session, copy) and reproducible, that the same transform is applied to every frame of a copy, that the augmented frame source keeps the original timestamps, and that augmented rows never enter a test fold.} \\ \hline
\textbf{Product/Ver/ Module:} & \multicolumn{3}{L{\dimexpr 0.63\linewidth+4\tabcolsep+2\arrayrulewidth\relax}|}{FYP-1 attention pipeline, commit e743616. Ingestion and training: makes augmented copies of whole sessions (flip, colour shift, noise, masks) that are read on the fly and used only to train models and tune the tracker.} \\ \hline
\textbf{Environment:} & \multicolumn{3}{L{\dimexpr 0.63\linewidth+4\tabcolsep+2\arrayrulewidth\relax}|}{Windows 11 (10.0.26200); conda environment fyp-local: Python 3.13.18, NumPy 2.5.4, SciPy 1.18.1, scikit-learn 1.9.1, pandas 3.1.1, OpenCV 4.11.0, PyTorch 2.9.0+cu126. CPU only; no GPU, dataset or model weights needed unless stated. Run from the repository root: python -m unittest discover -s tests -t .} \\ \hline
\textbf{Assumptions:} & \multicolumn{3}{L{\dimexpr 0.63\linewidth+4\tabcolsep+2\arrayrulewidth\relax}|}{The augmentation settings come from configs/\allowbreak{}deliverable3.json (hue\_\allowbreak{}deg 3 at the time of writing). Small synthetic ZIP files of 4 PNG frames (64 x 48 random noise, named images/\allowbreak{}10\_\allowbreak{}00\_\allowbreak{}0<i>\_\allowbreak{}000000.png) stand in for DIPSER camera ZIPs. A tiny DataFrame stands in for the training table.} \\ \hline
\textbf{Pre-Requisite:} & \multicolumn{3}{L{\dimexpr 0.63\linewidth+4\tabcolsep+2\arrayrulewidth\relax}|}{Run from the repository root; configs/\allowbreak{}deliverable3.json and the results files read by configs/\allowbreak{}derived.py must exist (the config is loaded at import). OpenCV, pandas and PyTorch installed. Test ZIPs are written only to temporary directories.} \\ \hline
\textbf{Step No.} & \textbf{Execution description} & \multicolumn{2}{L{\dimexpr 0.36\linewidth+2\tabcolsep+1\arrayrulewidth\relax}|}{\textbf{Procedure result}} \\ \hline
1 & \texttt{test\_\allowbreak{}hflip\_\allowbreak{}mirrors\_\allowbreak{}boxes}: Call hflip\_\allowbreak{}boxes on box [[10, 5, 30, 20]] with frame width 100, then flip the result again. & \multicolumn{2}{L{\dimexpr 0.36\linewidth+2\tabcolsep+1\arrayrulewidth\relax}|}{The flipped box is [[70, 5, 90, 20]]. Flipping twice returns the original box. (Passed)} \\ \hline
2 & \texttt{test\_\allowbreak{}parameters\_\allowbreak{}are\_\allowbreak{}fixed\_\allowbreak{}per\_\allowbreak{}session\_\allowbreak{}and\_\allowbreak{}copy}: Call sample\_\allowbreak{}augmentation for session 'group\_\allowbreak{}01/\allowbreak{}experiment\_\allowbreak{}01/\allowbreak{}context\_\allowbreak{}cam\_\allowbreak{}01' copy 1 twice, for the same session copy 2, and for 'group\_\allowbreak{}02/\allowbreak{}experiment\_\allowbreak{}01/\allowbreak{}context\_\allowbreak{}cam\_\allowbreak{}01' copy 1. & \multicolumn{2}{L{\dimexpr 0.36\linewidth+2\tabcolsep+1\arrayrulewidth\relax}|}{The two calls with the same session and copy give equal parameters. A different copy number, and a different session, each give different parameters. (Passed)} \\ \hline
3 & \texttt{test\_\allowbreak{}hue\_\allowbreak{}shift\_\allowbreak{}stays\_\allowbreak{}small}: Sample parameters for 50 sessions ('s0' to 's49') and copies 1 and 2, and collect hue\_\allowbreak{}deg. & \multicolumn{2}{L{\dimexpr 0.36\linewidth+2\tabcolsep+1\arrayrulewidth\relax}|}{Every |hue\_\allowbreak{}deg| is at most the configured hue\_\allowbreak{}deg, and the configured hue\_\allowbreak{}deg is at most 5 degrees. (Passed)} \\ \hline
4 & \texttt{test\_\allowbreak{}tag\_\allowbreak{}ignores\_\allowbreak{}copy\_\allowbreak{}count}: Compute aug\_\allowbreak{}tag for copy 1 with the config's settings at copies = 1 and at copies = 3, and with hue\_\allowbreak{}deg changed to 4. & \multicolumn{2}{L{\dimexpr 0.36\linewidth+2\tabcolsep+1\arrayrulewidth\relax}|}{The tag is the same for copies = 1 and copies = 3. Changing hue\_\allowbreak{}deg to 4 gives a different tag. (Passed)} \\ \hline
5 & \texttt{test\_\allowbreak{}same\_\allowbreak{}transform\_\allowbreak{}on\_\allowbreak{}every\_\allowbreak{}frame\_\allowbreak{}noise\_\allowbreak{}aside}: Take the sampled parameters for session 's' copy 1, with noise 0 and hflip True. Apply augment\_\allowbreak{}frame to two identical 8 x 8 frames of value 40, as frame indices 0 and 1. & \multicolumn{2}{L{\dimexpr 0.36\linewidth+2\tabcolsep+1\arrayrulewidth\relax}|}{The two output frames are identical: the flip, colour shift and masks do not change from frame to frame. (Passed)} \\ \hline
6 & \texttt{test\_\allowbreak{}source\_\allowbreak{}is\_\allowbreak{}a\_\allowbreak{}dataset\_\allowbreak{}with\_\allowbreak{}the\_\allowbreak{}original\_\allowbreak{}timestamps}: Write a 4-frame ZIP. Build an AugmentedZipFrameSource with only hflip on (noise 0, no masks, neutral colour: brightness 0, contrast 1, saturation 1, hue 0, gamma 1) and a plain ZipFrameSource on the same ZIP. & \multicolumn{2}{L{\dimexpr 0.36\linewidth+2\tabcolsep+1\arrayrulewidth\relax}|}{The augmented source is a torch Dataset with length 4. Its timestamps equal the plain source's. Item 2's image equals the plain frame 2 flipped horizontally. Iterating it yields frame indices [0, 1, 2, 3]. (Passed)} \\ \hline
7 & \texttt{test\_\allowbreak{}a\_\allowbreak{}copy\_\allowbreak{}is\_\allowbreak{}the\_\allowbreak{}same\_\allowbreak{}video\_\allowbreak{}every\_\allowbreak{}time\_\allowbreak{}it\_\allowbreak{}is\_\allowbreak{}read}: Write a 4-frame ZIP. Read all frames through AugmentedZipFrameSource with copy 1 parameters for 'group\_\allowbreak{}01/\allowbreak{}experiment\_\allowbreak{}01/\allowbreak{}context\_\allowbreak{}cam\_\allowbreak{}01' and noise 0.3, then read them again with freshly sampled copy 1 parameters and noise 0.3. Read frame 0 with copy 2 parameters and noise 0.3. & \multicolumn{2}{L{\dimexpr 0.36\linewidth+2\tabcolsep+1\arrayrulewidth\relax}|}{Every frame of the two copy 1 reads is identical, noise included. Frame 0 of copy 2 differs from frame 0 of copy 1. (Passed)} \\ \hline
8 & \texttt{test\_\allowbreak{}augmented\_\allowbreak{}rows\_\allowbreak{}never\_\allowbreak{}reach\_\allowbreak{}a\_\allowbreak{}test\_\allowbreak{}fold}: Build a table of 6 rows: groups [g1, g1, g2, g2, g1, g2], aug flags [False, False, False, False, True, True]. Iterate heldout\_\allowbreak{}splits from scripts/\allowbreak{}train\_\allowbreak{}models.py. & \multicolumn{2}{L{\dimexpr 0.36\linewidth+2\tabcolsep+1\arrayrulewidth\relax}|}{In every fold, no test row is augmented; the training rows include augmented rows (copies of the other group); and no training row belongs to the held-out group. (Passed)} \\ \hline
\multicolumn{4}{|L{\dimexpr 0.84\linewidth+6\tabcolsep+3\arrayrulewidth\relax}|}{\textbf{Comments:} Augmented copies are never written to disk. test\_\allowbreak{}a\_\allowbreak{}copy\_\allowbreak{}is\_\allowbreak{}the\_\allowbreak{}same\_\allowbreak{}video\_\allowbreak{}every\_\allowbreak{}time\_\allowbreak{}it\_\allowbreak{}is\_\allowbreak{}read checks that seeding alone makes every read of a copy identical.} \\ \hline
\multicolumn{4}{|C{\dimexpr 0.84\linewidth+6\tabcolsep+3\arrayrulewidth\relax}|}{\textbf{\fbox{Passed} \quad Failed \quad Not Executed}} \\ \hline
\end{longtable}

\begin{longtable}{|L{0.21\linewidth}|L{0.27\linewidth}|L{0.17\linewidth}|L{0.19\linewidth}|}
\caption{Test case TC-08: Colour augmentation (lookup-table implementation)}\label{tab:tc-08} \\ \hline
\multicolumn{4}{|C{\dimexpr 0.84\linewidth+6\tabcolsep+3\arrayrulewidth\relax}|}{\textbf{Colour augmentation (lookup-table implementation)}} \\ \hline
\multicolumn{4}{|C{\dimexpr 0.84\linewidth+6\tabcolsep+3\arrayrulewidth\relax}|}{\textbf{src/\allowbreak{}ingestion/\allowbreak{}augment.py; tests/\allowbreak{}test\_\allowbreak{}augment.py (TestColourLookupTables)}} \\ \hline
\textbf{Test Case ID:} & TC-08 & \textbf{QA Test Engineer:} & Abi \\ \hline
\textbf{Test case Version:} & 3 & \textbf{Reviewed By:} & - \\ \hline
\textbf{Test Date:} & 2026-11-18 & \textbf{Use Case Reference(s):} & - \\ \hline
\textbf{Revision History:} & \multicolumn{3}{L{\dimexpr 0.63\linewidth+4\tabcolsep+2\arrayrulewidth\relax}|}{v1 2026-11-15 (a497de6); v2 2026-11-15 (be0a468); v3 2026-11-16 (d3ab703)} \\ \hline
\textbf{Objective:} & \multicolumn{3}{L{\dimexpr 0.63\linewidth+4\tabcolsep+2\arrayrulewidth\relax}|}{Check that the fast lookup-table colour transform gives bit-identical output to the per-pixel floating-point reference.} \\ \hline
\textbf{Product/Ver/ Module:} & \multicolumn{3}{L{\dimexpr 0.63\linewidth+4\tabcolsep+2\arrayrulewidth\relax}|}{FYP-1 attention pipeline, commit e743616. Ingestion: the colour step (brightness and contrast, hue and saturation, gamma) applied to every frame of an augmented copy.} \\ \hline
\textbf{Environment:} & \multicolumn{3}{L{\dimexpr 0.63\linewidth+4\tabcolsep+2\arrayrulewidth\relax}|}{Windows 11 (10.0.26200); conda environment fyp-local: Python 3.13.18, NumPy 2.5.4, SciPy 1.18.1, scikit-learn 1.9.1, pandas 3.1.1, OpenCV 4.11.0, PyTorch 2.9.0+cu126. CPU only; no GPU, dataset or model weights needed unless stated. Run from the repository root: python -m unittest discover -s tests -t .} \\ \hline
\textbf{Assumptions:} & \multicolumn{3}{L{\dimexpr 0.63\linewidth+4\tabcolsep+2\arrayrulewidth\relax}|}{A random 48 x 64 BGR image (numpy seed 0) stands in for a video frame. Parameter ranges cover and slightly exceed the configured augmentation ranges.} \\ \hline
\textbf{Pre-Requisite:} & \multicolumn{3}{L{\dimexpr 0.63\linewidth+4\tabcolsep+2\arrayrulewidth\relax}|}{OpenCV installed. As for the rest of the file, the config is loaded at import.} \\ \hline
\textbf{Step No.} & \textbf{Execution description} & \multicolumn{2}{L{\dimexpr 0.36\linewidth+2\tabcolsep+1\arrayrulewidth\relax}|}{\textbf{Procedure result}} \\ \hline
1 & \texttt{test\_\allowbreak{}lookup\_\allowbreak{}tables\_\allowbreak{}reproduce\_\allowbreak{}the\_\allowbreak{}per\_\allowbreak{}pixel\_\allowbreak{}maths\_\allowbreak{}exactly}: For 25 random parameter sets (brightness in [-35, 35], contrast in [0.85, 1.15], saturation in [0.8, 1.2], hue in [-3, 3] degrees, gamma = exp of a value in [-0.14, 0.14]), call colour and \_\allowbreak{}colour\_\allowbreak{}reference on the same image. & \multicolumn{2}{L{\dimexpr 0.36\linewidth+2\tabcolsep+1\arrayrulewidth\relax}|}{For every parameter set the two output images are exactly equal. (Passed)} \\ \hline
\multicolumn{4}{|L{\dimexpr 0.84\linewidth+6\tabcolsep+3\arrayrulewidth\relax}|}{\textbf{Comments:} colour is about 20 times faster than the reference according to its docstring; the test checks only equality, not speed.} \\ \hline
\multicolumn{4}{|C{\dimexpr 0.84\linewidth+6\tabcolsep+3\arrayrulewidth\relax}|}{\textbf{\fbox{Passed} \quad Failed \quad Not Executed}} \\ \hline
\end{longtable}

\begin{longtable}{|L{0.21\linewidth}|L{0.27\linewidth}|L{0.17\linewidth}|L{0.19\linewidth}|}
\caption{Test case TC-09: Derived configuration values}\label{tab:tc-09} \\ \hline
\multicolumn{4}{|C{\dimexpr 0.84\linewidth+6\tabcolsep+3\arrayrulewidth\relax}|}{\textbf{Derived configuration values}} \\ \hline
\multicolumn{4}{|C{\dimexpr 0.84\linewidth+6\tabcolsep+3\arrayrulewidth\relax}|}{\textbf{src/\allowbreak{}core/\allowbreak{}derive.py, src/\allowbreak{}core/\allowbreak{}config.py; tests/\allowbreak{}test\_\allowbreak{}derive.py (TestDerive)}} \\ \hline
\textbf{Test Case ID:} & TC-09 & \textbf{QA Test Engineer:} & Abi \\ \hline
\textbf{Test case Version:} & 1 & \textbf{Reviewed By:} & - \\ \hline
\textbf{Test Date:} & 2026-11-18 & \textbf{Use Case Reference(s):} & Design \S{}9 (config hashing) \\ \hline
\textbf{Revision History:} & \multicolumn{3}{L{\dimexpr 0.63\linewidth+4\tabcolsep+2\arrayrulewidth\relax}|}{v1 2026-11-17 (cbcaa22)} \\ \hline
\textbf{Objective:} & \multicolumn{3}{L{\dimexpr 0.63\linewidth+4\tabcolsep+2\arrayrulewidth\relax}|}{Check that config values derived from experiment results files are read correctly with their sources, are applied by load\_\allowbreak{}config before command-line overrides, are hashed by value only, and fail loudly when a results file or key is missing.} \\ \hline
\textbf{Product/Ver/ Module:} & \multicolumn{3}{L{\dimexpr 0.63\linewidth+4\tabcolsep+2\arrayrulewidth\relax}|}{FYP-1 attention pipeline, commit e743616. Configuration: a derived module (such as configs/\allowbreak{}derived.py) sets config values such as the detector threshold and tracker options from experiment results, every time a config is loaded.} \\ \hline
\textbf{Environment:} & \multicolumn{3}{L{\dimexpr 0.63\linewidth+4\tabcolsep+2\arrayrulewidth\relax}|}{Windows 11 (10.0.26200); conda environment fyp-local: Python 3.13.18, NumPy 2.5.4, SciPy 1.18.1, scikit-learn 1.9.1, pandas 3.1.1, OpenCV 4.11.0, PyTorch 2.9.0+cu126. CPU only; no GPU, dataset or model weights needed unless stated. Run from the repository root: python -m unittest discover -s tests -t .} \\ \hline
\textbf{Assumptions:} & \multicolumn{3}{L{\dimexpr 0.63\linewidth+4\tabcolsep+2\arrayrulewidth\relax}|}{setUp writes, into a temporary directory, a results file res.json = \{optimal\_\allowbreak{}threshold: 0.11, config\_\allowbreak{}patch: \{'tracker.options': \{a: 1\}\}\}, a derived module whose OVERRIDES set detector.conf from optimal\_\allowbreak{}threshold, tracker.options from config\_\allowbreak{}patch.tracker.options and window.length\_\allowbreak{}s to the plain value 6, and a copy of configs/\allowbreak{}deliverable3.json whose 'derived' key points to that module. The results key 'tracker.options' contains a dot, as in the real results files.} \\ \hline
\textbf{Pre-Requisite:} & \multicolumn{3}{L{\dimexpr 0.63\linewidth+4\tabcolsep+2\arrayrulewidth\relax}|}{Run from the repository root (configs/\allowbreak{}deliverable3.json is read as the base config). The real configs/\allowbreak{}derived.py and its results files are not used. Files are written only to a temporary directory that tearDown removes.} \\ \hline
\textbf{Step No.} & \textbf{Execution description} & \multicolumn{2}{L{\dimexpr 0.36\linewidth+2\tabcolsep+1\arrayrulewidth\relax}|}{\textbf{Procedure result}} \\ \hline
1 & \texttt{test\_\allowbreak{}values\_\allowbreak{}sources\_\allowbreak{}and\_\allowbreak{}keys\_\allowbreak{}with\_\allowbreak{}dots}: Call result(res.json, 'config\_\allowbreak{}patch.tracker.options'), then derived\_\allowbreak{}values on the temporary derived module. & \multicolumn{2}{L{\dimexpr 0.36\linewidth+2\tabcolsep+1\arrayrulewidth\relax}|}{result returns Sourced(value \{'a': 1\}, source '<path to res.json>\#config\_\allowbreak{}patch.tracker.options'), resolving the dotted key 'tracker.options' as one key. derived\_\allowbreak{}values returns values \{detector.conf: 0.11, tracker.options: \{a: 1\}, window.length\_\allowbreak{}s: 6\}. The source of the plain value window.length\_\allowbreak{}s ends with 'derived.py' (the module itself). (Passed)} \\ \hline
2 & \texttt{test\_\allowbreak{}load\_\allowbreak{}config\_\allowbreak{}applies\_\allowbreak{}them\_\allowbreak{}before\_\allowbreak{}set\_\allowbreak{}and\_\allowbreak{}hashes\_\allowbreak{}values\_\allowbreak{}only}: Load the temporary config with load\_\allowbreak{}config, then again with the override detector.conf = 0.3. Make a copy of the loaded config with derived\_\allowbreak{}from and derived changed, and compare config\_\allowbreak{}hash of both. & \multicolumn{2}{L{\dimexpr 0.36\linewidth+2\tabcolsep+1\arrayrulewidth\relax}|}{The loaded config has detector.conf 0.11 and tracker.options \{'a': 1\}. With the override, detector.conf is 0.3 (the override wins over the derived value). The two hashes are equal: where a value came from is not hashed. (Passed)} \\ \hline
3 & \texttt{test\_\allowbreak{}a\_\allowbreak{}missing\_\allowbreak{}result\_\allowbreak{}fails\_\allowbreak{}loudly}: Call result on a non-existent file missing.json with key 'x', and on res.json with the missing key 'config\_\allowbreak{}patch.nope'. & \multicolumn{2}{L{\dimexpr 0.36\linewidth+2\tabcolsep+1\arrayrulewidth\relax}|}{The first call raises FileNotFoundError; the second raises KeyError. (Passed)} \\ \hline
\multicolumn{4}{|L{\dimexpr 0.84\linewidth+6\tabcolsep+3\arrayrulewidth\relax}|}{\textbf{Comments:} The last step documents a design choice: the pipeline never falls back to a stale constant when an experiment's results file is missing.} \\ \hline
\multicolumn{4}{|C{\dimexpr 0.84\linewidth+6\tabcolsep+3\arrayrulewidth\relax}|}{\textbf{\fbox{Passed} \quad Failed \quad Not Executed}} \\ \hline
\end{longtable}

\begin{longtable}{|L{0.21\linewidth}|L{0.27\linewidth}|L{0.17\linewidth}|L{0.19\linewidth}|}
\caption{Test case TC-10: Head-direction geometry from COCO keypoints}\label{tab:tc-10} \\ \hline
\multicolumn{4}{|C{\dimexpr 0.84\linewidth+6\tabcolsep+3\arrayrulewidth\relax}|}{\textbf{Head-direction geometry from COCO keypoints}} \\ \hline
\multicolumn{4}{|C{\dimexpr 0.84\linewidth+6\tabcolsep+3\arrayrulewidth\relax}|}{\textbf{src/\allowbreak{}bfmu/\allowbreak{}head\_\allowbreak{}pose.py; tests/\allowbreak{}test\_\allowbreak{}head\_\allowbreak{}pose.py (TestHeadGeometry)}} \\ \hline
\textbf{Test Case ID:} & TC-10 & \textbf{QA Test Engineer:} & Abi \\ \hline
\textbf{Test case Version:} & 1 & \textbf{Reviewed By:} & - \\ \hline
\textbf{Test Date:} & 2026-11-18 & \textbf{Use Case Reference(s):} & - \\ \hline
\textbf{Revision History:} & \multicolumn{3}{L{\dimexpr 0.63\linewidth+4\tabcolsep+2\arrayrulewidth\relax}|}{v1 (not yet committed)} \\ \hline
\textbf{Objective:} & \multicolumn{3}{L{\dimexpr 0.63\linewidth+4\tabcolsep+2\arrayrulewidth\relax}|}{Check that head\_\allowbreak{}geometry computes yaw, head height, face confidence and the matched flag correctly from one person's keypoints, including the fallback to shoulders when ears are not visible and NaN output when the nose is missing.} \\ \hline
\textbf{Product/Ver/ Module:} & \multicolumn{3}{L{\dimexpr 0.63\linewidth+4\tabcolsep+2\arrayrulewidth\relax}|}{FYP-1 attention pipeline, commit e743616. Head-pose feature unit (BFMU d, experiment 09 visual focus of attention): per-frame head direction for each tracked student.} \\ \hline
\textbf{Environment:} & \multicolumn{3}{L{\dimexpr 0.63\linewidth+4\tabcolsep+2\arrayrulewidth\relax}|}{Windows 11 (10.0.26200); conda environment fyp-local: Python 3.13.18, NumPy 2.5.4, SciPy 1.18.1, scikit-learn 1.9.1, pandas 3.1.1, OpenCV 4.11.0, PyTorch 2.9.0+cu126. CPU only; no GPU, dataset or model weights needed unless stated. Run from the repository root: python -m unittest discover -s tests -t .} \\ \hline
\textbf{Assumptions:} & \multicolumn{3}{L{\dimexpr 0.63\linewidth+4\tabcolsep+2\arrayrulewidth\relax}|}{Synthetic keypoints from the helper person(): an upright person with nose at (100 + nose\_\allowbreak{}dx, nose\_\allowbreak{}y), eyes at y 35, ears at x 85 and 115 (30 px apart), shoulders at (60, 100) and (140, 100) (80 px apart). Keypoint confidence 0.9; ears set to 0.1 when ears=False. Visibility threshold kp\_\allowbreak{}conf is the default 0.5.} \\ \hline
\textbf{Pre-Requisite:} & \multicolumn{3}{L{\dimexpr 0.63\linewidth+4\tabcolsep+2\arrayrulewidth\relax}|}{None beyond the Python environment (numpy). No pose model weights or GPU.} \\ \hline
\textbf{Step No.} & \textbf{Execution description} & \multicolumn{2}{L{\dimexpr 0.36\linewidth+2\tabcolsep+1\arrayrulewidth\relax}|}{\textbf{Procedure result}} \\ \hline
1 & \texttt{test\_\allowbreak{}facing\_\allowbreak{}forward}: Call head\_\allowbreak{}geometry on the default person (nose at (100, 40), centred between the ears). & \multicolumn{2}{L{\dimexpr 0.36\linewidth+2\tabcolsep+1\arrayrulewidth\relax}|}{hp.yaw = 0.0, hp.head\_\allowbreak{}height = 60/\allowbreak{}80 = 0.75 (shoulder y 100 minus nose y 40, over shoulder width 80), hp.matched = 1.0, hp.face\_\allowbreak{}conf = 0.9. Observed: as expected. (Passed)} \\ \hline
2 & \texttt{test\_\allowbreak{}turned\_\allowbreak{}and\_\allowbreak{}lowered}: Call head\_\allowbreak{}geometry on a person with the nose moved 9 px left (nose\_\allowbreak{}dx = -9), then on a person with the nose lowered to y 80. & \multicolumn{2}{L{\dimexpr 0.36\linewidth+2\tabcolsep+1\arrayrulewidth\relax}|}{Turned: hp.yaw is negative and equals -9/\allowbreak{}30 = -0.3, using the ears as the reference pair. Lowered: hp.head\_\allowbreak{}height (20/\allowbreak{}80) is smaller than the default person's 0.75. Observed: as expected. (Passed)} \\ \hline
3 & \texttt{test\_\allowbreak{}shoulders\_\allowbreak{}replace\_\allowbreak{}missing\_\allowbreak{}ears}: Call head\_\allowbreak{}geometry on a person with nose\_\allowbreak{}dx = 8 and ear confidence 0.1 (below kp\_\allowbreak{}conf). & \multicolumn{2}{L{\dimexpr 0.36\linewidth+2\tabcolsep+1\arrayrulewidth\relax}|}{hp.yaw = 8/\allowbreak{}80 = 0.1: the shoulders replace the ears as the reference pair. Observed: as expected. (Passed)} \\ \hline
4 & \texttt{test\_\allowbreak{}missing\_\allowbreak{}nose\_\allowbreak{}gives\_\allowbreak{}nan}: Take the default person and set the nose confidence to 0.1, then call head\_\allowbreak{}geometry. & \multicolumn{2}{L{\dimexpr 0.36\linewidth+2\tabcolsep+1\arrayrulewidth\relax}|}{hp.yaw and hp.head\_\allowbreak{}height are both NaN, and hp.matched is still 1.0 (a pose was present). Observed: as expected. (Passed)} \\ \hline
\multicolumn{4}{|L{\dimexpr 0.84\linewidth+6\tabcolsep+3\arrayrulewidth\relax}|}{\textbf{Comments:} The module docstring ties this unit to experiment 09 (visual focus of attention); it cites no design-document section.} \\ \hline
\multicolumn{4}{|C{\dimexpr 0.84\linewidth+6\tabcolsep+3\arrayrulewidth\relax}|}{\textbf{\fbox{Passed} \quad Failed \quad Not Executed}} \\ \hline
\end{longtable}

\begin{longtable}{|L{0.21\linewidth}|L{0.27\linewidth}|L{0.17\linewidth}|L{0.19\linewidth}|}
\caption{Test case TC-11: Pose-to-track matching}\label{tab:tc-11} \\ \hline
\multicolumn{4}{|C{\dimexpr 0.84\linewidth+6\tabcolsep+3\arrayrulewidth\relax}|}{\textbf{Pose-to-track matching}} \\ \hline
\multicolumn{4}{|C{\dimexpr 0.84\linewidth+6\tabcolsep+3\arrayrulewidth\relax}|}{\textbf{src/\allowbreak{}bfmu/\allowbreak{}head\_\allowbreak{}pose.py, src/\allowbreak{}tracking/\allowbreak{}association.py; tests/\allowbreak{}test\_\allowbreak{}head\_\allowbreak{}pose.py (TestMatching)}} \\ \hline
\textbf{Test Case ID:} & TC-11 & \textbf{QA Test Engineer:} & Abi \\ \hline
\textbf{Test case Version:} & 1 & \textbf{Reviewed By:} & - \\ \hline
\textbf{Test Date:} & 2026-11-18 & \textbf{Use Case Reference(s):} & - \\ \hline
\textbf{Revision History:} & \multicolumn{3}{L{\dimexpr 0.63\linewidth+4\tabcolsep+2\arrayrulewidth\relax}|}{v1 (not yet committed)} \\ \hline
\textbf{Objective:} & \multicolumn{3}{L{\dimexpr 0.63\linewidth+4\tabcolsep+2\arrayrulewidth\relax}|}{Check that match\_\allowbreak{}poses pairs pose instances with tracks one-to-one by box IoU, drops pairs below the IoU threshold, and handles an empty pose list.} \\ \hline
\textbf{Product/Ver/ Module:} & \multicolumn{3}{L{\dimexpr 0.63\linewidth+4\tabcolsep+2\arrayrulewidth\relax}|}{FYP-1 attention pipeline, commit e743616. Head-pose unit, 'frame' input mode: links each pose detected in the frame to the tracked student it belongs to.} \\ \hline
\textbf{Environment:} & \multicolumn{3}{L{\dimexpr 0.63\linewidth+4\tabcolsep+2\arrayrulewidth\relax}|}{Windows 11 (10.0.26200); conda environment fyp-local: Python 3.13.18, NumPy 2.5.4, SciPy 1.18.1, scikit-learn 1.9.1, pandas 3.1.1, OpenCV 4.11.0, PyTorch 2.9.0+cu126. CPU only; no GPU, dataset or model weights needed unless stated. Run from the repository root: python -m unittest discover -s tests -t .} \\ \hline
\textbf{Assumptions:} & \multicolumn{3}{L{\dimexpr 0.63\linewidth+4\tabcolsep+2\arrayrulewidth\relax}|}{Synthetic boxes. min\_\allowbreak{}iou is 0.3.} \\ \hline
\textbf{Pre-Requisite:} & \multicolumn{3}{L{\dimexpr 0.63\linewidth+4\tabcolsep+2\arrayrulewidth\relax}|}{None beyond the Python environment (numpy, scipy).} \\ \hline
\textbf{Step No.} & \textbf{Execution description} & \multicolumn{2}{L{\dimexpr 0.36\linewidth+2\tabcolsep+1\arrayrulewidth\relax}|}{\textbf{Procedure result}} \\ \hline
1 & \texttt{test\_\allowbreak{}one\_\allowbreak{}to\_\allowbreak{}one\_\allowbreak{}by\_\allowbreak{}iou}: Tracks [0,0,10,10] and [20,0,30,10]. Poses [21,0,31,10], [0,1,10,11] and [100,100,110,110]. Call match\_\allowbreak{}poses(poses, tracks, 0.3). Then call it with an empty (0, 4) pose array. & \multicolumn{2}{L{\dimexpr 0.36\linewidth+2\tabcolsep+1\arrayrulewidth\relax}|}{Returns \{0: 1, 1: 0\} (track index to pose index): track 0 gets pose 1, track 1 gets pose 0, and the far pose is unmatched. With no poses the result is \{\}. Observed: as expected. (Passed)} \\ \hline
\multicolumn{4}{|L{\dimexpr 0.84\linewidth+6\tabcolsep+3\arrayrulewidth\relax}|}{\textbf{Comments:} -} \\ \hline
\multicolumn{4}{|C{\dimexpr 0.84\linewidth+6\tabcolsep+3\arrayrulewidth\relax}|}{\textbf{\fbox{Passed} \quad Failed \quad Not Executed}} \\ \hline
\end{longtable}

\begin{longtable}{|L{0.21\linewidth}|L{0.27\linewidth}|L{0.17\linewidth}|L{0.19\linewidth}|}
\caption{Test case TC-12: Head-pose feature extractor (BFMU d)}\label{tab:tc-12} \\ \hline
\multicolumn{4}{|C{\dimexpr 0.84\linewidth+6\tabcolsep+3\arrayrulewidth\relax}|}{\textbf{Head-pose feature extractor (BFMU d)}} \\ \hline
\multicolumn{4}{|C{\dimexpr 0.84\linewidth+6\tabcolsep+3\arrayrulewidth\relax}|}{\textbf{src/\allowbreak{}bfmu/\allowbreak{}head\_\allowbreak{}pose.py, src/\allowbreak{}core/\allowbreak{}types.py, src/\allowbreak{}bfmu/\allowbreak{}base.py; tests/\allowbreak{}test\_\allowbreak{}head\_\allowbreak{}pose.py (TestUnit)}} \\ \hline
\textbf{Test Case ID:} & TC-12 & \textbf{QA Test Engineer:} & Abi \\ \hline
\textbf{Test case Version:} & 1 & \textbf{Reviewed By:} & - \\ \hline
\textbf{Test Date:} & 2026-11-18 & \textbf{Use Case Reference(s):} & - \\ \hline
\textbf{Revision History:} & \multicolumn{3}{L{\dimexpr 0.63\linewidth+4\tabcolsep+2\arrayrulewidth\relax}|}{v1 (not yet committed)} \\ \hline
\textbf{Objective:} & \multicolumn{3}{L{\dimexpr 0.63\linewidth+4\tabcolsep+2\arrayrulewidth\relax}|}{Check the extractor without a pose model: 'tracks' mode reads keypoints carried by the tracks, detections keep their keypoints when the tracker selects a subset, and window features aggregate the per-frame signals correctly.} \\ \hline
\textbf{Product/Ver/ Module:} & \multicolumn{3}{L{\dimexpr 0.63\linewidth+4\tabcolsep+2\arrayrulewidth\relax}|}{FYP-1 attention pipeline, commit e743616. Head-pose unit in the per-frame and feature passes; in 'tracks' mode (experiment 09's setting) the person detector is itself a pose model and passes keypoints through the tracker.} \\ \hline
\textbf{Environment:} & \multicolumn{3}{L{\dimexpr 0.63\linewidth+4\tabcolsep+2\arrayrulewidth\relax}|}{Windows 11 (10.0.26200); conda environment fyp-local: Python 3.13.18, NumPy 2.5.4, SciPy 1.18.1, scikit-learn 1.9.1, pandas 3.1.1, OpenCV 4.11.0, PyTorch 2.9.0+cu126. CPU only; no GPU, dataset or model weights needed unless stated. Run from the repository root: python -m unittest discover -s tests -t .} \\ \hline
\textbf{Assumptions:} & \multicolumn{3}{L{\dimexpr 0.63\linewidth+4\tabcolsep+2\arrayrulewidth\relax}|}{Synthetic tracks, detections and tracklets; a blank 4 x 4 frame. No pose network is run in any step.} \\ \hline
\textbf{Pre-Requisite:} & \multicolumn{3}{L{\dimexpr 0.63\linewidth+4\tabcolsep+2\arrayrulewidth\relax}|}{None beyond the Python environment (numpy, scipy). No pose model weights or GPU: the model loads lazily and is never accessed here.} \\ \hline
\textbf{Step No.} & \textbf{Execution description} & \multicolumn{2}{L{\dimexpr 0.36\linewidth+2\tabcolsep+1\arrayrulewidth\relax}|}{\textbf{Procedure result}} \\ \hline
1 & \texttt{test\_\allowbreak{}tracks\_\allowbreak{}mode\_\allowbreak{}reads\_\allowbreak{}keypoints\_\allowbreak{}from\_\allowbreak{}the\_\allowbreak{}tracks}: Build head\_\allowbreak{}pose in mode 'tracks'. Track 1 carries the keypoints of person(nose\_\allowbreak{}dx=-9); track 2 carries no keypoints. Call frame\_\allowbreak{}signals on a blank 4 x 4 frame. & \multicolumn{2}{L{\dimexpr 0.36\linewidth+2\tabcolsep+1\arrayrulewidth\relax}|}{Track 1 has hp.yaw = -9/\allowbreak{}30 = -0.3 and hp.matched = 1.0. Track 2 has hp.matched = 0.0 and hp.yaw NaN. unit.\_\allowbreak{}model is still None (no model of its own). Observed: as expected. (Passed)} \\ \hline
2 & \texttt{test\_\allowbreak{}tracker\_\allowbreak{}hands\_\allowbreak{}each\_\allowbreak{}track\_\allowbreak{}its\_\allowbreak{}detections\_\allowbreak{}keypoints}: Build Detections with 2 boxes, confidences [0.9, 0.8], keypoints arange(68) reshaped to (2, 17, 2) and keypoint confidences of ones. Call take with mask [False, True]. & \multicolumn{2}{L{\dimexpr 0.36\linewidth+2\tabcolsep+1\arrayrulewidth\relax}|}{The result has length 1 and its keypoints[0] equals the second detection's keypoints. Observed: as expected. (Passed)} \\ \hline
3 & \texttt{test\_\allowbreak{}registered\_\allowbreak{}and\_\allowbreak{}window\_\allowbreak{}features}: Build head\_\allowbreak{}pose with defaults from the registry. Make a 10-frame tracklet (fps 5) with hp.yaw = five 0.2 then five -0.2, hp.head\_\allowbreak{}height = four 0.05 then six 0.3, hp.face\_\allowbreak{}conf all 0.8, hp.ear\_\allowbreak{}asym all 0, hp.matched = eight 1 then two 0. Call window\_\allowbreak{}features. Then call it on a 3-frame tracklet with no signals. & \multicolumn{2}{L{\dimexpr 0.36\linewidth+2\tabcolsep+1\arrayrulewidth\relax}|}{uses\_\allowbreak{}frames is True. The feature names equal unit.keys. hp.yaw.mean = 0.0, hp.yaw.abs = 0.2, hp.head\_\allowbreak{}low = 0.4 (4 of 10 frames below the 0.1 threshold), hp.matched = 0.8. For the empty tracklet every feature is NaN. Observed: as expected. (Passed)} \\ \hline
\multicolumn{4}{|L{\dimexpr 0.84\linewidth+6\tabcolsep+3\arrayrulewidth\relax}|}{\textbf{Comments:} The second step tests Detections.take in src/\allowbreak{}core/\allowbreak{}types.py, which the tracker uses to keep each detection's keypoints aligned with its box.} \\ \hline
\multicolumn{4}{|C{\dimexpr 0.84\linewidth+6\tabcolsep+3\arrayrulewidth\relax}|}{\textbf{\fbox{Passed} \quad Failed \quad Not Executed}} \\ \hline
\end{longtable}

\begin{longtable}{|L{0.21\linewidth}|L{0.27\linewidth}|L{0.17\linewidth}|L{0.19\linewidth}|}
\caption{Test case TC-13: Inter-labeler agreement (Krippendorff's alpha)}\label{tab:tc-13} \\ \hline
\multicolumn{4}{|C{\dimexpr 0.84\linewidth+6\tabcolsep+3\arrayrulewidth\relax}|}{\textbf{Inter-labeler agreement (Krippendorff's alpha)}} \\ \hline
\multicolumn{4}{|C{\dimexpr 0.84\linewidth+6\tabcolsep+3\arrayrulewidth\relax}|}{\textbf{scripts/\allowbreak{}labels/\allowbreak{}label\_\allowbreak{}agreement.py, src/\allowbreak{}evaluation/\allowbreak{}agreement.py; tests/\allowbreak{}test\_\allowbreak{}label\_\allowbreak{}agreement.py (TestAlpha)}} \\ \hline
\textbf{Test Case ID:} & TC-13 & \textbf{QA Test Engineer:} & Abi \\ \hline
\textbf{Test case Version:} & 2 & \textbf{Reviewed By:} & - \\ \hline
\textbf{Test Date:} & 2026-11-18 & \textbf{Use Case Reference(s):} & Design \S{}11; docs/\allowbreak{}experiments/\allowbreak{}06\_\allowbreak{}full\_\allowbreak{}pipeline\_\allowbreak{}and\_\allowbreak{}labels.md Part B \\ \hline
\textbf{Revision History:} & \multicolumn{3}{L{\dimexpr 0.63\linewidth+4\tabcolsep+2\arrayrulewidth\relax}|}{v1 2026-11-18 (5f5458c); v2 2026-11-18 (e664d52)} \\ \hline
\textbf{Objective:} & \multicolumn{3}{L{\dimexpr 0.63\linewidth+4\tabcolsep+2\arrayrulewidth\relax}|}{Check that the vectorised ordinal alpha used by the agreement analysis matches the reference implementation, and that identical labelers score perfect agreement.} \\ \hline
\textbf{Product/Ver/ Module:} & \multicolumn{3}{L{\dimexpr 0.63\linewidth+4\tabcolsep+2\arrayrulewidth\relax}|}{FYP-1 attention pipeline, commit e743616. Label agreement analysis script, which measures how far DIPSER's labelers agree as a ceiling for the model.} \\ \hline
\textbf{Environment:} & \multicolumn{3}{L{\dimexpr 0.63\linewidth+4\tabcolsep+2\arrayrulewidth\relax}|}{Windows 11 (10.0.26200); conda environment fyp-local: Python 3.13.18, NumPy 2.5.4, SciPy 1.18.1, scikit-learn 1.9.1, pandas 3.1.1, OpenCV 4.11.0, PyTorch 2.9.0+cu126. CPU only; no GPU, dataset or model weights needed unless stated. Run from the repository root: python -m unittest discover -s tests -t .} \\ \hline
\textbf{Assumptions:} & \multicolumn{3}{L{\dimexpr 0.63\linewidth+4\tabcolsep+2\arrayrulewidth\relax}|}{Synthetic labelers stand in for DIPSER's. The reference check uses a random 300 x 4 matrix of ratings 1-5 (seed 0) with 30\% missing and the first 150 rows of column 1 copied from column 0.} \\ \hline
\textbf{Pre-Requisite:} & \multicolumn{3}{L{\dimexpr 0.63\linewidth+4\tabcolsep+2\arrayrulewidth\relax}|}{None beyond the Python environment. The script module must import (it imports the dataset and pipeline packages) but no data is read.} \\ \hline
\textbf{Step No.} & \textbf{Execution description} & \multicolumn{2}{L{\dimexpr 0.36\linewidth+2\tabcolsep+1\arrayrulewidth\relax}|}{\textbf{Procedure result}} \\ \hline
1 & \texttt{test\_\allowbreak{}fast\_\allowbreak{}alpha\_\allowbreak{}matches\_\allowbreak{}the\_\allowbreak{}reference\_\allowbreak{}implementation}: Compute fast\_\allowbreak{}alpha on the random matrix and krippendorff\_\allowbreak{}alpha\_\allowbreak{}ordinal on the same values as lists. & \multicolumn{2}{L{\dimexpr 0.36\linewidth+2\tabcolsep+1\arrayrulewidth\relax}|}{The two alphas agree to 10 decimal places. (Passed)} \\ \hline
2 & \texttt{test\_\allowbreak{}identical\_\allowbreak{}labelers\_\allowbreak{}agree\_\allowbreak{}perfectly}: Build a student where labeler\_\allowbreak{}01 and labeler\_\allowbreak{}02 have the same track A (ratings 3, 2, 4, 3 at 0, 20, 50, 80 s). Get per-second step values with unit\_\allowbreak{}values and compute fast\_\allowbreak{}alpha. & \multicolumn{2}{L{\dimexpr 0.36\linewidth+2\tabcolsep+1\arrayrulewidth\relax}|}{Alpha is 1.0. (Passed)} \\ \hline
\multicolumn{4}{|L{\dimexpr 0.84\linewidth+6\tabcolsep+3\arrayrulewidth\relax}|}{\textbf{Comments:} -} \\ \hline
\multicolumn{4}{|C{\dimexpr 0.84\linewidth+6\tabcolsep+3\arrayrulewidth\relax}|}{\textbf{\fbox{Passed} \quad Failed \quad Not Executed}} \\ \hline
\end{longtable}

\begin{longtable}{|L{0.21\linewidth}|L{0.27\linewidth}|L{0.17\linewidth}|L{0.19\linewidth}|}
\caption{Test case TC-14: Human ceiling: predicting one labeler from the others}\label{tab:tc-14} \\ \hline
\multicolumn{4}{|C{\dimexpr 0.84\linewidth+6\tabcolsep+3\arrayrulewidth\relax}|}{\textbf{Human ceiling: predicting one labeler from the others}} \\ \hline
\multicolumn{4}{|C{\dimexpr 0.84\linewidth+6\tabcolsep+3\arrayrulewidth\relax}|}{\textbf{scripts/\allowbreak{}labels/\allowbreak{}label\_\allowbreak{}agreement.py, src/\allowbreak{}labels/\allowbreak{}interp.py; tests/\allowbreak{}test\_\allowbreak{}label\_\allowbreak{}agreement.py (TestCeiling)}} \\ \hline
\textbf{Test Case ID:} & TC-14 & \textbf{QA Test Engineer:} & Abi \\ \hline
\textbf{Test case Version:} & 2 & \textbf{Reviewed By:} & - \\ \hline
\textbf{Test Date:} & 2026-11-18 & \textbf{Use Case Reference(s):} & docs/\allowbreak{}experiments/\allowbreak{}06\_\allowbreak{}full\_\allowbreak{}pipeline\_\allowbreak{}and\_\allowbreak{}labels.md Part B \\ \hline
\textbf{Revision History:} & \multicolumn{3}{L{\dimexpr 0.63\linewidth+4\tabcolsep+2\arrayrulewidth\relax}|}{v1 2026-11-18 (5f5458c); v2 2026-11-18 (e664d52)} \\ \hline
\textbf{Objective:} & \multicolumn{3}{L{\dimexpr 0.63\linewidth+4\tabcolsep+2\arrayrulewidth\relax}|}{Check the predictor table built by ceiling\_\allowbreak{}rows: other labelers' curves at lag 0 and at lags, the +/\allowbreak{}-k s best match, run midpoints, the exclusion of the session-start entry, self-consistency columns and the mean of other experts.} \\ \hline
\textbf{Product/Ver/ Module:} & \multicolumn{3}{L{\dimexpr 0.63\linewidth+4\tabcolsep+2\arrayrulewidth\relax}|}{FYP-1 attention pipeline, commit e743616. Label agreement analysis script, ceiling part: scores how well one labeler's entries are predicted from the other labelers, the same way the model is scored.} \\ \hline
\textbf{Environment:} & \multicolumn{3}{L{\dimexpr 0.63\linewidth+4\tabcolsep+2\arrayrulewidth\relax}|}{Windows 11 (10.0.26200); conda environment fyp-local: Python 3.13.18, NumPy 2.5.4, SciPy 1.18.1, scikit-learn 1.9.1, pandas 3.1.1, OpenCV 4.11.0, PyTorch 2.9.0+cu126. CPU only; no GPU, dataset or model weights needed unless stated. Run from the repository root: python -m unittest discover -s tests -t .} \\ \hline
\textbf{Assumptions:} & \multicolumn{3}{L{\dimexpr 0.63\linewidth+4\tabcolsep+2\arrayrulewidth\relax}|}{Synthetic students with t\_\allowbreak{}last 120 s. Track A has ratings 3, 2, 4, 3 at 0, 20, 50, 80 s. A 'late' labeler gives the same ratings 2 s later.} \\ \hline
\textbf{Pre-Requisite:} & \multicolumn{3}{L{\dimexpr 0.63\linewidth+4\tabcolsep+2\arrayrulewidth\relax}|}{None beyond the Python environment.} \\ \hline
\textbf{Step No.} & \textbf{Execution description} & \multicolumn{2}{L{\dimexpr 0.36\linewidth+2\tabcolsep+1\arrayrulewidth\relax}|}{\textbf{Procedure result}} \\ \hline
1 & \texttt{test\_\allowbreak{}identical\_\allowbreak{}labelers\_\allowbreak{}predict\_\allowbreak{}each\_\allowbreak{}other\_\allowbreak{}exactly}: Build ceiling rows for a student where labeler\_\allowbreak{}01 and labeler\_\allowbreak{}02 both have track A. Stack the 'step|0' predictions of entry rows with stack\_\allowbreak{}pairs. & \multicolumn{2}{L{\dimexpr 0.36\linewidth+2\tabcolsep+1\arrayrulewidth\relax}|}{The predictions equal the targets exactly. (Passed)} \\ \hline
2 & \texttt{test\_\allowbreak{}a\_\allowbreak{}late\_\allowbreak{}labeler\_\allowbreak{}is\_\allowbreak{}recovered\_\allowbreak{}by\_\allowbreak{}the\_\allowbreak{}matching\_\allowbreak{}lag}: Build ceiling rows with labeler\_\allowbreak{}01 = A and labeler\_\allowbreak{}02 = A shifted 2 s later. Take entry rows with target labeler\_\allowbreak{}01. & \multicolumn{2}{L{\dimexpr 0.36\linewidth+2\tabcolsep+1\arrayrulewidth\relax}|}{At lag 0 the absolute error summed over entries is above 0. The column 'labeler\_\allowbreak{}02|step|2' equals the targets exactly, and so does the best match within +/\allowbreak{}-2 s ('labeler\_\allowbreak{}02|tol2'). (Passed)} \\ \hline
3 & \texttt{test\_\allowbreak{}run\_\allowbreak{}midpoints\_\allowbreak{}are\_\allowbreak{}unaffected\_\allowbreak{}by\_\allowbreak{}a\_\allowbreak{}small\_\allowbreak{}shift}: With the same 2 s shift, take midpoint rows with target labeler\_\allowbreak{}01. & \multicolumn{2}{L{\dimexpr 0.36\linewidth+2\tabcolsep+1\arrayrulewidth\relax}|}{The 'labeler\_\allowbreak{}02|step|0' values equal the targets exactly. (Passed)} \\ \hline
4 & \texttt{test\_\allowbreak{}gauss\_\allowbreak{}puts\_\allowbreak{}half\_\allowbreak{}the\_\allowbreak{}jump\_\allowbreak{}at\_\allowbreak{}the\_\allowbreak{}entry}: Interpolate entries 2 at 0 s and 4 at 10 s in gauss mode at t = 10 s. & \multicolumn{2}{L{\dimexpr 0.36\linewidth+2\tabcolsep+1\arrayrulewidth\relax}|}{Returns 3.0, half way through the jump. (Passed)} \\ \hline
5 & \texttt{test\_\allowbreak{}session\_\allowbreak{}start\_\allowbreak{}is\_\allowbreak{}excluded\_\allowbreak{}and\_\allowbreak{}self\_\allowbreak{}consistency\_\allowbreak{}columns}: Build ceiling rows for two labelers with track A and take labeler\_\allowbreak{}01's entry rows. & \multicolumn{2}{L{\dimexpr 0.36\linewidth+2\tabcolsep+1\arrayrulewidth\relax}|}{Entry times are [20, 50, 80]; the t = 0 entry is the session start and is excluded. own\_\allowbreak{}prev is [3, 2, 4]. own\_\allowbreak{}loro for the first entry is (3*20 + 4*30 + 3*40) /\allowbreak{} 90, the time-weighted mean of the other runs. (Passed)} \\ \hline
6 & \texttt{test\_\allowbreak{}median\_\allowbreak{}and\_\allowbreak{}mean\_\allowbreak{}of\_\allowbreak{}other\_\allowbreak{}experts\_\allowbreak{}skip\_\allowbreak{}the\_\allowbreak{}target\_\allowbreak{}and\_\allowbreak{}self}: Build ceiling rows for a student with labeler\_\allowbreak{}01 = A, labeler\_\allowbreak{}02 constant 2, labeler\_\allowbreak{}03 constant 4 and self\_\allowbreak{}labeling constant 5. Call others on labeler\_\allowbreak{}01's entry rows with 'step|0' and 'mean'. & \multicolumn{2}{L{\dimexpr 0.36\linewidth+2\tabcolsep+1\arrayrulewidth\relax}|}{Returns [3.0, 3.0, 3.0], the mean of 2 and 4. The target and the self-rating are not included. (Passed)} \\ \hline
\multicolumn{4}{|L{\dimexpr 0.84\linewidth+6\tabcolsep+3\arrayrulewidth\relax}|}{\textbf{Comments:} test\_\allowbreak{}gauss\_\allowbreak{}puts\_\allowbreak{}half\_\allowbreak{}the\_\allowbreak{}jump\_\allowbreak{}at\_\allowbreak{}the\_\allowbreak{}entry tests src/\allowbreak{}labels/\allowbreak{}interp.py directly; it sits here because the ceiling uses gauss curves.} \\ \hline
\multicolumn{4}{|C{\dimexpr 0.84\linewidth+6\tabcolsep+3\arrayrulewidth\relax}|}{\textbf{\fbox{Passed} \quad Failed \quad Not Executed}} \\ \hline
\end{longtable}

\begin{longtable}{|L{0.21\linewidth}|L{0.27\linewidth}|L{0.17\linewidth}|L{0.19\linewidth}|}
\caption{Test case TC-15: Agreement on coarse units}\label{tab:tc-15} \\ \hline
\multicolumn{4}{|C{\dimexpr 0.84\linewidth+6\tabcolsep+3\arrayrulewidth\relax}|}{\textbf{Agreement on coarse units}} \\ \hline
\multicolumn{4}{|C{\dimexpr 0.84\linewidth+6\tabcolsep+3\arrayrulewidth\relax}|}{\textbf{scripts/\allowbreak{}labels/\allowbreak{}label\_\allowbreak{}agreement.py; tests/\allowbreak{}test\_\allowbreak{}label\_\allowbreak{}agreement.py (TestCoarse)}} \\ \hline
\textbf{Test Case ID:} & TC-15 & \textbf{QA Test Engineer:} & Abi \\ \hline
\textbf{Test case Version:} & 2 & \textbf{Reviewed By:} & - \\ \hline
\textbf{Test Date:} & 2026-11-18 & \textbf{Use Case Reference(s):} & docs/\allowbreak{}experiments/\allowbreak{}06\_\allowbreak{}full\_\allowbreak{}pipeline\_\allowbreak{}and\_\allowbreak{}labels.md Part B \\ \hline
\textbf{Revision History:} & \multicolumn{3}{L{\dimexpr 0.63\linewidth+4\tabcolsep+2\arrayrulewidth\relax}|}{v1 2026-11-18 (5f5458c); v2 2026-11-18 (e664d52)} \\ \hline
\textbf{Objective:} & \multicolumn{3}{L{\dimexpr 0.63\linewidth+4\tabcolsep+2\arrayrulewidth\relax}|}{Check that agreement on student x activity means is high when labelers differ only in timing, and that the labeler list has six entries.} \\ \hline
\textbf{Product/Ver/ Module:} & \multicolumn{3}{L{\dimexpr 0.63\linewidth+4\tabcolsep+2\arrayrulewidth\relax}|}{FYP-1 attention pipeline, commit e743616. Label agreement analysis script, coarse part: agreement on per-student session means and k-minute means.} \\ \hline
\textbf{Environment:} & \multicolumn{3}{L{\dimexpr 0.63\linewidth+4\tabcolsep+2\arrayrulewidth\relax}|}{Windows 11 (10.0.26200); conda environment fyp-local: Python 3.13.18, NumPy 2.5.4, SciPy 1.18.1, scikit-learn 1.9.1, pandas 3.1.1, OpenCV 4.11.0, PyTorch 2.9.0+cu126. CPU only; no GPU, dataset or model weights needed unless stated. Run from the repository root: python -m unittest discover -s tests -t .} \\ \hline
\textbf{Assumptions:} & \multicolumn{3}{L{\dimexpr 0.63\linewidth+4\tabcolsep+2\arrayrulewidth\relax}|}{30 synthetic students (seed 2), t\_\allowbreak{}last 300 s. Each has a base rating in 2-4 with ratings base, base+1, base, base at 0, 30, 60, 90 s. labeler\_\allowbreak{}01 and labeler\_\allowbreak{}03 use these times; labeler\_\allowbreak{}02 enters the same ratings 3 s later.} \\ \hline
\textbf{Pre-Requisite:} & \multicolumn{3}{L{\dimexpr 0.63\linewidth+4\tabcolsep+2\arrayrulewidth\relax}|}{None beyond the Python environment.} \\ \hline
\textbf{Step No.} & \textbf{Execution description} & \multicolumn{2}{L{\dimexpr 0.36\linewidth+2\tabcolsep+1\arrayrulewidth\relax}|}{\textbf{Procedure result}} \\ \hline
1 & \texttt{test\_\allowbreak{}coarse\_\allowbreak{}means\_\allowbreak{}agree\_\allowbreak{}when\_\allowbreak{}labelers\_\allowbreak{}differ\_\allowbreak{}only\_\allowbreak{}in\_\allowbreak{}timing}: Call coarse on the 30 students and read student\_\allowbreak{}x\_\allowbreak{}activity, overall, experts\_\allowbreak{}pairwise. & \multicolumn{2}{L{\dimexpr 0.36\linewidth+2\tabcolsep+1\arrayrulewidth\relax}|}{The mean pairwise Spearman correlation between experts is above 0.95. LABELERS has length 6 (five experts plus self\_\allowbreak{}labeling). (Passed)} \\ \hline
\multicolumn{4}{|L{\dimexpr 0.84\linewidth+6\tabcolsep+3\arrayrulewidth\relax}|}{\textbf{Comments:} -} \\ \hline
\multicolumn{4}{|C{\dimexpr 0.84\linewidth+6\tabcolsep+3\arrayrulewidth\relax}|}{\textbf{\fbox{Passed} \quad Failed \quad Not Executed}} \\ \hline
\end{longtable}

\begin{longtable}{|L{0.21\linewidth}|L{0.27\linewidth}|L{0.17\linewidth}|L{0.19\linewidth}|}
\caption{Test case TC-16: Rating curve interpolation}\label{tab:tc-16} \\ \hline
\multicolumn{4}{|C{\dimexpr 0.84\linewidth+6\tabcolsep+3\arrayrulewidth\relax}|}{\textbf{Rating curve interpolation}} \\ \hline
\multicolumn{4}{|C{\dimexpr 0.84\linewidth+6\tabcolsep+3\arrayrulewidth\relax}|}{\textbf{src/\allowbreak{}labels/\allowbreak{}interp.py; tests/\allowbreak{}test\_\allowbreak{}label\_\allowbreak{}targets.py (TestInterpolation)}} \\ \hline
\textbf{Test Case ID:} & TC-16 & \textbf{QA Test Engineer:} & Abi \\ \hline
\textbf{Test case Version:} & 4 & \textbf{Reviewed By:} & - \\ \hline
\textbf{Test Date:} & 2026-11-18 & \textbf{Use Case Reference(s):} & docs/\allowbreak{}misc/\allowbreak{}targets.md \S{}2.1, \S{}7; docs/\allowbreak{}guide/\allowbreak{}13\_\allowbreak{}from\_\allowbreak{}ratings\_\allowbreak{}to\_\allowbreak{}numbers.md \\ \hline
\textbf{Revision History:} & \multicolumn{3}{L{\dimexpr 0.63\linewidth+4\tabcolsep+2\arrayrulewidth\relax}|}{v1 2026-11-16 (1106a77); v2 2026-11-18 (5f5458c); v3 2026-11-18 (e664d52); v4 2026-11-18 (e743616)} \\ \hline
\textbf{Objective:} & \multicolumn{3}{L{\dimexpr 0.63\linewidth+4\tabcolsep+2\arrayrulewidth\relax}|}{Check how a labeler's change events become a value at any time and a mean over a window, for every interpolation mode.} \\ \hline
\textbf{Product/Ver/ Module:} & \multicolumn{3}{L{\dimexpr 0.63\linewidth+4\tabcolsep+2\arrayrulewidth\relax}|}{FYP-1 attention pipeline, commit e743616. Label side of the per-labeler training path: it turns rating entries into training targets and agreement values.} \\ \hline
\textbf{Environment:} & \multicolumn{3}{L{\dimexpr 0.63\linewidth+4\tabcolsep+2\arrayrulewidth\relax}|}{Windows 11 (10.0.26200); conda environment fyp-local: Python 3.13.18, NumPy 2.5.4, SciPy 1.18.1, scikit-learn 1.9.1, pandas 3.1.1, OpenCV 4.11.0, PyTorch 2.9.0+cu126. CPU only; no GPU, dataset or model weights needed unless stated. Run from the repository root: python -m unittest discover -s tests -t .} \\ \hline
\textbf{Assumptions:} & \multicolumn{3}{L{\dimexpr 0.63\linewidth+4\tabcolsep+2\arrayrulewidth\relax}|}{One fixed synthetic labeler with entries T = [0, 10, 11, 30] s and ratings V = [3, 2, 3, 4].} \\ \hline
\textbf{Pre-Requisite:} & \multicolumn{3}{L{\dimexpr 0.63\linewidth+4\tabcolsep+2\arrayrulewidth\relax}|}{None beyond the Python environment.} \\ \hline
\textbf{Step No.} & \textbf{Execution description} & \multicolumn{2}{L{\dimexpr 0.36\linewidth+2\tabcolsep+1\arrayrulewidth\relax}|}{\textbf{Procedure result}} \\ \hline
1 & \texttt{test\_\allowbreak{}step\_\allowbreak{}holds\_\allowbreak{}and\_\allowbreak{}is\_\allowbreak{}nan\_\allowbreak{}before\_\allowbreak{}the\_\allowbreak{}first\_\allowbreak{}entry}: Call interpolate(T, V, t) in step mode at t = [-1, 0, 5, 10, 10.5, 40], and at the single time -1.0. & \multicolumn{2}{L{\dimexpr 0.36\linewidth+2\tabcolsep+1\arrayrulewidth\relax}|}{Values from t = 0 onward are [3, 3, 2, 2, 4]: each rating holds until the next entry and the last holds after the end. The value at -1.0 is NaN. (Passed)} \\ \hline
2 & \texttt{test\_\allowbreak{}every\_\allowbreak{}mode\_\allowbreak{}stays\_\allowbreak{}within\_\allowbreak{}1\_\allowbreak{}to\_\allowbreak{}5}: Evaluate the curve at 400 times from -1 to 45 s in each mode: step, ramp, linear, pchip, cubic and gauss. & \multicolumn{2}{L{\dimexpr 0.36\linewidth+2\tabcolsep+1\arrayrulewidth\relax}|}{In every mode all finite values lie in [1, 5]. (Passed)} \\ \hline
3 & \texttt{test\_\allowbreak{}window\_\allowbreak{}mean\_\allowbreak{}of\_\allowbreak{}the\_\allowbreak{}step}: Call window\_\allowbreak{}mean(T, V, 8, 12) in step mode. & \multicolumn{2}{L{\dimexpr 0.36\linewidth+2\tabcolsep+1\arrayrulewidth\relax}|}{Returns 2.75 (one second each at 3, 3, 2, 3). (Passed)} \\ \hline
4 & \texttt{test\_\allowbreak{}a\_\allowbreak{}centred\_\allowbreak{}window\_\allowbreak{}blends\_\allowbreak{}the\_\allowbreak{}previous\_\allowbreak{}value}: Call window\_\allowbreak{}mean on entries 2 at 0 s and 3 at 11 s over the window [7, 15), an 8 s window centred on the change. & \multicolumn{2}{L{\dimexpr 0.36\linewidth+2\tabcolsep+1\arrayrulewidth\relax}|}{Returns 2.5: half the window still holds the old value 2. (Passed)} \\ \hline
\multicolumn{4}{|L{\dimexpr 0.84\linewidth+6\tabcolsep+3\arrayrulewidth\relax}|}{\textbf{Comments:} -} \\ \hline
\multicolumn{4}{|C{\dimexpr 0.84\linewidth+6\tabcolsep+3\arrayrulewidth\relax}|}{\textbf{\fbox{Passed} \quad Failed \quad Not Executed}} \\ \hline
\end{longtable}

\begin{longtable}{|L{0.21\linewidth}|L{0.27\linewidth}|L{0.17\linewidth}|L{0.19\linewidth}|}
\caption{Test case TC-17: Rating conversion and rating metrics}\label{tab:tc-17} \\ \hline
\multicolumn{4}{|C{\dimexpr 0.84\linewidth+6\tabcolsep+3\arrayrulewidth\relax}|}{\textbf{Rating conversion and rating metrics}} \\ \hline
\multicolumn{4}{|C{\dimexpr 0.84\linewidth+6\tabcolsep+3\arrayrulewidth\relax}|}{\textbf{src/\allowbreak{}labels/\allowbreak{}interp.py, src/\allowbreak{}evaluation/\allowbreak{}metrics.py, src/\allowbreak{}core/\allowbreak{}classifiers.py, src/\allowbreak{}experts/\allowbreak{}expert.py, scripts/\allowbreak{}train\_\allowbreak{}models.py; tests/\allowbreak{}test\_\allowbreak{}label\_\allowbreak{}targets.py (TestRatings)}} \\ \hline
\textbf{Test Case ID:} & TC-17 & \textbf{QA Test Engineer:} & Abi \\ \hline
\textbf{Test case Version:} & 4 & \textbf{Reviewed By:} & - \\ \hline
\textbf{Test Date:} & 2026-11-18 & \textbf{Use Case Reference(s):} & docs/\allowbreak{}misc/\allowbreak{}targets.md \S{}3, \S{}7.1 items 2 and 7; Design \S{}11 \\ \hline
\textbf{Revision History:} & \multicolumn{3}{L{\dimexpr 0.63\linewidth+4\tabcolsep+2\arrayrulewidth\relax}|}{v1 2026-11-16 (1106a77); v2 2026-11-18 (5f5458c); v3 2026-11-18 (e664d52); v4 2026-11-18 (e743616)} \\ \hline
\textbf{Objective:} & \multicolumn{3}{L{\dimexpr 0.63\linewidth+4\tabcolsep+2\arrayrulewidth\relax}|}{Check the conversions between ratings 1-5 and soft labels 0-1, the rounding and majority rules, the per-label scoring metrics, the optional PCA step in classifiers, and per-labeler scoring.} \\ \hline
\textbf{Product/Ver/ Module:} & \multicolumn{3}{L{\dimexpr 0.63\linewidth+4\tabcolsep+2\arrayrulewidth\relax}|}{FYP-1 attention pipeline, commit e743616. Shared by training (soft targets, classifiers) and evaluation (rating\_\allowbreak{}metrics, c\_\allowbreak{}index, per-labeler scores) on the per-labeler path.} \\ \hline
\textbf{Environment:} & \multicolumn{3}{L{\dimexpr 0.63\linewidth+4\tabcolsep+2\arrayrulewidth\relax}|}{Windows 11 (10.0.26200); conda environment fyp-local: Python 3.13.18, NumPy 2.5.4, SciPy 1.18.1, scikit-learn 1.9.1, pandas 3.1.1, OpenCV 4.11.0, PyTorch 2.9.0+cu126. CPU only; no GPU, dataset or model weights needed unless stated. Run from the repository root: python -m unittest discover -s tests -t .} \\ \hline
\textbf{Assumptions:} & \multicolumn{3}{L{\dimexpr 0.63\linewidth+4\tabcolsep+2\arrayrulewidth\relax}|}{Inputs are small fixed arrays or seeded random data (numpy default\_\allowbreak{}rng seeds 0 and 1). The PCA check uses 12 columns of which 6 are near copies of the other 6.} \\ \hline
\textbf{Pre-Requisite:} & \multicolumn{3}{L{\dimexpr 0.63\linewidth+4\tabcolsep+2\arrayrulewidth\relax}|}{None beyond the Python environment (scikit-learn and SciPy installed).} \\ \hline
\textbf{Step No.} & \textbf{Execution description} & \multicolumn{2}{L{\dimexpr 0.36\linewidth+2\tabcolsep+1\arrayrulewidth\relax}|}{\textbf{Procedure result}} \\ \hline
1 & \texttt{test\_\allowbreak{}normalise\_\allowbreak{}and\_\allowbreak{}back\_\allowbreak{}clip\_\allowbreak{}out\_\allowbreak{}of\_\allowbreak{}range\_\allowbreak{}values}: Call normalise on ratings [0, 1, 3, 5, 7] and to\_\allowbreak{}rating on probabilities [-0.5, 0, 0.5, 1, 2]. & \multicolumn{2}{L{\dimexpr 0.36\linewidth+2\tabcolsep+1\arrayrulewidth\relax}|}{normalise gives [0, 0, 0.5, 1, 1] and to\_\allowbreak{}rating gives [1, 1, 3, 5, 5]. Out-of-range inputs are clipped. (Passed)} \\ \hline
2 & \texttt{test\_\allowbreak{}halves\_\allowbreak{}round\_\allowbreak{}toward\_\allowbreak{}3}: Call round\_\allowbreak{}rating on [2.5, 3.5, 2.49, 4.5, 1.5, 7] and quantise on [2.24, 2.25, 2.75, 3.25, 3.3] with step 0.5. & \multicolumn{2}{L{\dimexpr 0.36\linewidth+2\tabcolsep+1\arrayrulewidth\relax}|}{round\_\allowbreak{}rating gives [3, 3, 2, 4, 2, 5]. quantise gives [2, 2.5, 3, 3, 3.5]. Exact halves go toward 3 and values above 5 are clipped. (Passed)} \\ \hline
3 & \texttt{test\_\allowbreak{}majority\_\allowbreak{}ties\_\allowbreak{}and\_\allowbreak{}resolution}: Call majority\_\allowbreak{}rating on rows [3, 3, 2, 4, 2, 5] and [2, 2, 4, 4, NaN, NaN]; on [2.4, 2.6, 2.5, 3.1, 4.0, 4.1] at step 0.5 and at step 1.0; and on a row of only NaN. & \multicolumn{2}{L{\dimexpr 0.36\linewidth+2\tabcolsep+1\arrayrulewidth\relax}|}{The two rows give [3, 2]: a tie goes to the value closest to 3, then to the lower value. The third row gives 2.5 at step 0.5 and 3.0 at step 1.0. The all-NaN row gives NaN. (Passed)} \\ \hline
4 & \texttt{test\_\allowbreak{}rating\_\allowbreak{}metrics}: Call rating\_\allowbreak{}metrics with tau\_\allowbreak{}att 3 on labels y = [2, 3, 3, 4, 2, 5, 3, 1] against y itself, against a constant 3, and against a constant 2.91 with 4e-16 of rounding noise on one element. & \multicolumn{2}{L{\dimexpr 0.36\linewidth+2\tabcolsep+1\arrayrulewidth\relax}|}{Perfect prediction: mae 0.0, accuracy 1.0, auc 1.0. Constant 3: auc 0.5, balanced\_\allowbreak{}acc 0.5, spearman NaN. Near-constant with noise: spearman and pearson are both NaN, so rounding noise is not treated as variation. (Passed)} \\ \hline
5 & \texttt{test\_\allowbreak{}expert\_\allowbreak{}pca\_\allowbreak{}step}: Build 400 x 12 random features where columns 6-11 copy columns 0-5 plus small noise, with y = (column 0 > 0). Build a logreg classifier without PCA and one with pca=0.99. Fit SklearnExpert('logreg', pca=3) on soft labels 0.1 /\allowbreak{} 0.9. & \multicolumn{2}{L{\dimexpr 0.36\linewidth+2\tabcolsep+1\arrayrulewidth\relax}|}{The pipeline without PCA has no 'pca' step. With pca=0.99 the fitted PCA keeps at most 7 components, so the duplicated columns collapse. The soft-label expert keeps exactly 3 components and predict\_\allowbreak{}proba returns shape (400,). (Passed)} \\ \hline
6 & \texttt{test\_\allowbreak{}scores\_\allowbreak{}per\_\allowbreak{}perspective}: Call score\_\allowbreak{}with\_\allowbreak{}perspectives on 5 prediction rows for entries a, a, b, c, d (ratings 2, 2, 3, 4, 3; labelers l1, l1, l1, l2, l2) with predictions [2, 3, 3, 4, 2] and tau 3. & \multicolumn{2}{L{\dimexpr 0.36\linewidth+2\tabcolsep+1\arrayrulewidth\relax}|}{Scores are stored under 'k', 'k/\allowbreak{}by\_\allowbreak{}labeler/\allowbreak{}l1' and 'k/\allowbreak{}by\_\allowbreak{}labeler/\allowbreak{}l2'. The pooled key counts 4 entries (entry a once). l1 has mae 0.25 (a predicted 2.5 against 2, b exact). l2 has mae 0.5 (c exact, d off by 1). (Passed)} \\ \hline
7 & \texttt{test\_\allowbreak{}c\_\allowbreak{}index}: Call c\_\allowbreak{}index on y = [1, 2, 2, 3, 3, 4, 5] with predictions p = [0.1, 0.2, 0.3, 0.25, 0.4, 0.5, 0.9], with a constant 3, with y itself and with -y. & \multicolumn{2}{L{\dimexpr 0.36\linewidth+2\tabcolsep+1\arrayrulewidth\relax}|}{c\_\allowbreak{}index(y, p) equals a brute-force pairwise concordance (ties count half). A constant gives 0.5, y gives 1.0 and -y gives 0.0. (Passed)} \\ \hline
8 & \texttt{test\_\allowbreak{}spread\_\allowbreak{}metrics\_\allowbreak{}decompose\_\allowbreak{}the\_\allowbreak{}squared\_\allowbreak{}error}: Draw 400 integer ratings 1-5. Score a shrunk noisy prediction and a constant 2.8, then a constant 3.0, then the MSE-optimal linear rescaling of a prediction. & \multicolumn{2}{L{\dimexpr 0.36\linewidth+2\tabcolsep+1\arrayrulewidth\relax}|}{For both first predictions mse\_\allowbreak{}bias2 + mse\_\allowbreak{}spread + mse\_\allowbreak{}floor equals rmse squared. The constant 3.0 has pred\_\allowbreak{}sd 0.0, share\_\allowbreak{}pred\_\allowbreak{}3 1.0, calib\_\allowbreak{}ratio NaN and pred\_\allowbreak{}shares summing to 1. The optimal prediction has calib\_\allowbreak{}ratio 1.0 (6 decimal places) and mse\_\allowbreak{}spread 0.0 (9 decimal places). (Passed)} \\ \hline
\multicolumn{4}{|L{\dimexpr 0.84\linewidth+6\tabcolsep+3\arrayrulewidth\relax}|}{\textbf{Comments:} test\_\allowbreak{}expert\_\allowbreak{}pca\_\allowbreak{}step and test\_\allowbreak{}scores\_\allowbreak{}per\_\allowbreak{}perspective sit in this class but test the classifier builder and scripts/\allowbreak{}train\_\allowbreak{}models.py respectively.} \\ \hline
\multicolumn{4}{|C{\dimexpr 0.84\linewidth+6\tabcolsep+3\arrayrulewidth\relax}|}{\textbf{\fbox{Passed} \quad Failed \quad Not Executed}} \\ \hline
\end{longtable}

\begin{longtable}{|L{0.21\linewidth}|L{0.27\linewidth}|L{0.17\linewidth}|L{0.19\linewidth}|}
\caption{Test case TC-18: Soft-label training of engagement experts}\label{tab:tc-18} \\ \hline
\multicolumn{4}{|C{\dimexpr 0.84\linewidth+6\tabcolsep+3\arrayrulewidth\relax}|}{\textbf{Soft-label training of engagement experts}} \\ \hline
\multicolumn{4}{|C{\dimexpr 0.84\linewidth+6\tabcolsep+3\arrayrulewidth\relax}|}{\textbf{src/\allowbreak{}experts/\allowbreak{}expert.py, src/\allowbreak{}core/\allowbreak{}classifiers.py; tests/\allowbreak{}test\_\allowbreak{}label\_\allowbreak{}targets.py (TestSoftLabels)}} \\ \hline
\textbf{Test Case ID:} & TC-18 & \textbf{QA Test Engineer:} & Abi \\ \hline
\textbf{Test case Version:} & 4 & \textbf{Reviewed By:} & - \\ \hline
\textbf{Test Date:} & 2026-11-18 & \textbf{Use Case Reference(s):} & docs/\allowbreak{}misc/\allowbreak{}targets.md \S{}7.1 item 7; Design \S{}8 \\ \hline
\textbf{Revision History:} & \multicolumn{3}{L{\dimexpr 0.63\linewidth+4\tabcolsep+2\arrayrulewidth\relax}|}{v1 2026-11-16 (1106a77); v2 2026-11-18 (5f5458c); v3 2026-11-18 (e664d52); v4 2026-11-18 (e743616)} \\ \hline
\textbf{Objective:} & \multicolumn{3}{L{\dimexpr 0.63\linewidth+4\tabcolsep+2\arrayrulewidth\relax}|}{Check that SklearnExpert learns from soft targets in [0, 1]: logistic regression directly, SVM and k-NN by thresholding at 0.5, and a constant target is predicted unchanged.} \\ \hline
\textbf{Product/Ver/ Module:} & \multicolumn{3}{L{\dimexpr 0.63\linewidth+4\tabcolsep+2\arrayrulewidth\relax}|}{FYP-1 attention pipeline, commit e743616. Engagement experts inside the mixture of experts, used on the per-labeler path where targets are normalised ratings.} \\ \hline
\textbf{Environment:} & \multicolumn{3}{L{\dimexpr 0.63\linewidth+4\tabcolsep+2\arrayrulewidth\relax}|}{Windows 11 (10.0.26200); conda environment fyp-local: Python 3.13.18, NumPy 2.5.4, SciPy 1.18.1, scikit-learn 1.9.1, pandas 3.1.1, OpenCV 4.11.0, PyTorch 2.9.0+cu126. CPU only; no GPU, dataset or model weights needed unless stated. Run from the repository root: python -m unittest discover -s tests -t .} \\ \hline
\textbf{Assumptions:} & \multicolumn{3}{L{\dimexpr 0.63\linewidth+4\tabcolsep+2\arrayrulewidth\relax}|}{Synthetic data: 1500 rows of 3 Gaussian features (seed 0). The true probability is sigmoid(1.5 * x0) and the soft target adds Gaussian noise with SD 0.05, clipped to [0, 1].} \\ \hline
\textbf{Pre-Requisite:} & \multicolumn{3}{L{\dimexpr 0.63\linewidth+4\tabcolsep+2\arrayrulewidth\relax}|}{None beyond the Python environment.} \\ \hline
\textbf{Step No.} & \textbf{Execution description} & \multicolumn{2}{L{\dimexpr 0.36\linewidth+2\tabcolsep+1\arrayrulewidth\relax}|}{\textbf{Procedure result}} \\ \hline
1 & \texttt{test\_\allowbreak{}logistic\_\allowbreak{}regression\_\allowbreak{}learns\_\allowbreak{}the\_\allowbreak{}soft\_\allowbreak{}target}: Fit SklearnExpert('logreg') on the soft targets and predict on the same rows. & \multicolumn{2}{L{\dimexpr 0.36\linewidth+2\tabcolsep+1\arrayrulewidth\relax}|}{The mean prediction equals the mean target to 2 decimal places, and the correlation with the true probability is above 0.99. (Passed)} \\ \hline
2 & \texttt{test\_\allowbreak{}svm\_\allowbreak{}and\_\allowbreak{}knn\_\allowbreak{}threshold\_\allowbreak{}soft\_\allowbreak{}labels\_\allowbreak{}and\_\allowbreak{}keep\_\allowbreak{}the\_\allowbreak{}sign}: Fit SklearnExpert('svm') and SklearnExpert('knn'), each with max\_\allowbreak{}train\_\allowbreak{}rows=600, on the soft targets and predict on the same rows. & \multicolumn{2}{L{\dimexpr 0.36\linewidth+2\tabcolsep+1\arrayrulewidth\relax}|}{For both models the correlation of the prediction with the true probability is above 0.5. (Passed)} \\ \hline
3 & \texttt{test\_\allowbreak{}a\_\allowbreak{}constant\_\allowbreak{}soft\_\allowbreak{}target\_\allowbreak{}is\_\allowbreak{}predicted\_\allowbreak{}as\_\allowbreak{}is}: Fit SklearnExpert('logreg') on 5 rows all with target 0.4 and predict on 2 rows. & \multicolumn{2}{L{\dimexpr 0.36\linewidth+2\tabcolsep+1\arrayrulewidth\relax}|}{The first prediction is exactly 0.4. (Passed)} \\ \hline
\multicolumn{4}{|L{\dimexpr 0.84\linewidth+6\tabcolsep+3\arrayrulewidth\relax}|}{\textbf{Comments:} -} \\ \hline
\multicolumn{4}{|C{\dimexpr 0.84\linewidth+6\tabcolsep+3\arrayrulewidth\relax}|}{\textbf{\fbox{Passed} \quad Failed \quad Not Executed}} \\ \hline
\end{longtable}

\begin{longtable}{|L{0.21\linewidth}|L{0.27\linewidth}|L{0.17\linewidth}|L{0.19\linewidth}|}
\caption{Test case TC-19: Per-labeler models and rating output}\label{tab:tc-19} \\ \hline
\multicolumn{4}{|C{\dimexpr 0.84\linewidth+6\tabcolsep+3\arrayrulewidth\relax}|}{\textbf{Per-labeler models and rating output}} \\ \hline
\multicolumn{4}{|C{\dimexpr 0.84\linewidth+6\tabcolsep+3\arrayrulewidth\relax}|}{\textbf{src/\allowbreak{}experts/\allowbreak{}per\_\allowbreak{}labeler.py, src/\allowbreak{}labels/\allowbreak{}interp.py, src/\allowbreak{}output/\allowbreak{}graph.py; tests/\allowbreak{}test\_\allowbreak{}label\_\allowbreak{}targets.py (TestPerLabeler)}} \\ \hline
\textbf{Test Case ID:} & TC-19 & \textbf{QA Test Engineer:} & Abi \\ \hline
\textbf{Test case Version:} & 4 & \textbf{Reviewed By:} & - \\ \hline
\textbf{Test Date:} & 2026-11-18 & \textbf{Use Case Reference(s):} & docs/\allowbreak{}misc/\allowbreak{}targets.md \S{}7, \S{}7.1 items 8-9 \\ \hline
\textbf{Revision History:} & \multicolumn{3}{L{\dimexpr 0.63\linewidth+4\tabcolsep+2\arrayrulewidth\relax}|}{v1 2026-11-16 (1106a77); v2 2026-11-18 (5f5458c); v3 2026-11-18 (e664d52); v4 2026-11-18 (e743616)} \\ \hline
\textbf{Objective:} & \multicolumn{3}{L{\dimexpr 0.63\linewidth+4\tabcolsep+2\arrayrulewidth\relax}|}{Check that per-labeler models each learn only their own labeler's targets, that the pooled variant answers the same for all labelers, that window curves are built correctly in both prediction modes, and that the rating graph records are well formed.} \\ \hline
\textbf{Product/Ver/ Module:} & \multicolumn{3}{L{\dimexpr 0.63\linewidth+4\tabcolsep+2\arrayrulewidth\relax}|}{FYP-1 attention pipeline, commit e743616. Per-labeler training and prediction (engagement.labelers = 'per\_\allowbreak{}labeler' or 'pooled') and the final retention graph output.} \\ \hline
\textbf{Environment:} & \multicolumn{3}{L{\dimexpr 0.63\linewidth+4\tabcolsep+2\arrayrulewidth\relax}|}{Windows 11 (10.0.26200); conda environment fyp-local: Python 3.13.18, NumPy 2.5.4, SciPy 1.18.1, scikit-learn 1.9.1, pandas 3.1.1, OpenCV 4.11.0, PyTorch 2.9.0+cu126. CPU only; no GPU, dataset or model weights needed unless stated. Run from the repository root: python -m unittest discover -s tests -t .} \\ \hline
\textbf{Assumptions:} & \multicolumn{3}{L{\dimexpr 0.63\linewidth+4\tabcolsep+2\arrayrulewidth\relax}|}{600 synthetic rows with 2 Gaussian features (seed 1), all in experiment '01', cycling through labeler\_\allowbreak{}01, labeler\_\allowbreak{}02 and self\_\allowbreak{}labeling. labeler\_\allowbreak{}01 and self\_\allowbreak{}labeling have target sigmoid(2 * x0); labeler\_\allowbreak{}02 has the opposite target. Each mixture is a shared, soft ExpertMixture of logistic regression, with gate input all ones.} \\ \hline
\textbf{Pre-Requisite:} & \multicolumn{3}{L{\dimexpr 0.63\linewidth+4\tabcolsep+2\arrayrulewidth\relax}|}{None beyond the Python environment. setUp fits the per-labeler model.} \\ \hline
\textbf{Step No.} & \textbf{Execution description} & \multicolumn{2}{L{\dimexpr 0.36\linewidth+2\tabcolsep+1\arrayrulewidth\relax}|}{\textbf{Procedure result}} \\ \hline
1 & \texttt{test\_\allowbreak{}each\_\allowbreak{}labeler\_\allowbreak{}model\_\allowbreak{}learns\_\allowbreak{}only\_\allowbreak{}its\_\allowbreak{}own\_\allowbreak{}labels}: Predict all rows with the per-labeler model. & \multicolumn{2}{L{\dimexpr 0.36\linewidth+2\tabcolsep+1\arrayrulewidth\relax}|}{labeler\_\allowbreak{}01's prediction correlates with x0 above 0.9. labeler\_\allowbreak{}02's correlates below -0.9. (Passed)} \\ \hline
2 & \texttt{test\_\allowbreak{}a\_\allowbreak{}pooled\_\allowbreak{}model\_\allowbreak{}answers\_\allowbreak{}the\_\allowbreak{}same\_\allowbreak{}for\_\allowbreak{}every\_\allowbreak{}labeler}: Fit PerLabelerModel with pooled=True on the same rows and predict all rows. & \multicolumn{2}{L{\dimexpr 0.36\linewidth+2\tabcolsep+1\arrayrulewidth\relax}|}{The output has keys labeler\_\allowbreak{}01, labeler\_\allowbreak{}02 and self\_\allowbreak{}labeling. labeler\_\allowbreak{}01 and labeler\_\allowbreak{}02 predictions are identical. labeler\_\allowbreak{}02's prediction correlates with x0 above 0.5, so the majority view wins. (Passed)} \\ \hline
3 & \texttt{test\_\allowbreak{}window\_\allowbreak{}curves\_\allowbreak{}in\_\allowbreak{}both\_\allowbreak{}modes}: Group rows into windows of 4 (150 windows) and add one more window with no rows. Call window\_\allowbreak{}curves in 'per\_\allowbreak{}labeler' and 'majority' mode with resolution 0.5. & \multicolumn{2}{L{\dimexpr 0.36\linewidth+2\tabcolsep+1\arrayrulewidth\relax}|}{Curves are returned for all three labelers. The per\_\allowbreak{}labeler final curve equals the median of the labeler curves. Both final curves lie in [1, 5]. The empty last window of labeler\_\allowbreak{}01's curve equals to\_\allowbreak{}rating of that labeler's training fill. (Passed)} \\ \hline
4 & \texttt{test\_\allowbreak{}rating\_\allowbreak{}graph\_\allowbreak{}records}: Call build\_\allowbreak{}rating\_\allowbreak{}graph for two windows with final ratings [3.2, 2.4], labeler curves labeler\_\allowbreak{}01 = [3, 2] and self\_\allowbreak{}labeling = [4, 3], tau\_\allowbreak{}att 3.0, mode 'per\_\allowbreak{}labeler'. & \multicolumn{2}{L{\dimexpr 0.36\linewidth+2\tabcolsep+1\arrayrulewidth\relax}|}{engaged is 'Y' then 'N'. The second record's labeler\_\allowbreak{}ratings are \{labeler\_\allowbreak{}01: 2.0, self\_\allowbreak{}labeling: 3.0\}. No record has a 'tid' key, so no track identity reaches the output. (Passed)} \\ \hline
\multicolumn{4}{|L{\dimexpr 0.84\linewidth+6\tabcolsep+3\arrayrulewidth\relax}|}{\textbf{Comments:} The 'no tid' check supports the privacy rule that the retention graph carries only window-level values.} \\ \hline
\multicolumn{4}{|C{\dimexpr 0.84\linewidth+6\tabcolsep+3\arrayrulewidth\relax}|}{\textbf{\fbox{Passed} \quad Failed \quad Not Executed}} \\ \hline
\end{longtable}

\begin{longtable}{|L{0.21\linewidth}|L{0.27\linewidth}|L{0.17\linewidth}|L{0.19\linewidth}|}
\caption{Test case TC-20: Joining rating entries to track features}\label{tab:tc-20} \\ \hline
\multicolumn{4}{|C{\dimexpr 0.84\linewidth+6\tabcolsep+3\arrayrulewidth\relax}|}{\textbf{Joining rating entries to track features}} \\ \hline
\multicolumn{4}{|C{\dimexpr 0.84\linewidth+6\tabcolsep+3\arrayrulewidth\relax}|}{\textbf{src/\allowbreak{}experts/\allowbreak{}per\_\allowbreak{}labeler.py; tests/\allowbreak{}test\_\allowbreak{}label\_\allowbreak{}targets.py (TestJoinEntries)}} \\ \hline
\textbf{Test Case ID:} & TC-20 & \textbf{QA Test Engineer:} & Abi \\ \hline
\textbf{Test case Version:} & 4 & \textbf{Reviewed By:} & - \\ \hline
\textbf{Test Date:} & 2026-11-18 & \textbf{Use Case Reference(s):} & docs/\allowbreak{}misc/\allowbreak{}targets.md \S{}7 \\ \hline
\textbf{Revision History:} & \multicolumn{3}{L{\dimexpr 0.63\linewidth+4\tabcolsep+2\arrayrulewidth\relax}|}{v1 2026-11-16 (1106a77); v2 2026-11-18 (5f5458c); v3 2026-11-18 (e664d52); v4 2026-11-18 (e743616)} \\ \hline
\textbf{Objective:} & \multicolumn{3}{L{\dimexpr 0.63\linewidth+4\tabcolsep+2\arrayrulewidth\relax}|}{Check that each rating entry is joined only to feature rows of tracks linked to the rated student in the entry's window.} \\ \hline
\textbf{Product/Ver/ Module:} & \multicolumn{3}{L{\dimexpr 0.63\linewidth+4\tabcolsep+2\arrayrulewidth\relax}|}{FYP-1 attention pipeline, commit e743616. Builds per-labeler training rows from label entries and label-centred window features.} \\ \hline
\textbf{Environment:} & \multicolumn{3}{L{\dimexpr 0.63\linewidth+4\tabcolsep+2\arrayrulewidth\relax}|}{Windows 11 (10.0.26200); conda environment fyp-local: Python 3.13.18, NumPy 2.5.4, SciPy 1.18.1, scikit-learn 1.9.1, pandas 3.1.1, OpenCV 4.11.0, PyTorch 2.9.0+cu126. CPU only; no GPU, dataset or model weights needed unless stated. Run from the repository root: python -m unittest discover -s tests -t .} \\ \hline
\textbf{Assumptions:} & \multicolumn{3}{L{\dimexpr 0.63\linewidth+4\tabcolsep+2\arrayrulewidth\relax}|}{Three synthetic entries and three synthetic feature rows; the track-to-student map is given directly.} \\ \hline
\textbf{Pre-Requisite:} & \multicolumn{3}{L{\dimexpr 0.63\linewidth+4\tabcolsep+2\arrayrulewidth\relax}|}{None beyond the Python environment.} \\ \hline
\textbf{Step No.} & \textbf{Execution description} & \multicolumn{2}{L{\dimexpr 0.36\linewidth+2\tabcolsep+1\arrayrulewidth\relax}|}{\textbf{Procedure result}} \\ \hline
1 & \texttt{test\_\allowbreak{}entries\_\allowbreak{}join\_\allowbreak{}their\_\allowbreak{}students\_\allowbreak{}linked\_\allowbreak{}tracks\_\allowbreak{}only}: Call join\_\allowbreak{}entries with entries (s1, labeler\_\allowbreak{}01, 10 s), (s1, self\_\allowbreak{}labeling, 10 s), (s2, labeler\_\allowbreak{}01, 20 s); feature rows for tracks 7 and 8 at 10 s and track 9 at 20 s; and the map \{7: s1, 9: s3\}. & \multicolumn{2}{L{\dimexpr 0.36\linewidth+2\tabcolsep+1\arrayrulewidth\relax}|}{Two rows are returned, both with tid 7 and sid s1. Track 8 is unlinked and dropped, and s2's entry has no linked track and is dropped. (Passed)} \\ \hline
\multicolumn{4}{|L{\dimexpr 0.84\linewidth+6\tabcolsep+3\arrayrulewidth\relax}|}{\textbf{Comments:} -} \\ \hline
\multicolumn{4}{|C{\dimexpr 0.84\linewidth+6\tabcolsep+3\arrayrulewidth\relax}|}{\textbf{\fbox{Passed} \quad Failed \quad Not Executed}} \\ \hline
\end{longtable}

\begin{longtable}{|L{0.21\linewidth}|L{0.27\linewidth}|L{0.17\linewidth}|L{0.19\linewidth}|}
\caption{Test case TC-21: Attention label parsing}\label{tab:tc-21} \\ \hline
\multicolumn{4}{|C{\dimexpr 0.84\linewidth+6\tabcolsep+3\arrayrulewidth\relax}|}{\textbf{Attention label parsing}} \\ \hline
\multicolumn{4}{|C{\dimexpr 0.84\linewidth+6\tabcolsep+3\arrayrulewidth\relax}|}{\textbf{src/\allowbreak{}labels/\allowbreak{}attention.py, src/\allowbreak{}core/\allowbreak{}clock.py; tests/\allowbreak{}test\_\allowbreak{}labels.py (TestAttentionTrack)}} \\ \hline
\textbf{Test Case ID:} & TC-21 & \textbf{QA Test Engineer:} & Abi \\ \hline
\textbf{Test case Version:} & 1 & \textbf{Reviewed By:} & - \\ \hline
\textbf{Test Date:} & 2026-11-18 & \textbf{Use Case Reference(s):} & Design \S{}8.1 step 1 \\ \hline
\textbf{Revision History:} & \multicolumn{3}{L{\dimexpr 0.63\linewidth+4\tabcolsep+2\arrayrulewidth\relax}|}{v1 2026-11-14 (178c9bc)} \\ \hline
\textbf{Objective:} & \multicolumn{3}{L{\dimexpr 0.63\linewidth+4\tabcolsep+2\arrayrulewidth\relax}|}{Check that a labeler's list of change events is turned into a correct attention step function, and that the wall-clock parser reads both timestamp formats. Covers clock parsing, event filtering, the time-weighted window mean, the unlabelled period and the clock offset.} \\ \hline
\textbf{Product/Ver/ Module:} & \multicolumn{3}{L{\dimexpr 0.63\linewidth+4\tabcolsep+2\arrayrulewidth\relax}|}{FYP-1 attention pipeline, commit e743616. First stage of the label side of the training pipeline: it reads DIPSER labeler JSON events before window and per-labeler targets are built.} \\ \hline
\textbf{Environment:} & \multicolumn{3}{L{\dimexpr 0.63\linewidth+4\tabcolsep+2\arrayrulewidth\relax}|}{Windows 11 (10.0.26200); conda environment fyp-local: Python 3.13.18, NumPy 2.5.4, SciPy 1.18.1, scikit-learn 1.9.1, pandas 3.1.1, OpenCV 4.11.0, PyTorch 2.9.0+cu126. CPU only; no GPU, dataset or model weights needed unless stated. Run from the repository root: python -m unittest discover -s tests -t .} \\ \hline
\textbf{Assumptions:} & \multicolumn{3}{L{\dimexpr 0.63\linewidth+4\tabcolsep+2\arrayrulewidth\relax}|}{Synthetic event lists built by a helper stand in for DIPSER labeler files. Every event has emotion '1'; the attention key is added only where the test gives a value. All events fall in the minute after 10:00:00 (T0).} \\ \hline
\textbf{Pre-Requisite:} & \multicolumn{3}{L{\dimexpr 0.63\linewidth+4\tabcolsep+2\arrayrulewidth\relax}|}{None beyond the Python environment.} \\ \hline
\textbf{Step No.} & \textbf{Execution description} & \multicolumn{2}{L{\dimexpr 0.36\linewidth+2\tabcolsep+1\arrayrulewidth\relax}|}{\textbf{Procedure result}} \\ \hline
1 & \texttt{test\_\allowbreak{}clock}: Call parse\_\allowbreak{}clock on the label format '10:40:43:047895' and the frame-filename format '10\_\allowbreak{}40\_\allowbreak{}43\_\allowbreak{}4'. & \multicolumn{2}{L{\dimexpr 0.36\linewidth+2\tabcolsep+1\arrayrulewidth\relax}|}{Returns seconds since midnight: 10*3600 + 40*60 + 43.047895 for the first and 10*3600 + 40*60 + 43.4 for the second. The fractional part is read with its own number of digits. (Passed)} \\ \hline
2 & \texttt{test\_\allowbreak{}emotion\_\allowbreak{}only\_\allowbreak{}entries\_\allowbreak{}do\_\allowbreak{}not\_\allowbreak{}change\_\allowbreak{}attention}: Call parse\_\allowbreak{}events on three events at T0, T0+5 s and T0+10 s with attention 3, no attention key, and 5. & \multicolumn{2}{L{\dimexpr 0.36\linewidth+2\tabcolsep+1\arrayrulewidth\relax}|}{The track values are [3, 5]. The emotion-only entry is skipped and does not add a value. (Passed)} \\ \hline
3 & \texttt{test\_\allowbreak{}time\_\allowbreak{}weighted\_\allowbreak{}mean}: Parse events with attention 2 at T0 and 4 at T0+4 s. Call mean\_\allowbreak{}over on [T0, T0+8) and on [T0+6, T0+10). & \multicolumn{2}{L{\dimexpr 0.36\linewidth+2\tabcolsep+1\arrayrulewidth\relax}|}{The first mean is 3.0 (4 s at 2 and 4 s at 4). The second is 4.0, because the last value holds after the last entry. (Passed)} \\ \hline
4 & \texttt{test\_\allowbreak{}unlabeled\_\allowbreak{}before\_\allowbreak{}first\_\allowbreak{}entry}: Parse a single event with attention 3 at T0+5 s. Call mean\_\allowbreak{}over on [T0, T0+8). & \multicolumn{2}{L{\dimexpr 0.36\linewidth+2\tabcolsep+1\arrayrulewidth\relax}|}{Returns NaN, because the window starts before the labeler's first entry. (Passed)} \\ \hline
5 & \texttt{test\_\allowbreak{}out\_\allowbreak{}of\_\allowbreak{}range\_\allowbreak{}values\_\allowbreak{}dropped}: Parse events with attention 3, 7 and 2 at T0, T0+5 s and T0+10 s. & \multicolumn{2}{L{\dimexpr 0.36\linewidth+2\tabcolsep+1\arrayrulewidth\relax}|}{The track values are [3, 2]. The value 7 is outside 1-5 and is dropped. (Passed)} \\ \hline
6 & \texttt{test\_\allowbreak{}offset}: Parse one event at T0 with offset\_\allowbreak{}s = 2.0. & \multicolumn{2}{L{\dimexpr 0.36\linewidth+2\tabcolsep+1\arrayrulewidth\relax}|}{The first track time equals T0 + 2.0 s, so the offset shifts labels onto the camera clock. (Passed)} \\ \hline
\multicolumn{4}{|L{\dimexpr 0.84\linewidth+6\tabcolsep+3\arrayrulewidth\relax}|}{\textbf{Comments:} The 'value 7' case mirrors a real DIPSER entry noted in the parse\_\allowbreak{}events docstring (group\_\allowbreak{}01/\allowbreak{}experiment\_\allowbreak{}03/\allowbreak{}subject\_\allowbreak{}18 self\_\allowbreak{}labeling).} \\ \hline
\multicolumn{4}{|C{\dimexpr 0.84\linewidth+6\tabcolsep+3\arrayrulewidth\relax}|}{\textbf{\fbox{Passed} \quad Failed \quad Not Executed}} \\ \hline
\end{longtable}

\begin{longtable}{|L{0.21\linewidth}|L{0.27\linewidth}|L{0.17\linewidth}|L{0.19\linewidth}|}
\caption{Test case TC-22: Window engagement labels}\label{tab:tc-22} \\ \hline
\multicolumn{4}{|C{\dimexpr 0.84\linewidth+6\tabcolsep+3\arrayrulewidth\relax}|}{\textbf{Window engagement labels}} \\ \hline
\multicolumn{4}{|C{\dimexpr 0.84\linewidth+6\tabcolsep+3\arrayrulewidth\relax}|}{\textbf{src/\allowbreak{}labels/\allowbreak{}window\_\allowbreak{}labels.py, src/\allowbreak{}labels/\allowbreak{}attention.py; tests/\allowbreak{}test\_\allowbreak{}labels.py (TestWindowLabeler)}} \\ \hline
\textbf{Test Case ID:} & TC-22 & \textbf{QA Test Engineer:} & Abi \\ \hline
\textbf{Test case Version:} & 1 & \textbf{Reviewed By:} & - \\ \hline
\textbf{Test Date:} & 2026-11-18 & \textbf{Use Case Reference(s):} & Design \S{}8.1 steps 3-5 \\ \hline
\textbf{Revision History:} & \multicolumn{3}{L{\dimexpr 0.63\linewidth+4\tabcolsep+2\arrayrulewidth\relax}|}{v1 2026-11-14 (178c9bc)} \\ \hline
\textbf{Objective:} & \multicolumn{3}{L{\dimexpr 0.63\linewidth+4\tabcolsep+2\arrayrulewidth\relax}|}{Check how per-student labeler tracks become a window label: the median consensus per student, the share of engaged students, the minimum number of labelled students, and the option to leave out the self-rating.} \\ \hline
\textbf{Product/Ver/ Module:} & \multicolumn{3}{L{\dimexpr 0.63\linewidth+4\tabcolsep+2\arrayrulewidth\relax}|}{FYP-1 attention pipeline, commit e743616. Label side of the training pipeline: produces the engaged /\allowbreak{} not engaged window labels and the share engaged used for training and evaluation.} \\ \hline
\textbf{Environment:} & \multicolumn{3}{L{\dimexpr 0.63\linewidth+4\tabcolsep+2\arrayrulewidth\relax}|}{Windows 11 (10.0.26200); conda environment fyp-local: Python 3.13.18, NumPy 2.5.4, SciPy 1.18.1, scikit-learn 1.9.1, pandas 3.1.1, OpenCV 4.11.0, PyTorch 2.9.0+cu126. CPU only; no GPU, dataset or model weights needed unless stated. Run from the repository root: python -m unittest discover -s tests -t .} \\ \hline
\textbf{Assumptions:} & \multicolumn{3}{L{\dimexpr 0.63\linewidth+4\tabcolsep+2\arrayrulewidth\relax}|}{Each synthetic labeler gives one constant rating from T0, so the window mean equals that rating. Labeler names follow the DIPSER pattern (labeler\_\allowbreak{}00, labeler\_\allowbreak{}01, ..., self\_\allowbreak{}labeling).} \\ \hline
\textbf{Pre-Requisite:} & \multicolumn{3}{L{\dimexpr 0.63\linewidth+4\tabcolsep+2\arrayrulewidth\relax}|}{None beyond the Python environment.} \\ \hline
\textbf{Step No.} & \textbf{Execution description} & \multicolumn{2}{L{\dimexpr 0.36\linewidth+2\tabcolsep+1\arrayrulewidth\relax}|}{\textbf{Procedure result}} \\ \hline
1 & \texttt{test\_\allowbreak{}median\_\allowbreak{}consensus\_\allowbreak{}and\_\allowbreak{}share}: Label the window [T0, T0+8) for three students with labeler ratings s1 = (5, 4, 1), s2 = (2, 2, 5), s3 = (4, 4, 4), using WindowLabeler(tau\_\allowbreak{}att=4, tau\_\allowbreak{}share=0.5, min\_\allowbreak{}labeled\_\allowbreak{}students=3). & \multicolumn{2}{L{\dimexpr 0.36\linewidth+2\tabcolsep+1\arrayrulewidth\relax}|}{Student medians are 4, 2 and 4, so two of three students are engaged. engaged = 1, n\_\allowbreak{}labeled = 3 and share\_\allowbreak{}engaged = 2/\allowbreak{}3. (Passed)} \\ \hline
2 & \texttt{test\_\allowbreak{}too\_\allowbreak{}few\_\allowbreak{}students\_\allowbreak{}dropped}: Label the window [T0, T0+8) for one student rated 5, with min\_\allowbreak{}labeled\_\allowbreak{}students=3. & \multicolumn{2}{L{\dimexpr 0.36\linewidth+2\tabcolsep+1\arrayrulewidth\relax}|}{engaged is None (the window is dropped) and n\_\allowbreak{}labeled is 1. (Passed)} \\ \hline
3 & \texttt{test\_\allowbreak{}exclude\_\allowbreak{}self}: Compute consensus for one student with labeler\_\allowbreak{}01 = 2 and self\_\allowbreak{}labeling = 5 over [T0, T0+8), once with exclude\_\allowbreak{}self=False and once with exclude\_\allowbreak{}self=True. & \multicolumn{2}{L{\dimexpr 0.36\linewidth+2\tabcolsep+1\arrayrulewidth\relax}|}{Without exclusion the consensus is the median 3.5. With exclusion the self-rating is ignored and the consensus is 2.0. (Passed)} \\ \hline
\multicolumn{4}{|L{\dimexpr 0.84\linewidth+6\tabcolsep+3\arrayrulewidth\relax}|}{\textbf{Comments:} -} \\ \hline
\multicolumn{4}{|C{\dimexpr 0.84\linewidth+6\tabcolsep+3\arrayrulewidth\relax}|}{\textbf{\fbox{Passed} \quad Failed \quad Not Executed}} \\ \hline
\end{longtable}

\begin{longtable}{|L{0.21\linewidth}|L{0.27\linewidth}|L{0.17\linewidth}|L{0.19\linewidth}|}
\caption{Test case TC-23: Gated mixture of engagement experts}\label{tab:tc-23} \\ \hline
\multicolumn{4}{|C{\dimexpr 0.84\linewidth+6\tabcolsep+3\arrayrulewidth\relax}|}{\textbf{Gated mixture of engagement experts}} \\ \hline
\multicolumn{4}{|C{\dimexpr 0.84\linewidth+6\tabcolsep+3\arrayrulewidth\relax}|}{\textbf{src/\allowbreak{}experts/\allowbreak{}mixture.py, src/\allowbreak{}experts/\allowbreak{}expert.py, src/\allowbreak{}spu/\allowbreak{}context.py; tests/\allowbreak{}test\_\allowbreak{}mixture.py (TestMixture)}} \\ \hline
\textbf{Test Case ID:} & TC-23 & \textbf{QA Test Engineer:} & Abi \\ \hline
\textbf{Test case Version:} & 1 & \textbf{Reviewed By:} & - \\ \hline
\textbf{Test Date:} & 2026-11-18 & \textbf{Use Case Reference(s):} & Design \S{}8 (mixture), \S{}8.2 (hand-weighted baseline), \S{}8.3 (per-experiment training), \S{}7 (oracle context) \\ \hline
\textbf{Revision History:} & \multicolumn{3}{L{\dimexpr 0.63\linewidth+4\tabcolsep+2\arrayrulewidth\relax}|}{v1 2026-11-14 (178c9bc)} \\ \hline
\textbf{Objective:} & \multicolumn{3}{L{\dimexpr 0.63\linewidth+4\tabcolsep+2\arrayrulewidth\relax}|}{Check that ExpertMixture combines per-experiment expert probabilities correctly in soft, hard, oracle and shared configurations, that each expert is trained only on its own experiment's windows, and that the hand-weighted baseline applies the per-experiment phone-use sign.} \\ \hline
\textbf{Product/Ver/ Module:} & \multicolumn{3}{L{\dimexpr 0.63\linewidth+4\tabcolsep+2\arrayrulewidth\relax}|}{FYP-1 attention pipeline, commit e743616. Engagement-scoring stage: takes each window's feature vector (Vec) and the SPU context gate's class probabilities and returns P(window engaged).} \\ \hline
\textbf{Environment:} & \multicolumn{3}{L{\dimexpr 0.63\linewidth+4\tabcolsep+2\arrayrulewidth\relax}|}{Windows 11 (10.0.26200); conda environment fyp-local: Python 3.13.18, NumPy 2.5.4, SciPy 1.18.1, scikit-learn 1.9.1, pandas 3.1.1, OpenCV 4.11.0, PyTorch 2.9.0+cu126. CPU only; no GPU, dataset or model weights needed unless stated. Run from the repository root: python -m unittest discover -s tests -t .} \\ \hline
\textbf{Assumptions:} & \multicolumn{3}{L{\dimexpr 0.63\linewidth+4\tabcolsep+2\arrayrulewidth\relax}|}{Synthetic inputs only. Experts are replaced by a stub class 'Constant' that returns a fixed probability per experiment (01: 0.1, 02: 0.5, 03: 0.9, shared: 0.3), so the test checks the mixing arithmetic, not any learned model. Feature matrices are zeros where their values do not matter.} \\ \hline
\textbf{Pre-Requisite:} & \multicolumn{3}{L{\dimexpr 0.63\linewidth+4\tabcolsep+2\arrayrulewidth\relax}|}{None beyond the Python environment (numpy). No data files, model weights or GPU.} \\ \hline
\textbf{Step No.} & \textbf{Execution description} & \multicolumn{2}{L{\dimexpr 0.36\linewidth+2\tabcolsep+1\arrayrulewidth\relax}|}{\textbf{Procedure result}} \\ \hline
1 & \texttt{test\_\allowbreak{}soft}: Build a soft-mode, per-experiment mixture over classes ['01', '02', '03'] and call predict\_\allowbreak{}proba on 2 rows with gate probabilities g = [[0.5, 0.5, 0.0], [0.1, 0.2, 0.7]]. & \multicolumn{2}{L{\dimexpr 0.36\linewidth+2\tabcolsep+1\arrayrulewidth\relax}|}{Output equals the gate-weighted sum of expert probabilities: [0.5*0.1 + 0.5*0.5, 0.1*0.1 + 0.2*0.5 + 0.7*0.9] = [0.30, 0.74] (assert\_\allowbreak{}allclose). Observed: as expected. (Passed)} \\ \hline
2 & \texttt{test\_\allowbreak{}hard\_\allowbreak{}routes\_\allowbreak{}to\_\allowbreak{}argmax}: Same gate probabilities, mixture built in hard mode. & \multicolumn{2}{L{\dimexpr 0.36\linewidth+2\tabcolsep+1\arrayrulewidth\relax}|}{Each row takes the probability of the argmax expert only: [0.1, 0.9] (row 0 ties 01/\allowbreak{}02 and argmax picks the first, 01). Observed: as expected. (Passed)} \\ \hline
3 & \texttt{test\_\allowbreak{}oracle\_\allowbreak{}is\_\allowbreak{}one\_\allowbreak{}hot}: Build g with context\_\allowbreak{}probs('oracle', classes, ['03', '02']) and pass it to a soft-mode mixture. & \multicolumn{2}{L{\dimexpr 0.36\linewidth+2\tabcolsep+1\arrayrulewidth\relax}|}{The oracle gate is one-hot on the true experiment, so the output is the matching expert's value: [0.9, 0.5]. Observed: as expected. (Passed)} \\ \hline
4 & \texttt{test\_\allowbreak{}shared\_\allowbreak{}scope\_\allowbreak{}uses\_\allowbreak{}one\_\allowbreak{}expert}: Build a soft-mode mixture with scope 'shared' (one expert built with exp\_\allowbreak{}id None) and call predict\_\allowbreak{}proba with the same g. & \multicolumn{2}{L{\dimexpr 0.36\linewidth+2\tabcolsep+1\arrayrulewidth\relax}|}{Every class maps to the single shared expert, so both rows return 0.3 regardless of g. Observed: as expected. (Passed)} \\ \hline
5 & \texttt{test\_\allowbreak{}each\_\allowbreak{}expert\_\allowbreak{}trains\_\allowbreak{}on\_\allowbreak{}its\_\allowbreak{}own\_\allowbreak{}experiment}: Build a per-experiment mixture with a Recorder expert that logs the first feature column it receives in fit(). Call fit with X = [[1], [2], [3], [4]], y = [0, 1, 0, 1], experiment ids ['01', '02', '01', '03']. & \multicolumn{2}{L{\dimexpr 0.36\linewidth+2\tabcolsep+1\arrayrulewidth\relax}|}{Each expert sees only its own experiment's rows: \{'01': [1.0, 3.0], '02': [2.0], '03': [4.0]\}. Observed: as expected. (Passed)} \\ \hline
6 & \texttt{test\_\allowbreak{}hand\_\allowbreak{}weighted\_\allowbreak{}phone\_\allowbreak{}sign}: Build two HandWeightedExpert instances on features ['vec.still', 'vec.phone'] with weight vec.still = 1.0 and phone\_\allowbreak{}sign \{'01': +1, '02': -1\}: one for experiment 01 (phone on-task), one for 02 (phone off-task). Fit both on X = [[0, -1], [0, 1]] and predict on the same X. & \multicolumn{2}{L{\dimexpr 0.36\linewidth+2\tabcolsep+1\arrayrulewidth\relax}|}{For experiment 01, the row with higher phone use (row 1) gets a higher P(engaged) than row 0. For experiment 02 the order is reversed: row 1 gets a lower P(engaged). Observed: as expected. (Passed)} \\ \hline
\multicolumn{4}{|L{\dimexpr 0.84\linewidth+6\tabcolsep+3\arrayrulewidth\relax}|}{\textbf{Comments:} The helper class 'Constant' and the factory function 'mixture' are fixtures, not test cases. In the phone-sign step vec.phone has no explicit weight, so it takes the default 1.0 multiplied by the experiment's sign.} \\ \hline
\multicolumn{4}{|C{\dimexpr 0.84\linewidth+6\tabcolsep+3\arrayrulewidth\relax}|}{\textbf{\fbox{Passed} \quad Failed \quad Not Executed}} \\ \hline
\end{longtable}

\begin{longtable}{|L{0.21\linewidth}|L{0.27\linewidth}|L{0.17\linewidth}|L{0.19\linewidth}|}
\caption{Test case TC-24: Leave-one-group-out split}\label{tab:tc-24} \\ \hline
\multicolumn{4}{|C{\dimexpr 0.84\linewidth+6\tabcolsep+3\arrayrulewidth\relax}|}{\textbf{Leave-one-group-out split}} \\ \hline
\multicolumn{4}{|C{\dimexpr 0.84\linewidth+6\tabcolsep+3\arrayrulewidth\relax}|}{\textbf{src/\allowbreak{}spu/\allowbreak{}splits.py; tests/\allowbreak{}test\_\allowbreak{}mixture.py (TestLogo)}} \\ \hline
\textbf{Test Case ID:} & TC-24 & \textbf{QA Test Engineer:} & Abi \\ \hline
\textbf{Test case Version:} & 1 & \textbf{Reviewed By:} & - \\ \hline
\textbf{Test Date:} & 2026-11-18 & \textbf{Use Case Reference(s):} & Design \S{}7 \\ \hline
\textbf{Revision History:} & \multicolumn{3}{L{\dimexpr 0.63\linewidth+4\tabcolsep+2\arrayrulewidth\relax}|}{v1 2026-11-14 (178c9bc)} \\ \hline
\textbf{Objective:} & \multicolumn{3}{L{\dimexpr 0.63\linewidth+4\tabcolsep+2\arrayrulewidth\relax}|}{Check that logo\_\allowbreak{}splits yields one fold per group and that each fold's test set is exactly the held-out group, with none of that group in training.} \\ \hline
\textbf{Product/Ver/ Module:} & \multicolumn{3}{L{\dimexpr 0.63\linewidth+4\tabcolsep+2\arrayrulewidth\relax}|}{FYP-1 attention pipeline, commit e743616. Evaluation protocol used by the SPU gate and expert experiments, so that all cameras of a held-out group stay together in the test fold.} \\ \hline
\textbf{Environment:} & \multicolumn{3}{L{\dimexpr 0.63\linewidth+4\tabcolsep+2\arrayrulewidth\relax}|}{Windows 11 (10.0.26200); conda environment fyp-local: Python 3.13.18, NumPy 2.5.4, SciPy 1.18.1, scikit-learn 1.9.1, pandas 3.1.1, OpenCV 4.11.0, PyTorch 2.9.0+cu126. CPU only; no GPU, dataset or model weights needed unless stated. Run from the repository root: python -m unittest discover -s tests -t .} \\ \hline
\textbf{Assumptions:} & \multicolumn{3}{L{\dimexpr 0.63\linewidth+4\tabcolsep+2\arrayrulewidth\relax}|}{Synthetic group labels; no features or models involved.} \\ \hline
\textbf{Pre-Requisite:} & \multicolumn{3}{L{\dimexpr 0.63\linewidth+4\tabcolsep+2\arrayrulewidth\relax}|}{None beyond the Python environment (numpy).} \\ \hline
\textbf{Step No.} & \textbf{Execution description} & \multicolumn{2}{L{\dimexpr 0.36\linewidth+2\tabcolsep+1\arrayrulewidth\relax}|}{\textbf{Procedure result}} \\ \hline
1 & \texttt{test\_\allowbreak{}held\_\allowbreak{}out\_\allowbreak{}group\_\allowbreak{}is\_\allowbreak{}whole\_\allowbreak{}test\_\allowbreak{}fold}: Call logo\_\allowbreak{}splits on groups ['g1', 'g1', 'g2', 'g3', 'g2'] and collect the folds as (group, train\_\allowbreak{}idx, test\_\allowbreak{}idx). & \multicolumn{2}{L{\dimexpr 0.36\linewidth+2\tabcolsep+1\arrayrulewidth\relax}|}{Folds come in sorted group order ['g1', 'g2', 'g3']. In every fold all test rows belong to the held-out group and no training row does. Observed: as expected. (Passed)} \\ \hline
\multicolumn{4}{|L{\dimexpr 0.84\linewidth+6\tabcolsep+3\arrayrulewidth\relax}|}{\textbf{Comments:} -} \\ \hline
\multicolumn{4}{|C{\dimexpr 0.84\linewidth+6\tabcolsep+3\arrayrulewidth\relax}|}{\textbf{\fbox{Passed} \quad Failed \quad Not Executed}} \\ \hline
\end{longtable}

\begin{longtable}{|L{0.21\linewidth}|L{0.27\linewidth}|L{0.17\linewidth}|L{0.19\linewidth}|}
\caption{Test case TC-25: Per-labeler ensemble on the window grid}\label{tab:tc-25} \\ \hline
\multicolumn{4}{|C{\dimexpr 0.84\linewidth+6\tabcolsep+3\arrayrulewidth\relax}|}{\textbf{Per-labeler ensemble on the window grid}} \\ \hline
\multicolumn{4}{|C{\dimexpr 0.84\linewidth+6\tabcolsep+3\arrayrulewidth\relax}|}{\textbf{src/\allowbreak{}experts/\allowbreak{}levels.py, src/\allowbreak{}experts/\allowbreak{}mixture.py, src/\allowbreak{}experts/\allowbreak{}expert.py; tests/\allowbreak{}test\_\allowbreak{}per\_\allowbreak{}labeler.py (TestLabelerEnsemble)}} \\ \hline
\textbf{Test Case ID:} & TC-25 & \textbf{QA Test Engineer:} & Abi \\ \hline
\textbf{Test case Version:} & 1 & \textbf{Reviewed By:} & - \\ \hline
\textbf{Test Date:} & 2026-11-18 & \textbf{Use Case Reference(s):} & Design \S{}8.1; config engagement.labelers = 'ensemble' \\ \hline
\textbf{Revision History:} & \multicolumn{3}{L{\dimexpr 0.63\linewidth+4\tabcolsep+2\arrayrulewidth\relax}|}{v1 2026-11-15 (a497de6)} \\ \hline
\textbf{Objective:} & \multicolumn{3}{L{\dimexpr 0.63\linewidth+4\tabcolsep+2\arrayrulewidth\relax}|}{Check that the labeler ensemble trains one mixture per labeler (excluding the self-rating), and that at prediction a labeler votes only for students it rated, while all labelers vote when nothing is known. Also check that consensus mode fits a single mixture.} \\ \hline
\textbf{Product/Ver/ Module:} & \multicolumn{3}{L{\dimexpr 0.63\linewidth+4\tabcolsep+2\arrayrulewidth\relax}|}{FYP-1 attention pipeline, commit e743616. Engagement model fitting and window prediction in late aggregation (fit\_\allowbreak{}engagement, predict\_\allowbreak{}windows).} \\ \hline
\textbf{Environment:} & \multicolumn{3}{L{\dimexpr 0.63\linewidth+4\tabcolsep+2\arrayrulewidth\relax}|}{Windows 11 (10.0.26200); conda environment fyp-local: Python 3.13.18, NumPy 2.5.4, SciPy 1.18.1, scikit-learn 1.9.1, pandas 3.1.1, OpenCV 4.11.0, PyTorch 2.9.0+cu126. CPU only; no GPU, dataset or model weights needed unless stated. Run from the repository root: python -m unittest discover -s tests -t .} \\ \hline
\textbf{Assumptions:} & \multicolumn{3}{L{\dimexpr 0.63\linewidth+4\tabcolsep+2\arrayrulewidth\relax}|}{400 synthetic student rows (seed 0) with one feature x, grouped into 100 windows of 4. labeler\_\allowbreak{}01 and self\_\allowbreak{}labeling label a row engaged when x > 0; labeler\_\allowbreak{}02 gives the opposite. All rows are in experiment '01'. Windows 0-49 train and 50-99 test, with gate probabilities [1, 0]. The score is the correlation between predicted window scores and the true share of rows with x > 0.} \\ \hline
\textbf{Pre-Requisite:} & \multicolumn{3}{L{\dimexpr 0.63\linewidth+4\tabcolsep+2\arrayrulewidth\relax}|}{None beyond the Python environment.} \\ \hline
\textbf{Step No.} & \textbf{Execution description} & \multicolumn{2}{L{\dimexpr 0.36\linewidth+2\tabcolsep+1\arrayrulewidth\relax}|}{\textbf{Procedure result}} \\ \hline
1 & \texttt{test\_\allowbreak{}a\_\allowbreak{}labeler\_\allowbreak{}who\_\allowbreak{}did\_\allowbreak{}not\_\allowbreak{}rate\_\allowbreak{}a\_\allowbreak{}student\_\allowbreak{}does\_\allowbreak{}not\_\allowbreak{}vote}: Build rows where labeler\_\allowbreak{}02 rated no student (rated\_\allowbreak{}by false for labeler\_\allowbreak{}02). Fit with engagement config \{labelers: 'ensemble', ensemble\_\allowbreak{}exclude: ['self\_\allowbreak{}labeling']\} and predict the test windows. & \multicolumn{2}{L{\dimexpr 0.36\linewidth+2\tabcolsep+1\arrayrulewidth\relax}|}{The model is a LabelerEnsemble with mixtures for labeler\_\allowbreak{}01 and labeler\_\allowbreak{}02 only. The correlation with the true share is above 0.8, so only labeler\_\allowbreak{}01's view is used. (Passed)} \\ \hline
2 & \texttt{test\_\allowbreak{}both\_\allowbreak{}labelers\_\allowbreak{}vote\_\allowbreak{}where\_\allowbreak{}both\_\allowbreak{}rated}: Same as above, but labeler\_\allowbreak{}02 rated every student. & \multicolumn{2}{L{\dimexpr 0.36\linewidth+2\tabcolsep+1\arrayrulewidth\relax}|}{The ensemble has labeler\_\allowbreak{}01 and labeler\_\allowbreak{}02. The absolute correlation is below 0.5, because the two opposite views cancel. (Passed)} \\ \hline
3 & \texttt{test\_\allowbreak{}without\_\allowbreak{}labeler\_\allowbreak{}information\_\allowbreak{}every\_\allowbreak{}labeler\_\allowbreak{}votes}: Same as above, but rated\_\allowbreak{}by is None (as at inference). & \multicolumn{2}{L{\dimexpr 0.36\linewidth+2\tabcolsep+1\arrayrulewidth\relax}|}{The ensemble has labeler\_\allowbreak{}01 and labeler\_\allowbreak{}02. The absolute correlation is below 0.5, because all labelers are averaged. (Passed)} \\ \hline
4 & \texttt{test\_\allowbreak{}consensus\_\allowbreak{}mode\_\allowbreak{}fits\_\allowbreak{}one\_\allowbreak{}mixture}: Call fit\_\allowbreak{}engagement with \{labelers: 'consensus'\} on the same rows. & \multicolumn{2}{L{\dimexpr 0.36\linewidth+2\tabcolsep+1\arrayrulewidth\relax}|}{The returned model is a single ExpertMixture. (Passed)} \\ \hline
\multicolumn{4}{|L{\dimexpr 0.84\linewidth+6\tabcolsep+3\arrayrulewidth\relax}|}{\textbf{Comments:} -} \\ \hline
\multicolumn{4}{|C{\dimexpr 0.84\linewidth+6\tabcolsep+3\arrayrulewidth\relax}|}{\textbf{\fbox{Passed} \quad Failed \quad Not Executed}} \\ \hline
\end{longtable}

\begin{longtable}{|L{0.21\linewidth}|L{0.27\linewidth}|L{0.17\linewidth}|L{0.19\linewidth}|}
\caption{Test case TC-26: SCB student-behaviour feature extractor (BFMU c)}\label{tab:tc-26} \\ \hline
\multicolumn{4}{|C{\dimexpr 0.84\linewidth+6\tabcolsep+3\arrayrulewidth\relax}|}{\textbf{SCB student-behaviour feature extractor (BFMU c)}} \\ \hline
\multicolumn{4}{|C{\dimexpr 0.84\linewidth+6\tabcolsep+3\arrayrulewidth\relax}|}{\textbf{src/\allowbreak{}bfmu/\allowbreak{}scb.py, src/\allowbreak{}bfmu/\allowbreak{}base.py, src/\allowbreak{}tracking/\allowbreak{}association.py; tests/\allowbreak{}test\_\allowbreak{}scb.py (TestScbExtractor)}} \\ \hline
\textbf{Test Case ID:} & TC-26 & \textbf{QA Test Engineer:} & Abi \\ \hline
\textbf{Test case Version:} & 1 & \textbf{Reviewed By:} & - \\ \hline
\textbf{Test Date:} & 2026-11-18 & \textbf{Use Case Reference(s):} & Design \S{}5c \\ \hline
\textbf{Revision History:} & \multicolumn{3}{L{\dimexpr 0.63\linewidth+4\tabcolsep+2\arrayrulewidth\relax}|}{v1 2026-11-14 (178c9bc)} \\ \hline
\textbf{Objective:} & \multicolumn{3}{L{\dimexpr 0.63\linewidth+4\tabcolsep+2\arrayrulewidth\relax}|}{Check the parts of the SCB extractor that do not need the detector network: feature key naming, assignment of behaviour boxes and candidate scores to tracks, per-class confidence thresholds, and window feature computation, in both 'soft' and 'flags' modes.} \\ \hline
\textbf{Product/Ver/ Module:} & \multicolumn{3}{L{\dimexpr 0.63\linewidth+4\tabcolsep+2\arrayrulewidth\relax}|}{FYP-1 attention pipeline, commit e743616. Behaviour feature unit in the per-frame pass: turns SCB YOLOv7 behaviour detections (hand raising, read, write) into per-track, per-frame signals, then into window features for the engagement experts.} \\ \hline
\textbf{Environment:} & \multicolumn{3}{L{\dimexpr 0.63\linewidth+4\tabcolsep+2\arrayrulewidth\relax}|}{Windows 11 (10.0.26200); conda environment fyp-local: Python 3.13.18, NumPy 2.5.4, SciPy 1.18.1, scikit-learn 1.9.1, pandas 3.1.1, OpenCV 4.11.0, PyTorch 2.9.0+cu126. CPU only; no GPU, dataset or model weights needed unless stated. Run from the repository root: python -m unittest discover -s tests -t .} \\ \hline
\textbf{Assumptions:} & \multicolumn{3}{L{\dimexpr 0.63\linewidth+4\tabcolsep+2\arrayrulewidth\relax}|}{Synthetic boxes, scores and tracklets. The extractor is built with weights 'unused.pt' and classes ['hand\_\allowbreak{}raising', 'read', 'write']; the detector is never loaded, so the YOLOv7 inference path (frame\_\allowbreak{}signals) is not exercised here. min\_\allowbreak{}overlap is 0.5.} \\ \hline
\textbf{Pre-Requisite:} & \multicolumn{3}{L{\dimexpr 0.63\linewidth+4\tabcolsep+2\arrayrulewidth\relax}|}{None beyond the Python environment (numpy, scipy). No SCB weights and no GPU are needed, because the detectors load lazily on first use.} \\ \hline
\textbf{Step No.} & \textbf{Execution description} & \multicolumn{2}{L{\dimexpr 0.36\linewidth+2\tabcolsep+1\arrayrulewidth\relax}|}{\textbf{Procedure result}} \\ \hline
1 & \texttt{test\_\allowbreak{}keys\_\allowbreak{}known\_\allowbreak{}without\_\allowbreak{}loading\_\allowbreak{}weights}: Build one extractor in 'flags' mode and one in 'soft' mode from the registry and read their keys. & \multicolumn{2}{L{\dimexpr 0.36\linewidth+2\tabcolsep+1\arrayrulewidth\relax}|}{Flags mode keys are ('scb.hand\_\allowbreak{}raising', 'scb.read', 'scb.write'). Soft mode keys are each class plus its '.max' variant, then the objectness key: ('scb.hand\_\allowbreak{}raising', 'scb.hand\_\allowbreak{}raising.max', 'scb.read', 'scb.read.max', 'scb.write', 'scb.write.max', 'scb.m0.obj', 'scb.m0.obj.max'). \_\allowbreak{}detectors is still None. Observed: as expected. (Passed)} \\ \hline
2 & \texttt{test\_\allowbreak{}soft\_\allowbreak{}scores\_\allowbreak{}take\_\allowbreak{}max\_\allowbreak{}of\_\allowbreak{}candidates\_\allowbreak{}per\_\allowbreak{}track}: Two track boxes [0,0,100,200] and [120,0,220,200]. Four candidate boxes: two inside track 0, one inside track 1, one at [400,0,450,50] that overlaps no track. Each candidate has 4 scores. Call soft\_\allowbreak{}scores. & \multicolumn{2}{L{\dimexpr 0.36\linewidth+2\tabcolsep+1\arrayrulewidth\relax}|}{Track 0 gets the column-wise maximum of its two candidates, [0.2, 0.7, 0.3, 0.8]. Track 1 gets its single candidate's scores, [0.4, 0.0, 0.1, 0.5]. The outside candidate (scores 0.9) is discarded. Observed: as expected. (Passed)} \\ \hline
3 & \texttt{test\_\allowbreak{}soft\_\allowbreak{}candidate\_\allowbreak{}goes\_\allowbreak{}to\_\allowbreak{}one\_\allowbreak{}track\_\allowbreak{}only}: Two identical track boxes [0,0,100,200] (duplicate tracks on one student) and one candidate [10,10,50,50] with score 1.0. Call soft\_\allowbreak{}scores. & \multicolumn{2}{L{\dimexpr 0.36\linewidth+2\tabcolsep+1\arrayrulewidth\relax}|}{The output sums to 1.0, so the candidate is credited to exactly one track, not both. Observed: as expected. (Passed)} \\ \hline
4 & \texttt{test\_\allowbreak{}soft\_\allowbreak{}window\_\allowbreak{}features\_\allowbreak{}mean\_\allowbreak{}and\_\allowbreak{}peak}: Tracklet of 4 frames with signal scb.hand\_\allowbreak{}raising = [0, 0, 0.8, 0]. Call window\_\allowbreak{}features in soft mode with fps 5.0. & \multicolumn{2}{L{\dimexpr 0.36\linewidth+2\tabcolsep+1\arrayrulewidth\relax}|}{scb.hand\_\allowbreak{}raising = 0.2 (window mean) and scb.hand\_\allowbreak{}raising.max = 0.8 (window peak). Observed: as expected. (Passed)} \\ \hline
5 & \texttt{test\_\allowbreak{}normalize\_\allowbreak{}class}: Call normalize\_\allowbreak{}class on 'hand-raising', 'Read', 'Using Phone'. & \multicolumn{2}{L{\dimexpr 0.36\linewidth+2\tabcolsep+1\arrayrulewidth\relax}|}{Returns 'hand\_\allowbreak{}raising', 'read', 'using\_\allowbreak{}phone' (lower case, '-' and ' ' replaced by '\_\allowbreak{}'). Observed: as expected. (Passed)} \\ \hline
6 & \texttt{test\_\allowbreak{}behaviour\_\allowbreak{}box\_\allowbreak{}assigned\_\allowbreak{}to\_\allowbreak{}containing\_\allowbreak{}track}: Flags mode. Track boxes [0,0,100,200] and [120,0,220,200] with track ids [7, 9]. Behaviour boxes [130,40,200,120] (inside the second track) and [400,0,450,50] (outside both). Call match. & \multicolumn{2}{L{\dimexpr 0.36\linewidth+2\tabcolsep+1\arrayrulewidth\relax}|}{Returns [9]: only the track that contains the box is flagged; the outside box flags nobody. Observed: as expected. (Passed)} \\ \hline
7 & \texttt{test\_\allowbreak{}one\_\allowbreak{}box\_\allowbreak{}flags\_\allowbreak{}one\_\allowbreak{}track}: Flags mode. Two identical track boxes [0,0,100,200] with ids [1, 2] and one behaviour box [10,10,50,50]. Call match. & \multicolumn{2}{L{\dimexpr 0.36\linewidth+2\tabcolsep+1\arrayrulewidth\relax}|}{The returned list has length 1: one-to-one (Hungarian) matching flags one track, not both duplicates. Observed: as expected. (Passed)} \\ \hline
8 & \texttt{test\_\allowbreak{}class\_\allowbreak{}conf\_\allowbreak{}overrides\_\allowbreak{}default}: Build a flags-mode extractor whose model spec has conf 0.25 and class\_\allowbreak{}conf \{'hand\_\allowbreak{}raising': 0.6\}. Read thresholds. & \multicolumn{2}{L{\dimexpr 0.36\linewidth+2\tabcolsep+1\arrayrulewidth\relax}|}{thresholds == [[0.6, 0.25, 0.25]]: hand\_\allowbreak{}raising uses its own threshold, read and write fall back to conf. Observed: as expected. (Passed)} \\ \hline
9 & \texttt{test\_\allowbreak{}window\_\allowbreak{}feature\_\allowbreak{}is\_\allowbreak{}share\_\allowbreak{}of\_\allowbreak{}flagged\_\allowbreak{}frames}: Flags mode. Tracklet of 10 frames with scb.write = [1, 1, 1, 0, 0, 0, 0, 0, 0, 0] and no scb.read signal. Call window\_\allowbreak{}features with fps 5.0. & \multicolumn{2}{L{\dimexpr 0.36\linewidth+2\tabcolsep+1\arrayrulewidth\relax}|}{scb.write = 0.3 (3 of 10 frames flagged). scb.read is NaN because the signal is absent. Observed: as expected. (Passed)} \\ \hline
\multicolumn{4}{|L{\dimexpr 0.84\linewidth+6\tabcolsep+3\arrayrulewidth\relax}|}{\textbf{Comments:} The thresholds (conf, class\_\allowbreak{}conf) only apply in flags mode; soft mode uses no thresholds or NMS.} \\ \hline
\multicolumn{4}{|C{\dimexpr 0.84\linewidth+6\tabcolsep+3\arrayrulewidth\relax}|}{\textbf{\fbox{Passed} \quad Failed \quad Not Executed}} \\ \hline
\end{longtable}

\begin{longtable}{|L{0.21\linewidth}|L{0.27\linewidth}|L{0.17\linewidth}|L{0.19\linewidth}|}
\caption{Test case TC-27: SCB extractor wiring in the pipeline and frame-pass cache key}\label{tab:tc-27} \\ \hline
\multicolumn{4}{|C{\dimexpr 0.84\linewidth+6\tabcolsep+3\arrayrulewidth\relax}|}{\textbf{SCB extractor wiring in the pipeline and frame-pass cache key}} \\ \hline
\multicolumn{4}{|C{\dimexpr 0.84\linewidth+6\tabcolsep+3\arrayrulewidth\relax}|}{\textbf{src/\allowbreak{}core/\allowbreak{}pipeline.py, src/\allowbreak{}bfmu/\allowbreak{}scb.py, configs/\allowbreak{}deliverable3.json; tests/\allowbreak{}test\_\allowbreak{}scb.py (TestScbInPipeline)}} \\ \hline
\textbf{Test Case ID:} & TC-27 & \textbf{QA Test Engineer:} & Abi \\ \hline
\textbf{Test case Version:} & 1 & \textbf{Reviewed By:} & - \\ \hline
\textbf{Test Date:} & 2026-11-18 & \textbf{Use Case Reference(s):} & Design \S{}5c \\ \hline
\textbf{Revision History:} & \multicolumn{3}{L{\dimexpr 0.63\linewidth+4\tabcolsep+2\arrayrulewidth\relax}|}{v1 2026-11-14 (178c9bc)} \\ \hline
\textbf{Objective:} & \multicolumn{3}{L{\dimexpr 0.63\linewidth+4\tabcolsep+2\arrayrulewidth\relax}|}{Check that the Deliverable III configuration builds a pipeline containing the SCB extractor, and that SCB settings are part of the frame-pass (tracking) hash, so changing them forces a new frame pass instead of reusing a stale cache.} \\ \hline
\textbf{Product/Ver/ Module:} & \multicolumn{3}{L{\dimexpr 0.63\linewidth+4\tabcolsep+2\arrayrulewidth\relax}|}{FYP-1 attention pipeline, commit e743616. Pipeline construction and result caching: tracking\_\allowbreak{}hash keys the cached per-frame pass, which includes every BFMU unit that reads frames.} \\ \hline
\textbf{Environment:} & \multicolumn{3}{L{\dimexpr 0.63\linewidth+4\tabcolsep+2\arrayrulewidth\relax}|}{Windows 11 (10.0.26200); conda environment fyp-local: Python 3.13.18, NumPy 2.5.4, SciPy 1.18.1, scikit-learn 1.9.1, pandas 3.1.1, OpenCV 4.11.0, PyTorch 2.9.0+cu126. CPU only; no GPU, dataset or model weights needed unless stated. Run from the repository root: python -m unittest discover -s tests -t .} \\ \hline
\textbf{Assumptions:} & \multicolumn{3}{L{\dimexpr 0.63\linewidth+4\tabcolsep+2\arrayrulewidth\relax}|}{Uses the real configs/\allowbreak{}deliverable3.json, whose third bfmu entry (index 2) is the SCB extractor in soft mode with min\_\allowbreak{}overlap 0.5. No video is processed.} \\ \hline
\textbf{Pre-Requisite:} & \multicolumn{3}{L{\dimexpr 0.63\linewidth+4\tabcolsep+2\arrayrulewidth\relax}|}{configs/\allowbreak{}deliverable3.json present. No model weights or GPU: the pipeline builds its detector and SCB networks lazily.} \\ \hline
\textbf{Step No.} & \textbf{Execution description} & \multicolumn{2}{L{\dimexpr 0.36\linewidth+2\tabcolsep+1\arrayrulewidth\relax}|}{\textbf{Procedure result}} \\ \hline
1 & \texttt{test\_\allowbreak{}pipeline\_\allowbreak{}builds\_\allowbreak{}and\_\allowbreak{}scb\_\allowbreak{}is\_\allowbreak{}frame\_\allowbreak{}level}: Load configs/\allowbreak{}deliverable3.json and construct Pipeline(cfg). Load a second copy, set bfmu[2].min\_\allowbreak{}overlap to 0.6, and compute tracking\_\allowbreak{}hash for both. & \multicolumn{2}{L{\dimexpr 0.36\linewidth+2\tabcolsep+1\arrayrulewidth\relax}|}{'ScbBehaviour' is among the pipeline's extractor class names, and the two tracking hashes differ. Observed: as expected. (Passed)} \\ \hline
\multicolumn{4}{|L{\dimexpr 0.84\linewidth+6\tabcolsep+3\arrayrulewidth\relax}|}{\textbf{Comments:} The hash differs because ScbBehaviour has uses\_\allowbreak{}frames = True, so its config is included in the frame-pass hash.} \\ \hline
\multicolumn{4}{|C{\dimexpr 0.84\linewidth+6\tabcolsep+3\arrayrulewidth\relax}|}{\textbf{\fbox{Passed} \quad Failed \quad Not Executed}} \\ \hline
\end{longtable}

\begin{longtable}{|L{0.21\linewidth}|L{0.27\linewidth}|L{0.17\linewidth}|L{0.19\linewidth}|}
\caption{Test case TC-28: Safe YOLOv7 checkpoint loader}\label{tab:tc-28} \\ \hline
\multicolumn{4}{|C{\dimexpr 0.84\linewidth+6\tabcolsep+3\arrayrulewidth\relax}|}{\textbf{Safe YOLOv7 checkpoint loader}} \\ \hline
\multicolumn{4}{|C{\dimexpr 0.84\linewidth+6\tabcolsep+3\arrayrulewidth\relax}|}{\textbf{src/\allowbreak{}bfmu/\allowbreak{}yolov7.py; tests/\allowbreak{}test\_\allowbreak{}scb.py (TestSafeLoader)}} \\ \hline
\textbf{Test Case ID:} & TC-28 & \textbf{QA Test Engineer:} & Abi \\ \hline
\textbf{Test case Version:} & 1 & \textbf{Reviewed By:} & - \\ \hline
\textbf{Test Date:} & 2026-11-18 & \textbf{Use Case Reference(s):} & Design \S{}5c \\ \hline
\textbf{Revision History:} & \multicolumn{3}{L{\dimexpr 0.63\linewidth+4\tabcolsep+2\arrayrulewidth\relax}|}{v1 2026-11-14 (178c9bc)} \\ \hline
\textbf{Objective:} & \multicolumn{3}{L{\dimexpr 0.63\linewidth+4\tabcolsep+2\arrayrulewidth\relax}|}{Check that load\_\allowbreak{}yolov7 refuses a checkpoint that references a Python global outside its allow-list, so a tampered third-party pickle cannot run code.} \\ \hline
\textbf{Product/Ver/ Module:} & \multicolumn{3}{L{\dimexpr 0.63\linewidth+4\tabcolsep+2\arrayrulewidth\relax}|}{FYP-1 attention pipeline, commit e743616. Model-loading layer of the SCB behaviour detector; checkpoints are scanned for referenced globals before torch's weights\_\allowbreak{}only unpickler loads them.} \\ \hline
\textbf{Environment:} & \multicolumn{3}{L{\dimexpr 0.63\linewidth+4\tabcolsep+2\arrayrulewidth\relax}|}{Windows 11 (10.0.26200); conda environment fyp-local: Python 3.13.18, NumPy 2.5.4, SciPy 1.18.1, scikit-learn 1.9.1, pandas 3.1.1, OpenCV 4.11.0, PyTorch 2.9.0+cu126. CPU only; no GPU, dataset or model weights needed unless stated. Run from the repository root: python -m unittest discover -s tests -t .} \\ \hline
\textbf{Assumptions:} & \multicolumn{3}{L{\dimexpr 0.63\linewidth+4\tabcolsep+2\arrayrulewidth\relax}|}{The checkpoint is a temporary file created by the test: torch.save of \{'model': fractions.Fraction(1, 3)\}. fractions.Fraction is not in the ALLOWED prefixes.} \\ \hline
\textbf{Pre-Requisite:} & \multicolumn{3}{L{\dimexpr 0.63\linewidth+4\tabcolsep+2\arrayrulewidth\relax}|}{PyTorch installed and the YOLOv7 repository cloned at lib/\allowbreak{}yolov7 (lib/\allowbreak{}yolov7/\allowbreak{}models/\allowbreak{}yolo.py must exist); otherwise the class is skipped. No SCB weights or GPU needed. Writes only to a temporary directory that is deleted afterwards.} \\ \hline
\textbf{Step No.} & \textbf{Execution description} & \multicolumn{2}{L{\dimexpr 0.36\linewidth+2\tabcolsep+1\arrayrulewidth\relax}|}{\textbf{Procedure result}} \\ \hline
1 & \texttt{test\_\allowbreak{}refuses\_\allowbreak{}unexpected\_\allowbreak{}globals}: Save \{'model': Fraction(1, 3)\} to bad.pt in a temporary directory and call load\_\allowbreak{}yolov7 on it. & \multicolumn{2}{L{\dimexpr 0.36\linewidth+2\tabcolsep+1\arrayrulewidth\relax}|}{load\_\allowbreak{}yolov7 raises RuntimeError (refusing to load because of the unexpected global). Observed: as expected; the test ran and passed (lib/\allowbreak{}yolov7 was present, so it was not skipped). (Passed)} \\ \hline
\multicolumn{4}{|L{\dimexpr 0.84\linewidth+6\tabcolsep+3\arrayrulewidth\relax}|}{\textbf{Comments:} Conditional test: decorated with skipUnless on the presence of lib/\allowbreak{}yolov7/\allowbreak{}models/\allowbreak{}yolo.py.} \\ \hline
\multicolumn{4}{|C{\dimexpr 0.84\linewidth+6\tabcolsep+3\arrayrulewidth\relax}|}{\textbf{\fbox{Passed} \quad Failed \quad Not Executed}} \\ \hline
\end{longtable}

\begin{longtable}{|L{0.21\linewidth}|L{0.27\linewidth}|L{0.17\linewidth}|L{0.19\linewidth}|}
\caption{Test case TC-29: Within/\allowbreak{}between-cluster spread decomposition}\label{tab:tc-29} \\ \hline
\multicolumn{4}{|C{\dimexpr 0.84\linewidth+6\tabcolsep+3\arrayrulewidth\relax}|}{\textbf{Within/\allowbreak{}between-cluster spread decomposition}} \\ \hline
\multicolumn{4}{|C{\dimexpr 0.84\linewidth+6\tabcolsep+3\arrayrulewidth\relax}|}{\textbf{src/\allowbreak{}cfmu/\allowbreak{}spread.py; tests/\allowbreak{}test\_\allowbreak{}spread.py (TestSpread)}} \\ \hline
\textbf{Test Case ID:} & TC-29 & \textbf{QA Test Engineer:} & Abi \\ \hline
\textbf{Test case Version:} & 1 & \textbf{Reviewed By:} & - \\ \hline
\textbf{Test Date:} & 2026-11-18 & \textbf{Use Case Reference(s):} & Design \S{}10.4 \\ \hline
\textbf{Revision History:} & \multicolumn{3}{L{\dimexpr 0.63\linewidth+4\tabcolsep+2\arrayrulewidth\relax}|}{v1 2026-11-14 (178c9bc)} \\ \hline
\textbf{Objective:} & \multicolumn{3}{L{\dimexpr 0.63\linewidth+4\tabcolsep+2\arrayrulewidth\relax}|}{Check that spread\_\allowbreak{}decomposition splits the population variance of per-track descriptors exactly into within-cluster and between-cluster parts, and that the per-block between-cluster share behaves correctly in the limiting cases.} \\ \hline
\textbf{Product/Ver/ Module:} & \multicolumn{3}{L{\dimexpr 0.63\linewidth+4\tabcolsep+2\arrayrulewidth\relax}|}{FYP-1 attention pipeline, commit e743616. CFMU (cluster feature unit), part of FYP-1. Not active in Deliverable III; the maths is implemented and tested ahead of the clustering step (\S{}10.3), which is not built yet.} \\ \hline
\textbf{Environment:} & \multicolumn{3}{L{\dimexpr 0.63\linewidth+4\tabcolsep+2\arrayrulewidth\relax}|}{Windows 11 (10.0.26200); conda environment fyp-local: Python 3.13.18, NumPy 2.5.4, SciPy 1.18.1, scikit-learn 1.9.1, pandas 3.1.1, OpenCV 4.11.0, PyTorch 2.9.0+cu126. CPU only; no GPU, dataset or model weights needed unless stated. Run from the repository root: python -m unittest discover -s tests -t .} \\ \hline
\textbf{Assumptions:} & \multicolumn{3}{L{\dimexpr 0.63\linewidth+4\tabcolsep+2\arrayrulewidth\relax}|}{Synthetic inputs: x is a 20 x 4 matrix of standard-normal values and labels are random cluster ids in \{0, 1, 2\}, both from numpy default\_\allowbreak{}rng(0). Some steps replace the labels or data with degenerate cases.} \\ \hline
\textbf{Pre-Requisite:} & \multicolumn{3}{L{\dimexpr 0.63\linewidth+4\tabcolsep+2\arrayrulewidth\relax}|}{None beyond the Python environment (numpy).} \\ \hline
\textbf{Step No.} & \textbf{Execution description} & \multicolumn{2}{L{\dimexpr 0.36\linewidth+2\tabcolsep+1\arrayrulewidth\relax}|}{\textbf{Procedure result}} \\ \hline
1 & \texttt{test\_\allowbreak{}within\_\allowbreak{}plus\_\allowbreak{}between\_\allowbreak{}equals\_\allowbreak{}total\_\allowbreak{}population\_\allowbreak{}variance}: Call spread\_\allowbreak{}decomposition(x, labels) on the random 20 x 4 data with 3 clusters. & \multicolumn{2}{L{\dimexpr 0.36\linewidth+2\tabcolsep+1\arrayrulewidth\relax}|}{total\_\allowbreak{}sd squared equals the population variance x.var(axis=0) per feature (rtol 1e-10), and within\_\allowbreak{}sd squared plus between\_\allowbreak{}sd squared equals total\_\allowbreak{}sd squared. Observed: as expected. (Passed)} \\ \hline
2 & \texttt{test\_\allowbreak{}single\_\allowbreak{}cluster\_\allowbreak{}has\_\allowbreak{}no\_\allowbreak{}between\_\allowbreak{}spread}: Call spread\_\allowbreak{}decomposition with all 20 labels set to 0 and one block 'all' covering columns [0, 1, 2, 3]. & \multicolumn{2}{L{\dimexpr 0.36\linewidth+2\tabcolsep+1\arrayrulewidth\relax}|}{between\_\allowbreak{}sd is 0 for every feature and between\_\allowbreak{}share['all'] is exactly 0.0. Observed: as expected. (Passed)} \\ \hline
3 & \texttt{test\_\allowbreak{}singletons\_\allowbreak{}have\_\allowbreak{}no\_\allowbreak{}within\_\allowbreak{}spread}: Call spread\_\allowbreak{}decomposition with labels 0..19 (every track its own cluster) and block 'a' covering columns [0, 1]. & \multicolumn{2}{L{\dimexpr 0.36\linewidth+2\tabcolsep+1\arrayrulewidth\relax}|}{within\_\allowbreak{}sd is 0 for every feature (atol 1e-12) and between\_\allowbreak{}share['a'] is 1.0 (assertAlmostEqual). Observed: as expected. (Passed)} \\ \hline
4 & \texttt{test\_\allowbreak{}zero\_\allowbreak{}total\_\allowbreak{}spread\_\allowbreak{}gives\_\allowbreak{}zero\_\allowbreak{}share}: Call spread\_\allowbreak{}decomposition on a 5 x 2 matrix of ones, labels [0, 0, 1, 1, 2], block 'b' covering columns [0, 1]. & \multicolumn{2}{L{\dimexpr 0.36\linewidth+2\tabcolsep+1\arrayrulewidth\relax}|}{Total variance is zero, so between\_\allowbreak{}share['b'] is returned as 0.0 instead of a division by zero. Observed: as expected. (Passed)} \\ \hline
\multicolumn{4}{|L{\dimexpr 0.84\linewidth+6\tabcolsep+3\arrayrulewidth\relax}|}{\textbf{Comments:} -} \\ \hline
\multicolumn{4}{|C{\dimexpr 0.84\linewidth+6\tabcolsep+3\arrayrulewidth\relax}|}{\textbf{\fbox{Passed} \quad Failed \quad Not Executed}} \\ \hline
\end{longtable}

\begin{longtable}{|L{0.21\linewidth}|L{0.27\linewidth}|L{0.17\linewidth}|L{0.19\linewidth}|}
\caption{Test case TC-30: Windowing (aggregation)}\label{tab:tc-30} \\ \hline
\multicolumn{4}{|C{\dimexpr 0.84\linewidth+6\tabcolsep+3\arrayrulewidth\relax}|}{\textbf{Windowing (aggregation)}} \\ \hline
\multicolumn{4}{|C{\dimexpr 0.84\linewidth+6\tabcolsep+3\arrayrulewidth\relax}|}{\textbf{src/\allowbreak{}aggregation/\allowbreak{}windows.py; tests/\allowbreak{}test\_\allowbreak{}stages.py (TestWindows)}} \\ \hline
\textbf{Test Case ID:} & TC-30 & \textbf{QA Test Engineer:} & Abi \\ \hline
\textbf{Test case Version:} & 3 & \textbf{Reviewed By:} & - \\ \hline
\textbf{Test Date:} & 2026-11-18 & \textbf{Use Case Reference(s):} & - \\ \hline
\textbf{Revision History:} & \multicolumn{3}{L{\dimexpr 0.63\linewidth+4\tabcolsep+2\arrayrulewidth\relax}|}{v1 2026-11-14 (178c9bc); v2 2026-11-15 (be0a468); v3 2026-11-16 (d3ab703)} \\ \hline
\textbf{Objective:} & \multicolumn{3}{L{\dimexpr 0.63\linewidth+4\tabcolsep+2\arrayrulewidth\relax}|}{Check that a session's time span is cut into fixed-length windows with the configured fractional overlap, and that a span shorter than one window gives no windows.} \\ \hline
\textbf{Product/Ver/ Module:} & \multicolumn{3}{L{\dimexpr 0.63\linewidth+4\tabcolsep+2\arrayrulewidth\relax}|}{FYP-1 attention pipeline, commit e743616. Feature pass: splits a session into overlapping fixed-length windows before per-track features are aggregated into one vector per window.} \\ \hline
\textbf{Environment:} & \multicolumn{3}{L{\dimexpr 0.63\linewidth+4\tabcolsep+2\arrayrulewidth\relax}|}{Windows 11 (10.0.26200); conda environment fyp-local: Python 3.13.18, NumPy 2.5.4, SciPy 1.18.1, scikit-learn 1.9.1, pandas 3.1.1, OpenCV 4.11.0, PyTorch 2.9.0+cu126. CPU only; no GPU, dataset or model weights needed unless stated. Run from the repository root: python -m unittest discover -s tests -t .} \\ \hline
\textbf{Assumptions:} & \multicolumn{3}{L{\dimexpr 0.63\linewidth+4\tabcolsep+2\arrayrulewidth\relax}|}{Times are plain seconds. A trailing partial window is dropped by design.} \\ \hline
\textbf{Pre-Requisite:} & \multicolumn{3}{L{\dimexpr 0.63\linewidth+4\tabcolsep+2\arrayrulewidth\relax}|}{None beyond the Python environment (numpy).} \\ \hline
\textbf{Step No.} & \textbf{Execution description} & \multicolumn{2}{L{\dimexpr 0.36\linewidth+2\tabcolsep+1\arrayrulewidth\relax}|}{\textbf{Procedure result}} \\ \hline
1 & \texttt{test\_\allowbreak{}overlap}: Call make\_\allowbreak{}windows(t\_\allowbreak{}first=0, t\_\allowbreak{}last=16, length\_\allowbreak{}s=8, overlap=0.5). & \multicolumn{2}{L{\dimexpr 0.36\linewidth+2\tabcolsep+1\arrayrulewidth\relax}|}{Three windows are returned with (t\_\allowbreak{}start, t\_\allowbreak{}end) = (0, 8), (4, 12), (8, 16): a step of 4 s, as expected for 50\% overlap. (Passed)} \\ \hline
2 & \texttt{test\_\allowbreak{}too\_\allowbreak{}short}: Call make\_\allowbreak{}windows(0, 5, 8, 0.5): a 5 s span with 8 s windows. & \multicolumn{2}{L{\dimexpr 0.36\linewidth+2\tabcolsep+1\arrayrulewidth\relax}|}{An empty list is returned. (Passed)} \\ \hline
\multicolumn{4}{|L{\dimexpr 0.84\linewidth+6\tabcolsep+3\arrayrulewidth\relax}|}{\textbf{Comments:} -} \\ \hline
\multicolumn{4}{|C{\dimexpr 0.84\linewidth+6\tabcolsep+3\arrayrulewidth\relax}|}{\textbf{\fbox{Passed} \quad Failed \quad Not Executed}} \\ \hline
\end{longtable}

\begin{longtable}{|L{0.21\linewidth}|L{0.27\linewidth}|L{0.17\linewidth}|L{0.19\linewidth}|}
\caption{Test case TC-31: Tracklet building and visibility filtering}\label{tab:tc-31} \\ \hline
\multicolumn{4}{|C{\dimexpr 0.84\linewidth+6\tabcolsep+3\arrayrulewidth\relax}|}{\textbf{Tracklet building and visibility filtering}} \\ \hline
\multicolumn{4}{|C{\dimexpr 0.84\linewidth+6\tabcolsep+3\arrayrulewidth\relax}|}{\textbf{src/\allowbreak{}tracking/\allowbreak{}tracklets.py, src/\allowbreak{}core/\allowbreak{}types.py; tests/\allowbreak{}test\_\allowbreak{}stages.py (TestTracklets)}} \\ \hline
\textbf{Test Case ID:} & TC-31 & \textbf{QA Test Engineer:} & Abi \\ \hline
\textbf{Test case Version:} & 3 & \textbf{Reviewed By:} & - \\ \hline
\textbf{Test Date:} & 2026-11-18 & \textbf{Use Case Reference(s):} & Design \S{}4 (stage 4) \\ \hline
\textbf{Revision History:} & \multicolumn{3}{L{\dimexpr 0.63\linewidth+4\tabcolsep+2\arrayrulewidth\relax}|}{v1 2026-11-14 (178c9bc); v2 2026-11-15 (be0a468); v3 2026-11-16 (d3ab703)} \\ \hline
\textbf{Objective:} & \multicolumn{3}{L{\dimexpr 0.63\linewidth+4\tabcolsep+2\arrayrulewidth\relax}|}{Check that per-frame tracker states are assembled into tracklets with their per-frame signals, and that filter\_\allowbreak{}visible splits a tracklet only where the tracker lost it, not where the camera dropped frames.} \\ \hline
\textbf{Product/Ver/ Module:} & \multicolumn{3}{L{\dimexpr 0.63\linewidth+4\tabcolsep+2\arrayrulewidth\relax}|}{FYP-1 attention pipeline, commit e743616. Feature pass, first step: turns the frame-pass tracker output into tracklets and removes short or broken pieces before windowing.} \\ \hline
\textbf{Environment:} & \multicolumn{3}{L{\dimexpr 0.63\linewidth+4\tabcolsep+2\arrayrulewidth\relax}|}{Windows 11 (10.0.26200); conda environment fyp-local: Python 3.13.18, NumPy 2.5.4, SciPy 1.18.1, scikit-learn 1.9.1, pandas 3.1.1, OpenCV 4.11.0, PyTorch 2.9.0+cu126. CPU only; no GPU, dataset or model weights needed unless stated. Run from the repository root: python -m unittest discover -s tests -t .} \\ \hline
\textbf{Assumptions:} & \multicolumn{3}{L{\dimexpr 0.63\linewidth+4\tabcolsep+2\arrayrulewidth\relax}|}{A synthetic tracker output stands in for real tracks: one track (TID 7) with a fixed box [0, 0, 10, 10], confidence 0.9 and a per-frame signal 'sig' equal to the frame time. Gaps are counted in existing frames that lack the track, converted to seconds with the given fps.} \\ \hline
\textbf{Pre-Requisite:} & \multicolumn{3}{L{\dimexpr 0.63\linewidth+4\tabcolsep+2\arrayrulewidth\relax}|}{None beyond the Python environment.} \\ \hline
\textbf{Step No.} & \textbf{Execution description} & \multicolumn{2}{L{\dimexpr 0.36\linewidth+2\tabcolsep+1\arrayrulewidth\relax}|}{\textbf{Procedure result}} \\ \hline
1 & \texttt{test\_\allowbreak{}split\_\allowbreak{}where\_\allowbreak{}tracker\_\allowbreak{}lost\_\allowbreak{}the\_\allowbreak{}track}: Frames exist every 0.2 s from 0 to 3.8 s (5 fps). The track is present at t = 0, 0.2, 0.4, 0.6, 2.4, 2.6, so it is absent from 8 existing frames (1.6 s). Build the tracklet, then call filter\_\allowbreak{}visible with fps 5, max\_\allowbreak{}gap\_\allowbreak{}s 1.0 and min\_\allowbreak{}visible\_\allowbreak{}s 0.0, then again with min\_\allowbreak{}visible\_\allowbreak{}s 0.7. & \multicolumn{2}{L{\dimexpr 0.36\linewidth+2\tabcolsep+1\arrayrulewidth\relax}|}{The built tracklet's 'sig' signal equals the input times [0, 0.2, 0.4, 0.6, 2.4, 2.6]. With min\_\allowbreak{}visible\_\allowbreak{}s 0.0 the tracklet is split into two pieces of 4 and 2 observations. With min\_\allowbreak{}visible\_\allowbreak{}s 0.7 only the 4-observation piece (0.8 s visible) is kept; the 2-observation piece (0.4 s) is dropped. (Passed)} \\ \hline
2 & \texttt{test\_\allowbreak{}camera\_\allowbreak{}frame\_\allowbreak{}drop\_\allowbreak{}is\_\allowbreak{}not\_\allowbreak{}a\_\allowbreak{}gap}: Frame times are [0, 0.2, 0.4, 0.6, 3.6, 3.8]: the camera dropped 3 s of frames, but the track is present in every frame that exists. Call filter\_\allowbreak{}visible with fps 5, min\_\allowbreak{}visible\_\allowbreak{}s 0.0, max\_\allowbreak{}gap\_\allowbreak{}s 1.0. & \multicolumn{2}{L{\dimexpr 0.36\linewidth+2\tabcolsep+1\arrayrulewidth\relax}|}{One piece with all 6 observations is returned; the wall-clock jump of 3 s is not treated as a tracking gap. (Passed)} \\ \hline
\multicolumn{4}{|L{\dimexpr 0.84\linewidth+6\tabcolsep+3\arrayrulewidth\relax}|}{\textbf{Comments:} The second step documents a property of the DIPSER footage: cameras drop frames for several seconds, and a track that keeps its TID across such a drop was not lost by the tracker.} \\ \hline
\multicolumn{4}{|C{\dimexpr 0.84\linewidth+6\tabcolsep+3\arrayrulewidth\relax}|}{\textbf{\fbox{Passed} \quad Failed \quad Not Executed}} \\ \hline
\end{longtable}

\begin{longtable}{|L{0.21\linewidth}|L{0.27\linewidth}|L{0.17\linewidth}|L{0.19\linewidth}|}
\caption{Test case TC-32: SID-TID association}\label{tab:tc-32} \\ \hline
\multicolumn{4}{|C{\dimexpr 0.84\linewidth+6\tabcolsep+3\arrayrulewidth\relax}|}{\textbf{SID-TID association}} \\ \hline
\multicolumn{4}{|C{\dimexpr 0.84\linewidth+6\tabcolsep+3\arrayrulewidth\relax}|}{\textbf{src/\allowbreak{}tracking/\allowbreak{}association.py; tests/\allowbreak{}test\_\allowbreak{}stages.py (TestAssociation)}} \\ \hline
\textbf{Test Case ID:} & TC-32 & \textbf{QA Test Engineer:} & Abi \\ \hline
\textbf{Test case Version:} & 3 & \textbf{Reviewed By:} & - \\ \hline
\textbf{Test Date:} & 2026-11-18 & \textbf{Use Case Reference(s):} & Design \S{}4 \\ \hline
\textbf{Revision History:} & \multicolumn{3}{L{\dimexpr 0.63\linewidth+4\tabcolsep+2\arrayrulewidth\relax}|}{v1 2026-11-14 (178c9bc); v2 2026-11-15 (be0a468); v3 2026-11-16 (d3ab703)} \\ \hline
\textbf{Objective:} & \multicolumn{3}{L{\dimexpr 0.63\linewidth+4\tabcolsep+2\arrayrulewidth\relax}|}{Check that tracks are linked to annotated student boxes by per-frame Hungarian IoU matching and a majority vote, and that a person who only passes over a student's box is not linked to that student.} \\ \hline
\textbf{Product/Ver/ Module:} & \multicolumn{3}{L{\dimexpr 0.63\linewidth+4\tabcolsep+2\arrayrulewidth\relax}|}{FYP-1 attention pipeline, commit e743616. Training and evaluation only: links tracker IDs (TIDs) to the manually annotated subject IDs (SIDs) so per-student labels can be attached to tracks. Not used at inference.} \\ \hline
\textbf{Environment:} & \multicolumn{3}{L{\dimexpr 0.63\linewidth+4\tabcolsep+2\arrayrulewidth\relax}|}{Windows 11 (10.0.26200); conda environment fyp-local: Python 3.13.18, NumPy 2.5.4, SciPy 1.18.1, scikit-learn 1.9.1, pandas 3.1.1, OpenCV 4.11.0, PyTorch 2.9.0+cu126. CPU only; no GPU, dataset or model weights needed unless stated. Run from the repository root: python -m unittest discover -s tests -t .} \\ \hline
\textbf{Assumptions:} & \multicolumn{3}{L{\dimexpr 0.63\linewidth+4\tabcolsep+2\arrayrulewidth\relax}|}{Subject annotations are one static box per subject, as in DIPSER. Synthetic tracklets of 20 frames (0.4 s apart) with 50 x 100 px boxes stand in for tracker output. min\_\allowbreak{}iou = 0.3.} \\ \hline
\textbf{Pre-Requisite:} & \multicolumn{3}{L{\dimexpr 0.63\linewidth+4\tabcolsep+2\arrayrulewidth\relax}|}{None beyond the Python environment (numpy, scipy).} \\ \hline
\textbf{Step No.} & \textbf{Execution description} & \multicolumn{2}{L{\dimexpr 0.36\linewidth+2\tabcolsep+1\arrayrulewidth\relax}|}{\textbf{Procedure result}} \\ \hline
1 & \texttt{test\_\allowbreak{}majority\_\allowbreak{}vote}: Two subject boxes: subject\_\allowbreak{}01 at [100, 100, 150, 200] and subject\_\allowbreak{}02 at [400, 100, 450, 200]. Three static tracks: TID 1 on subject\_\allowbreak{}01's box, TID 2 on subject\_\allowbreak{}02's box, TID 3 at x = 900 (overlapping no subject). Call associate\_\allowbreak{}sids with min\_\allowbreak{}iou 0.3 and min\_\allowbreak{}match\_\allowbreak{}share 0.5. & \multicolumn{2}{L{\dimexpr 0.36\linewidth+2\tabcolsep+1\arrayrulewidth\relax}|}{The result is exactly \{1: 'subject\_\allowbreak{}01', 2: 'subject\_\allowbreak{}02'\}; TID 3 is left unassociated. (Passed)} \\ \hline
2 & \texttt{test\_\allowbreak{}walker\_\allowbreak{}passing\_\allowbreak{}a\_\allowbreak{}student\_\allowbreak{}stays\_\allowbreak{}unassociated}: One subject box (subject\_\allowbreak{}01). TID 1 sits still on that box. TID 9 has 40 frames and moves right at 50 px/\allowbreak{}s from x = 0, so it sweeps across the student's box for a few frames. Call associate\_\allowbreak{}sids on both tracks with min\_\allowbreak{}match\_\allowbreak{}share 0.5, then on the walker alone with min\_\allowbreak{}match\_\allowbreak{}share 0.0. & \multicolumn{2}{L{\dimexpr 0.36\linewidth+2\tabcolsep+1\arrayrulewidth\relax}|}{With both tracks and min\_\allowbreak{}match\_\allowbreak{}share 0.5 the result is exactly \{1: 'subject\_\allowbreak{}01'\}; the walker (TID 9) gets no SID. With the walker alone and min\_\allowbreak{}match\_\allowbreak{}share 0.0 (the earlier behaviour without the share guard), TID 9 is associated. (Passed)} \\ \hline
\multicolumn{4}{|L{\dimexpr 0.84\linewidth+6\tabcolsep+3\arrayrulewidth\relax}|}{\textbf{Comments:} The second call in test\_\allowbreak{}walker\_\allowbreak{}passing\_\allowbreak{}a\_\allowbreak{}student\_\allowbreak{}stays\_\allowbreak{}unassociated reproduces the earlier behaviour on purpose, to show the case the min\_\allowbreak{}match\_\allowbreak{}share guard was added for (a teaching assistant walking past a seated student).} \\ \hline
\multicolumn{4}{|C{\dimexpr 0.84\linewidth+6\tabcolsep+3\arrayrulewidth\relax}|}{\textbf{\fbox{Passed} \quad Failed \quad Not Executed}} \\ \hline
\end{longtable}

\begin{longtable}{|L{0.21\linewidth}|L{0.27\linewidth}|L{0.17\linewidth}|L{0.19\linewidth}|}
\caption{Test case TC-33: BFMU (a): Kalman/\allowbreak{}geometry motion features}\label{tab:tc-33} \\ \hline
\multicolumn{4}{|C{\dimexpr 0.84\linewidth+6\tabcolsep+3\arrayrulewidth\relax}|}{\textbf{BFMU (a): Kalman/\allowbreak{}geometry motion features}} \\ \hline
\multicolumn{4}{|C{\dimexpr 0.84\linewidth+6\tabcolsep+3\arrayrulewidth\relax}|}{\textbf{src/\allowbreak{}bfmu/\allowbreak{}kalman\_\allowbreak{}geometry.py, src/\allowbreak{}bfmu/\allowbreak{}base.py; tests/\allowbreak{}test\_\allowbreak{}stages.py (TestKalmanGeometry)}} \\ \hline
\textbf{Test Case ID:} & TC-33 & \textbf{QA Test Engineer:} & Abi \\ \hline
\textbf{Test case Version:} & 3 & \textbf{Reviewed By:} & - \\ \hline
\textbf{Test Date:} & 2026-11-18 & \textbf{Use Case Reference(s):} & Design \S{}5a \\ \hline
\textbf{Revision History:} & \multicolumn{3}{L{\dimexpr 0.63\linewidth+4\tabcolsep+2\arrayrulewidth\relax}|}{v1 2026-11-14 (178c9bc); v2 2026-11-15 (be0a468); v3 2026-11-16 (d3ab703)} \\ \hline
\textbf{Objective:} & \multicolumn{3}{L{\dimexpr 0.63\linewidth+4\tabcolsep+2\arrayrulewidth\relax}|}{Check the per-track window features computed from box motion: speed normalised by the box diagonal, stillness, and aspect trend, and the choice between the tracker's Kalman velocity and finite differences.} \\ \hline
\textbf{Product/Ver/ Module:} & \multicolumn{3}{L{\dimexpr 0.63\linewidth+4\tabcolsep+2\arrayrulewidth\relax}|}{FYP-1 attention pipeline, commit e743616. Feature pass: per-track behaviour features computed from tracker output only, later aggregated per window.} \\ \hline
\textbf{Environment:} & \multicolumn{3}{L{\dimexpr 0.63\linewidth+4\tabcolsep+2\arrayrulewidth\relax}|}{Windows 11 (10.0.26200); conda environment fyp-local: Python 3.13.18, NumPy 2.5.4, SciPy 1.18.1, scikit-learn 1.9.1, pandas 3.1.1, OpenCV 4.11.0, PyTorch 2.9.0+cu126. CPU only; no GPU, dataset or model weights needed unless stated. Run from the repository root: python -m unittest discover -s tests -t .} \\ \hline
\textbf{Assumptions:} & \multicolumn{3}{L{\dimexpr 0.63\linewidth+4\tabcolsep+2\arrayrulewidth\relax}|}{A synthetic tracklet of 20 frames, 0.4 s apart, with a 50 x 100 px box moving right at 10 px/\allowbreak{}s stands in for a real track. Default still\_\allowbreak{}threshold is 0.05 diagonals/\allowbreak{}s. The extractor is built through the FEATURE\_\allowbreak{}EXTRACTORS registry.} \\ \hline
\textbf{Pre-Requisite:} & \multicolumn{3}{L{\dimexpr 0.63\linewidth+4\tabcolsep+2\arrayrulewidth\relax}|}{None beyond the Python environment.} \\ \hline
\textbf{Step No.} & \textbf{Execution description} & \multicolumn{2}{L{\dimexpr 0.36\linewidth+2\tabcolsep+1\arrayrulewidth\relax}|}{\textbf{Procedure result}} \\ \hline
1 & \texttt{test\_\allowbreak{}finite\_\allowbreak{}difference\_\allowbreak{}speed}: Build 'kalman\_\allowbreak{}geometry' with use\_\allowbreak{}kalman=True. The tracklet's Kalman velocity is NaN, so the extractor must fall back to finite differences. Call window\_\allowbreak{}features(tracklet, fps=2.5). & \multicolumn{2}{L{\dimexpr 0.36\linewidth+2\tabcolsep+1\arrayrulewidth\relax}|}{kg.speed equals 10 /\allowbreak{} hypot(50, 100) (about 0.0894 diagonals/\allowbreak{}s) to 6 decimal places; kg.stillness is 0.0 (the track is always above the threshold); kg.aspect\_\allowbreak{}trend is 0.0 (the box shape does not change). (Passed)} \\ \hline
2 & \texttt{test\_\allowbreak{}kalman\_\allowbreak{}velocity\_\allowbreak{}preferred\_\allowbreak{}when\_\allowbreak{}present}: Build 'kalman\_\allowbreak{}geometry' with defaults. Take the same moving tracklet but set its Kalman velocity to [0, 0] on every frame (the tracker reports the person as still). Call window\_\allowbreak{}features with fps 2.5. & \multicolumn{2}{L{\dimexpr 0.36\linewidth+2\tabcolsep+1\arrayrulewidth\relax}|}{kg.stillness is 1.0: the finite Kalman velocity is used instead of the box displacement. (Passed)} \\ \hline
\multicolumn{4}{|L{\dimexpr 0.84\linewidth+6\tabcolsep+3\arrayrulewidth\relax}|}{\textbf{Comments:} -} \\ \hline
\multicolumn{4}{|C{\dimexpr 0.84\linewidth+6\tabcolsep+3\arrayrulewidth\relax}|}{\textbf{\fbox{Passed} \quad Failed \quad Not Executed}} \\ \hline
\end{longtable}

\begin{longtable}{|L{0.21\linewidth}|L{0.27\linewidth}|L{0.17\linewidth}|L{0.19\linewidth}|}
\caption{Test case TC-34: Window aggregation (class-level vector Vec)}\label{tab:tc-34} \\ \hline
\multicolumn{4}{|C{\dimexpr 0.84\linewidth+6\tabcolsep+3\arrayrulewidth\relax}|}{\textbf{Window aggregation (class-level vector Vec)}} \\ \hline
\multicolumn{4}{|C{\dimexpr 0.84\linewidth+6\tabcolsep+3\arrayrulewidth\relax}|}{\textbf{src/\allowbreak{}aggregation/\allowbreak{}vec.py, src/\allowbreak{}cfmu/\allowbreak{}extractor.py; tests/\allowbreak{}test\_\allowbreak{}stages.py (TestAggregation)}} \\ \hline
\textbf{Test Case ID:} & TC-34 & \textbf{QA Test Engineer:} & Abi \\ \hline
\textbf{Test case Version:} & 3 & \textbf{Reviewed By:} & - \\ \hline
\textbf{Test Date:} & 2026-11-18 & \textbf{Use Case Reference(s):} & Design \S{}6 \\ \hline
\textbf{Revision History:} & \multicolumn{3}{L{\dimexpr 0.63\linewidth+4\tabcolsep+2\arrayrulewidth\relax}|}{v1 2026-11-14 (178c9bc); v2 2026-11-15 (be0a468); v3 2026-11-16 (d3ab703)} \\ \hline
\textbf{Objective:} & \multicolumn{3}{L{\dimexpr 0.63\linewidth+4\tabcolsep+2\arrayrulewidth\relax}|}{Check that every window produces the same set of keys, that windows without visible tracks give NaN statistics, and that track count, threshold fractions and the occupancy heatmap are computed correctly.} \\ \hline
\textbf{Product/Ver/ Module:} & \multicolumn{3}{L{\dimexpr 0.63\linewidth+4\tabcolsep+2\arrayrulewidth\relax}|}{FYP-1 attention pipeline, commit e743616. Feature pass: combines the per-track features of all visible tracks in a window into one fixed-length class-level vector (Vec).} \\ \hline
\textbf{Environment:} & \multicolumn{3}{L{\dimexpr 0.63\linewidth+4\tabcolsep+2\arrayrulewidth\relax}|}{Windows 11 (10.0.26200); conda environment fyp-local: Python 3.13.18, NumPy 2.5.4, SciPy 1.18.1, scikit-learn 1.9.1, pandas 3.1.1, OpenCV 4.11.0, PyTorch 2.9.0+cu126. CPU only; no GPU, dataset or model weights needed unless stated. Run from the repository root: python -m unittest discover -s tests -t .} \\ \hline
\textbf{Assumptions:} & \multicolumn{3}{L{\dimexpr 0.63\linewidth+4\tabcolsep+2\arrayrulewidth\relax}|}{Two synthetic tracklets (20 frames, 0.4 s apart, so t = 0 to 7.6 s): TID 1 moving at 10 px/\allowbreak{}s, TID 2 still. One BFMU extractor (kalman\_\allowbreak{}geometry), the null CFMU extractor, a threshold of 0.5 on kg.stillness, frame size 1280 x 720, default 3 x 3 heatmap and min\_\allowbreak{}coverage 0.5.} \\ \hline
\textbf{Pre-Requisite:} & \multicolumn{3}{L{\dimexpr 0.63\linewidth+4\tabcolsep+2\arrayrulewidth\relax}|}{None beyond the Python environment.} \\ \hline
\textbf{Step No.} & \textbf{Execution description} & \multicolumn{2}{L{\dimexpr 0.36\linewidth+2\tabcolsep+1\arrayrulewidth\relax}|}{\textbf{Procedure result}} \\ \hline
1 & \texttt{test\_\allowbreak{}vec\_\allowbreak{}keys\_\allowbreak{}stable\_\allowbreak{}and\_\allowbreak{}nan\_\allowbreak{}when\_\allowbreak{}empty}: Aggregate window [0, 8) with both tracklets (fps 2.5), and window [100, 108) with TID 1, which has no observations there. & \multicolumn{2}{L{\dimexpr 0.36\linewidth+2\tabcolsep+1\arrayrulewidth\relax}|}{Both vectors have the same sorted key set. The full window has n\_\allowbreak{}tracks = 2.0 and kg.stillness.frac = 0.5 (one still track of two). The empty window has kg.speed.mean = NaN. The occupancy keys (occ\_\allowbreak{}*) of the full window sum to 1.0. (Passed)} \\ \hline
\multicolumn{4}{|L{\dimexpr 0.84\linewidth+6\tabcolsep+3\arrayrulewidth\relax}|}{\textbf{Comments:} -} \\ \hline
\multicolumn{4}{|C{\dimexpr 0.84\linewidth+6\tabcolsep+3\arrayrulewidth\relax}|}{\textbf{\fbox{Passed} \quad Failed \quad Not Executed}} \\ \hline
\end{longtable}

\begin{longtable}{|L{0.21\linewidth}|L{0.27\linewidth}|L{0.17\linewidth}|L{0.19\linewidth}|}
\caption{Test case TC-35: Configuration hashing, cache keys, pipeline wiring, track store and output graph}\label{tab:tc-35} \\ \hline
\multicolumn{4}{|C{\dimexpr 0.84\linewidth+6\tabcolsep+3\arrayrulewidth\relax}|}{\textbf{Configuration hashing, cache keys, pipeline wiring, track store and output graph}} \\ \hline
\multicolumn{4}{|C{\dimexpr 0.84\linewidth+6\tabcolsep+3\arrayrulewidth\relax}|}{\textbf{src/\allowbreak{}core/\allowbreak{}config.py, src/\allowbreak{}core/\allowbreak{}pipeline.py, src/\allowbreak{}core/\allowbreak{}store.py, src/\allowbreak{}output/\allowbreak{}graph.py; tests/\allowbreak{}test\_\allowbreak{}stages.py (TestConfigAndWiring)}} \\ \hline
\textbf{Test Case ID:} & TC-35 & \textbf{QA Test Engineer:} & Abi \\ \hline
\textbf{Test case Version:} & 3 & \textbf{Reviewed By:} & - \\ \hline
\textbf{Test Date:} & 2026-11-18 & \textbf{Use Case Reference(s):} & Design \S{}3, \S{}9 \\ \hline
\textbf{Revision History:} & \multicolumn{3}{L{\dimexpr 0.63\linewidth+4\tabcolsep+2\arrayrulewidth\relax}|}{v1 2026-11-14 (178c9bc); v2 2026-11-15 (be0a468); v3 2026-11-16 (d3ab703)} \\ \hline
\textbf{Objective:} & \multicolumn{3}{L{\dimexpr 0.63\linewidth+4\tabcolsep+2\arrayrulewidth\relax}|}{Check that config hashes and the two cache keys (frame pass and feature pass) change exactly when the settings they depend on change, that the pipeline can be built and run from a cached track without a GPU or video, that the track cache round-trips, and that the output graph carries no identity fields.} \\ \hline
\textbf{Product/Ver/ Module:} & \multicolumn{3}{L{\dimexpr 0.63\linewidth+4\tabcolsep+2\arrayrulewidth\relax}|}{FYP-1 attention pipeline, commit e743616. Core wiring: config loading and hashing select cached results; the Pipeline builds every stage from the config; the track store persists the frame pass; build\_\allowbreak{}graph writes the final retention graph.} \\ \hline
\textbf{Environment:} & \multicolumn{3}{L{\dimexpr 0.63\linewidth+4\tabcolsep+2\arrayrulewidth\relax}|}{Windows 11 (10.0.26200); conda environment fyp-local: Python 3.13.18, NumPy 2.5.4, SciPy 1.18.1, scikit-learn 1.9.1, pandas 3.1.1, OpenCV 4.11.0, PyTorch 2.9.0+cu126. CPU only; no GPU, dataset or model weights needed unless stated. Run from the repository root: python -m unittest discover -s tests -t .} \\ \hline
\textbf{Assumptions:} & \multicolumn{3}{L{\dimexpr 0.63\linewidth+4\tabcolsep+2\arrayrulewidth\relax}|}{configs/\allowbreak{}deliverable3.json is the reference config. It names configs/\allowbreak{}derived.py, which reads values from experiment results files. Synthetic tracklets and scene rows stand in for a real frame pass. The SCB behaviour detector loads its weights on first use, so building the pipeline needs no weights.} \\ \hline
\textbf{Pre-Requisite:} & \multicolumn{3}{L{\dimexpr 0.63\linewidth+4\tabcolsep+2\arrayrulewidth\relax}|}{Run from the repository root. configs/\allowbreak{}deliverable3.json and configs/\allowbreak{}derived.py must be present, and the results files read by configs/\allowbreak{}derived.py (outputs/\allowbreak{}tuning/\allowbreak{}detector\_\allowbreak{}search\_\allowbreak{}augmented.json and outputs/\allowbreak{}tuning/\allowbreak{}tracker\_\allowbreak{}botsort\_\allowbreak{}full\_\allowbreak{}reid\_\allowbreak{}cmc\_\allowbreak{}aug1-4f270e71\_\allowbreak{}bisect.json) must exist. No GPU, video or model weights are needed. test\_\allowbreak{}track\_\allowbreak{}store\_\allowbreak{}round\_\allowbreak{}trip writes only to a temporary directory.} \\ \hline
\textbf{Step No.} & \textbf{Execution description} & \multicolumn{2}{L{\dimexpr 0.36\linewidth+2\tabcolsep+1\arrayrulewidth\relax}|}{\textbf{Procedure result}} \\ \hline
1 & \texttt{test\_\allowbreak{}hash\_\allowbreak{}excludes\_\allowbreak{}output\_\allowbreak{}dir\_\allowbreak{}and\_\allowbreak{}tuning}: Load the config three times: unchanged; with output\_\allowbreak{}dir = 'elsewhere' and tuning.seed = 7; with labels.tau\_\allowbreak{}att = 3. Compute config\_\allowbreak{}hash for each. & \multicolumn{2}{L{\dimexpr 0.36\linewidth+2\tabcolsep+1\arrayrulewidth\relax}|}{The first two hashes are equal (output\_\allowbreak{}dir and tuning are excluded). The hash with labels.tau\_\allowbreak{}att = 3 differs from the base hash. (Passed)} \\ \hline
2 & \texttt{test\_\allowbreak{}cache\_\allowbreak{}keys\_\allowbreak{}split\_\allowbreak{}frame\_\allowbreak{}pass\_\allowbreak{}from\_\allowbreak{}feature\_\allowbreak{}pass}: Build config variants: labels.tau\_\allowbreak{}att = 3; window.length\_\allowbreak{}s = 10; kalman\_\allowbreak{}geometry still\_\allowbreak{}threshold = 0.1 (bfmu[0], a window-level setting); residual\_\allowbreak{}flow grid + 1 (bfmu[1], a pixel-level setting); tracker.name = 'strongsort'. Compute tracking\_\allowbreak{}hash and feature\_\allowbreak{}hash for each and compare with the base config. & \multicolumn{2}{L{\dimexpr 0.36\linewidth+2\tabcolsep+1\arrayrulewidth\relax}|}{The label change leaves feature\_\allowbreak{}hash unchanged. The window length and still\_\allowbreak{}threshold changes leave tracking\_\allowbreak{}hash unchanged but change feature\_\allowbreak{}hash (feature pass only). The residual\_\allowbreak{}flow grid and tracker changes change both tracking\_\allowbreak{}hash and feature\_\allowbreak{}hash (a new frame pass is needed). (Passed)} \\ \hline
3 & \texttt{test\_\allowbreak{}pipeline\_\allowbreak{}builds\_\allowbreak{}without\_\allowbreak{}gpu}: Construct Pipeline(cfg) from the reference config. & \multicolumn{2}{L{\dimexpr 0.36\linewidth+2\tabcolsep+1\arrayrulewidth\relax}|}{Construction succeeds and the extractor classes, in order, are ['KalmanGeometry', 'ResidualFlow', 'ScbBehaviour']. (Passed)} \\ \hline
4 & \texttt{test\_\allowbreak{}track\_\allowbreak{}store\_\allowbreak{}round\_\allowbreak{}trip}: Build a SessionTrack (group\_\allowbreak{}01 /\allowbreak{} experiment\_\allowbreak{}02 /\allowbreak{} context\_\allowbreak{}cam\_\allowbreak{}03, 30 frame times 0.4 s apart, frame size 1280 x 720) with two tracklets: TID 1 with velocity [1.5, -2.0] and a signal 'rf.mag\_\allowbreak{}00' (linspace 0 to 1), and TID 4 (5 frames) with no signals. Save it with save\_\allowbreak{}track to a temporary .npz and load it with load\_\allowbreak{}track. & \multicolumn{2}{L{\dimexpr 0.36\linewidth+2\tabcolsep+1\arrayrulewidth\relax}|}{Key, frame size, timing and scene rows are equal to the original. Tracklet TIDs are [1, 4]. For each tracklet the t, xyxy, conf and velocity arrays are equal to the original. TID 1's 'rf.mag\_\allowbreak{}00' signal is equal to the original; TID 4's 'rf.mag\_\allowbreak{}00' is all NaN. (Passed)} \\ \hline
5 & \texttt{test\_\allowbreak{}feature\_\allowbreak{}pass\_\allowbreak{}runs\_\allowbreak{}from\_\allowbreak{}a\_\allowbreak{}track\_\allowbreak{}without\_\allowbreak{}video}: Build a SessionTrack with 100 frame times 0.2 s apart (0 to 19.8 s, fps 5), three still tracklets of 100 frames at x = 100, 300 and 500, scene rows \{'motion': 0.0\}, and no session key. Override window to length 8 s with 50\% overlap and call Pipeline(cfg).features(track). & \multicolumn{2}{L{\dimexpr 0.36\linewidth+2\tabcolsep+1\arrayrulewidth\relax}|}{Three windows are produced. Each window's vector has n\_\allowbreak{}tracks = 3.0. The first window's scene features contain 'motion.mean'. (Passed)} \\ \hline
6 & \texttt{test\_\allowbreak{}graph\_\allowbreak{}has\_\allowbreak{}no\_\allowbreak{}identity\_\allowbreak{}fields}: Call build\_\allowbreak{}graph with one window [0, 8), context probabilities [[0.2, 0.8]], classes ['01', '02'], engagement probability [0.7], track count [12] and tau 0.5. & \multicolumn{2}{L{\dimexpr 0.36\linewidth+2\tabcolsep+1\arrayrulewidth\relax}|}{The record's keys are exactly \{t\_\allowbreak{}start, t\_\allowbreak{}end, predicted\_\allowbreak{}context, layout\_\allowbreak{}family, engaged, engagement\_\allowbreak{}prob, n\_\allowbreak{}visible\_\allowbreak{}tracks\}: no box, TID or SID. predicted\_\allowbreak{}context.label is '02'. (Passed)} \\ \hline
\multicolumn{4}{|L{\dimexpr 0.84\linewidth+6\tabcolsep+3\arrayrulewidth\relax}|}{\textbf{Comments:} test\_\allowbreak{}feature\_\allowbreak{}pass\_\allowbreak{}runs\_\allowbreak{}from\_\allowbreak{}a\_\allowbreak{}track\_\allowbreak{}without\_\allowbreak{}video pins the window settings (8 s, 50\%) so the expected window count does not depend on the value in the config file (4 s at the time of writing).} \\ \hline
\multicolumn{4}{|C{\dimexpr 0.84\linewidth+6\tabcolsep+3\arrayrulewidth\relax}|}{\textbf{\fbox{Passed} \quad Failed \quad Not Executed}} \\ \hline
\end{longtable}

\begin{longtable}{|L{0.21\linewidth}|L{0.27\linewidth}|L{0.17\linewidth}|L{0.19\linewidth}|}
\caption{Test case TC-36: Inter-labeler agreement (Krippendorff's alpha, ordinal)}\label{tab:tc-36} \\ \hline
\multicolumn{4}{|C{\dimexpr 0.84\linewidth+6\tabcolsep+3\arrayrulewidth\relax}|}{\textbf{Inter-labeler agreement (Krippendorff's alpha, ordinal)}} \\ \hline
\multicolumn{4}{|C{\dimexpr 0.84\linewidth+6\tabcolsep+3\arrayrulewidth\relax}|}{\textbf{src/\allowbreak{}evaluation/\allowbreak{}agreement.py; tests/\allowbreak{}test\_\allowbreak{}tuning.py (TestAlpha)}} \\ \hline
\textbf{Test Case ID:} & TC-36 & \textbf{QA Test Engineer:} & Abi \\ \hline
\textbf{Test case Version:} & 5 & \textbf{Reviewed By:} & - \\ \hline
\textbf{Test Date:} & 2026-11-18 & \textbf{Use Case Reference(s):} & Design \S{}11 \\ \hline
\textbf{Revision History:} & \multicolumn{3}{L{\dimexpr 0.63\linewidth+4\tabcolsep+2\arrayrulewidth\relax}|}{v1 2026-11-14 (178c9bc); v2 2026-11-15 (a497de6); v3 2026-11-15 (be0a468); v4 2026-11-16 (d3ab703); v5 2026-11-18 (e664d52)} \\ \hline
\textbf{Objective:} & \multicolumn{3}{L{\dimexpr 0.63\linewidth+4\tabcolsep+2\arrayrulewidth\relax}|}{Check the ordinal Krippendorff's alpha implementation against a published reference value and against the perfect-agreement case.} \\ \hline
\textbf{Product/Ver/ Module:} & \multicolumn{3}{L{\dimexpr 0.63\linewidth+4\tabcolsep+2\arrayrulewidth\relax}|}{FYP-1 attention pipeline, commit e743616. Evaluation: measures agreement between human labelers, which sets the accuracy ceiling the attention models are compared against.} \\ \hline
\textbf{Environment:} & \multicolumn{3}{L{\dimexpr 0.63\linewidth+4\tabcolsep+2\arrayrulewidth\relax}|}{Windows 11 (10.0.26200); conda environment fyp-local: Python 3.13.18, NumPy 2.5.4, SciPy 1.18.1, scikit-learn 1.9.1, pandas 3.1.1, OpenCV 4.11.0, PyTorch 2.9.0+cu126. CPU only; no GPU, dataset or model weights needed unless stated. Run from the repository root: python -m unittest discover -s tests -t .} \\ \hline
\textbf{Assumptions:} & \multicolumn{3}{L{\dimexpr 0.63\linewidth+4\tabcolsep+2\arrayrulewidth\relax}|}{Missing ratings are given as NaN. Units with fewer than two ratings are ignored. The reference value is the ordinal alpha of the worked example in Krippendorff (2011), 'Computing Krippendorff's Alpha-Reliability'.} \\ \hline
\textbf{Pre-Requisite:} & \multicolumn{3}{L{\dimexpr 0.63\linewidth+4\tabcolsep+2\arrayrulewidth\relax}|}{None beyond the Python environment (numpy).} \\ \hline
\textbf{Step No.} & \textbf{Execution description} & \multicolumn{2}{L{\dimexpr 0.36\linewidth+2\tabcolsep+1\arrayrulewidth\relax}|}{\textbf{Procedure result}} \\ \hline
1 & \texttt{test\_\allowbreak{}krippendorff\_\allowbreak{}reference\_\allowbreak{}example}: Build the reference data set of 4 coders by 12 units (values 1 to 5, with NaN for missing ratings), transpose it to one list of ratings per unit, and call krippendorff\_\allowbreak{}alpha\_\allowbreak{}ordinal. & \multicolumn{2}{L{\dimexpr 0.36\linewidth+2\tabcolsep+1\arrayrulewidth\relax}|}{Alpha equals 0.815 to 3 decimal places. (Passed)} \\ \hline
2 & \texttt{test\_\allowbreak{}perfect\_\allowbreak{}agreement}: Call krippendorff\_\allowbreak{}alpha\_\allowbreak{}ordinal on units [[1, 1], [2, 2, 2], [3, 3], [5, 5]], where every unit's ratings agree. & \multicolumn{2}{L{\dimexpr 0.36\linewidth+2\tabcolsep+1\arrayrulewidth\relax}|}{Alpha equals 1.0. (Passed)} \\ \hline
\multicolumn{4}{|L{\dimexpr 0.84\linewidth+6\tabcolsep+3\arrayrulewidth\relax}|}{\textbf{Comments:} -} \\ \hline
\multicolumn{4}{|C{\dimexpr 0.84\linewidth+6\tabcolsep+3\arrayrulewidth\relax}|}{\textbf{\fbox{Passed} \quad Failed \quad Not Executed}} \\ \hline
\end{longtable}

\begin{longtable}{|L{0.21\linewidth}|L{0.27\linewidth}|L{0.17\linewidth}|L{0.19\linewidth}|}
\caption{Test case TC-37: Tracking quality metrics}\label{tab:tc-37} \\ \hline
\multicolumn{4}{|C{\dimexpr 0.84\linewidth+6\tabcolsep+3\arrayrulewidth\relax}|}{\textbf{Tracking quality metrics}} \\ \hline
\multicolumn{4}{|C{\dimexpr 0.84\linewidth+6\tabcolsep+3\arrayrulewidth\relax}|}{\textbf{src/\allowbreak{}evaluation/\allowbreak{}tracking.py, src/\allowbreak{}tracking/\allowbreak{}association.py; tests/\allowbreak{}test\_\allowbreak{}tuning.py (TestTrackingMetrics)}} \\ \hline
\textbf{Test Case ID:} & TC-37 & \textbf{QA Test Engineer:} & Abi \\ \hline
\textbf{Test case Version:} & 5 & \textbf{Reviewed By:} & - \\ \hline
\textbf{Test Date:} & 2026-11-18 & \textbf{Use Case Reference(s):} & Design \S{}11 \\ \hline
\textbf{Revision History:} & \multicolumn{3}{L{\dimexpr 0.63\linewidth+4\tabcolsep+2\arrayrulewidth\relax}|}{v1 2026-11-14 (178c9bc); v2 2026-11-15 (a497de6); v3 2026-11-15 (be0a468); v4 2026-11-16 (d3ab703); v5 2026-11-18 (e664d52)} \\ \hline
\textbf{Objective:} & \multicolumn{3}{L{\dimexpr 0.63\linewidth+4\tabcolsep+2\arrayrulewidth\relax}|}{Check IDF1, coverage, ID switches, TIDs per subject and subjects found, computed against static subject boxes, for perfect tracking, a fragmented track and a missed subject.} \\ \hline
\textbf{Product/Ver/ Module:} & \multicolumn{3}{L{\dimexpr 0.63\linewidth+4\tabcolsep+2\arrayrulewidth\relax}|}{FYP-1 attention pipeline, commit e743616. Evaluation and tracker tuning: scores raw tracker output against the manual subject boxes; it is the objective of the tracker parameter search.} \\ \hline
\textbf{Environment:} & \multicolumn{3}{L{\dimexpr 0.63\linewidth+4\tabcolsep+2\arrayrulewidth\relax}|}{Windows 11 (10.0.26200); conda environment fyp-local: Python 3.13.18, NumPy 2.5.4, SciPy 1.18.1, scikit-learn 1.9.1, pandas 3.1.1, OpenCV 4.11.0, PyTorch 2.9.0+cu126. CPU only; no GPU, dataset or model weights needed unless stated. Run from the repository root: python -m unittest discover -s tests -t .} \\ \hline
\textbf{Assumptions:} & \multicolumn{3}{L{\dimexpr 0.63\linewidth+4\tabcolsep+2\arrayrulewidth\relax}|}{Two subjects with static boxes: a at [0, 0, 100, 100] and b at [200, 0, 300, 100]. Ground truth assumes every subject is present in every frame. Synthetic static tracks sit exactly on the subject boxes. 10 frames, min\_\allowbreak{}iou 0.3.} \\ \hline
\textbf{Pre-Requisite:} & \multicolumn{3}{L{\dimexpr 0.63\linewidth+4\tabcolsep+2\arrayrulewidth\relax}|}{None beyond the Python environment (numpy, scipy).} \\ \hline
\textbf{Step No.} & \textbf{Execution description} & \multicolumn{2}{L{\dimexpr 0.36\linewidth+2\tabcolsep+1\arrayrulewidth\relax}|}{\textbf{Procedure result}} \\ \hline
1 & \texttt{test\_\allowbreak{}perfect\_\allowbreak{}tracking}: TID 1 covers subject a and TID 2 covers subject b in all 10 frames. Call tracking\_\allowbreak{}metrics with n\_\allowbreak{}frames 10 and min\_\allowbreak{}iou 0.3. & \multicolumn{2}{L{\dimexpr 0.36\linewidth+2\tabcolsep+1\arrayrulewidth\relax}|}{idf1 is 1.0; coverage is 1.0; id\_\allowbreak{}switches\_\allowbreak{}per\_\allowbreak{}subject is 0.0; tids\_\allowbreak{}per\_\allowbreak{}subject is 1.0. (Passed)} \\ \hline
2 & \texttt{test\_\allowbreak{}id\_\allowbreak{}switch\_\allowbreak{}and\_\allowbreak{}fragmentation}: Subject a is covered by TID 1 in frames 0 to 5 and by TID 3 in frames 6 to 9; subject b is covered by TID 2 in all 10 frames. Call tracking\_\allowbreak{}metrics. & \multicolumn{2}{L{\dimexpr 0.36\linewidth+2\tabcolsep+1\arrayrulewidth\relax}|}{id\_\allowbreak{}switches\_\allowbreak{}per\_\allowbreak{}subject is 0.5 (one switch over two subjects); tids\_\allowbreak{}per\_\allowbreak{}subject is 1.5; idf1 equals 2*16 /\allowbreak{} (2*16 + 4 + 4): 16 correctly identified frames, and TID 3's 4 frames count both as false positives and as false negatives under the one-to-one TID-SID mapping. (Passed)} \\ \hline
3 & \texttt{test\_\allowbreak{}missed\_\allowbreak{}subject}: Only TID 1 is present, on subject a, for 10 frames. Call tracking\_\allowbreak{}metrics. & \multicolumn{2}{L{\dimexpr 0.36\linewidth+2\tabcolsep+1\arrayrulewidth\relax}|}{coverage is 0.5 and subjects\_\allowbreak{}found is 0.5. (Passed)} \\ \hline
\multicolumn{4}{|L{\dimexpr 0.84\linewidth+6\tabcolsep+3\arrayrulewidth\relax}|}{\textbf{Comments:} -} \\ \hline
\multicolumn{4}{|C{\dimexpr 0.84\linewidth+6\tabcolsep+3\arrayrulewidth\relax}|}{\textbf{\fbox{Passed} \quad Failed \quad Not Executed}} \\ \hline
\end{longtable}

\begin{longtable}{|L{0.21\linewidth}|L{0.27\linewidth}|L{0.17\linewidth}|L{0.19\linewidth}|}
\caption{Test case TC-38: Hyperparameter search (random, beam bisection, nested leave-one-group-out)}\label{tab:tc-38} \\ \hline
\multicolumn{4}{|C{\dimexpr 0.84\linewidth+6\tabcolsep+3\arrayrulewidth\relax}|}{\textbf{Hyperparameter search (random, beam bisection, nested leave-one-group-out)}} \\ \hline
\multicolumn{4}{|C{\dimexpr 0.84\linewidth+6\tabcolsep+3\arrayrulewidth\relax}|}{\textbf{src/\allowbreak{}tuning/\allowbreak{}search.py, src/\allowbreak{}spu/\allowbreak{}splits.py; tests/\allowbreak{}test\_\allowbreak{}tuning.py (TestSearch)}} \\ \hline
\textbf{Test Case ID:} & TC-38 & \textbf{QA Test Engineer:} & Abi \\ \hline
\textbf{Test case Version:} & 5 & \textbf{Reviewed By:} & - \\ \hline
\textbf{Test Date:} & 2026-11-18 & \textbf{Use Case Reference(s):} & Design \S{}7 (leave-one-group-out splits) \\ \hline
\textbf{Revision History:} & \multicolumn{3}{L{\dimexpr 0.63\linewidth+4\tabcolsep+2\arrayrulewidth\relax}|}{v1 2026-11-14 (178c9bc); v2 2026-11-15 (a497de6); v3 2026-11-15 (be0a468); v4 2026-11-16 (d3ab703); v5 2026-11-18 (e664d52)} \\ \hline
\textbf{Objective:} & \multicolumn{3}{L{\dimexpr 0.63\linewidth+4\tabcolsep+2\arrayrulewidth\relax}|}{Check that parameter sampling is seeded and in range, that the searches evaluate the baseline as trial 0 and return trials best first, that beam bisection finds a known optimum within its budget and is deterministic, and that nested search never selects parameters on its outer test group.} \\ \hline
\textbf{Product/Ver/ Module:} & \multicolumn{3}{L{\dimexpr 0.63\linewidth+4\tabcolsep+2\arrayrulewidth\relax}|}{FYP-1 attention pipeline, commit e743616. Tuning: searches over config-defined parameter spaces for the tracker, detector and models.} \\ \hline
\textbf{Environment:} & \multicolumn{3}{L{\dimexpr 0.63\linewidth+4\tabcolsep+2\arrayrulewidth\relax}|}{Windows 11 (10.0.26200); conda environment fyp-local: Python 3.13.18, NumPy 2.5.4, SciPy 1.18.1, scikit-learn 1.9.1, pandas 3.1.1, OpenCV 4.11.0, PyTorch 2.9.0+cu126. CPU only; no GPU, dataset or model weights needed unless stated. Run from the repository root: python -m unittest discover -s tests -t .} \\ \hline
\textbf{Assumptions:} & \multicolumn{3}{L{\dimexpr 0.63\linewidth+4\tabcolsep+2\arrayrulewidth\relax}|}{Objectives are cheap synthetic functions with known optima, not real pipeline runs. The search space format is: a list for a categorical parameter, or \{low, high, log, int\} for a range.} \\ \hline
\textbf{Pre-Requisite:} & \multicolumn{3}{L{\dimexpr 0.63\linewidth+4\tabcolsep+2\arrayrulewidth\relax}|}{None beyond the Python environment.} \\ \hline
\textbf{Step No.} & \textbf{Execution description} & \multicolumn{2}{L{\dimexpr 0.36\linewidth+2\tabcolsep+1\arrayrulewidth\relax}|}{\textbf{Procedure result}} \\ \hline
1 & \texttt{test\_\allowbreak{}sampling\_\allowbreak{}is\_\allowbreak{}seeded\_\allowbreak{}and\_\allowbreak{}in\_\allowbreak{}range}: Space: c categorical [1, 2, 3]; x log range 0.01 to 1.0; n integer range 1 to 9. Call sample\_\allowbreak{}params twice with numpy generators seeded 1. & \multicolumn{2}{L{\dimexpr 0.36\linewidth+2\tabcolsep+1\arrayrulewidth\relax}|}{Both samples are equal. c is in [1, 2, 3]; x is in [0.01, 1.0]; n is an int. (Passed)} \\ \hline
2 & \texttt{test\_\allowbreak{}baseline\_\allowbreak{}is\_\allowbreak{}trial\_\allowbreak{}zero\_\allowbreak{}and\_\allowbreak{}best\_\allowbreak{}first}: Run random\_\allowbreak{}search on the same space for 20 trials, seed 0, baseline \{'x': 0.9\}, with objective score = -|x - 0.3| (x defaults to 0.5 when absent). & \multicolumn{2}{L{\dimexpr 0.36\linewidth+2\tabcolsep+1\arrayrulewidth\relax}|}{The trial numbered 0 has params exactly \{'x': 0.9\}. The first trial in the returned list has the highest score of all trials. (Passed)} \\ \hline
3 & \texttt{test\_\allowbreak{}bisection\_\allowbreak{}finds\_\allowbreak{}the\_\allowbreak{}optimum\_\allowbreak{}and\_\allowbreak{}respects\_\allowbreak{}the\_\allowbreak{}budget}: Space: a in [0, 1]; b integer in [1, 100] on a log scale. Objective: -(a - 0.73)\textasciicircum{}2 - (ln b - ln 20)\textasciicircum{}2 /\allowbreak{} 10, maximal at a = 0.73, b = 20. Run beam\_\allowbreak{}bisection\_\allowbreak{}search with a budget of 60 trials, baseline \{a: 0.5, b: 50\} and beam 2. & \multicolumn{2}{L{\dimexpr 0.36\linewidth+2\tabcolsep+1\arrayrulewidth\relax}|}{At most 60 trials are run. Trial 0 has params \{a: 0.5, b: 50\}. The best trial has a within 0.03 of 0.73 and b between 15 and 27. (Passed)} \\ \hline
4 & \texttt{test\_\allowbreak{}bisection\_\allowbreak{}is\_\allowbreak{}deterministic\_\allowbreak{}and\_\allowbreak{}holds\_\allowbreak{}inactive\_\allowbreak{}and\_\allowbreak{}categorical\_\allowbreak{}params}: Space: a and z in [0, 1], c categorical [1, 2]. Objective -|a - 0.2|, recording every point it is called with. Run beam\_\allowbreak{}bisection\_\allowbreak{}search twice with a budget of 12, baseline \{a: 0.9, z: 0.4, c: 2\} and z marked inactive. & \multicolumn{2}{L{\dimexpr 0.36\linewidth+2\tabcolsep+1\arrayrulewidth\relax}|}{Both runs return the same list of params in the same order. Every evaluated point has z = 0.4 and c = 2: the inactive and categorical parameters keep their baseline values. (Passed)} \\ \hline
5 & \texttt{test\_\allowbreak{}nested\_\allowbreak{}search\_\allowbreak{}never\_\allowbreak{}scores\_\allowbreak{}selection\_\allowbreak{}on\_\allowbreak{}outer\_\allowbreak{}test}: 12 rows in three groups (g1, g2, g3, four rows each). The score function returns params['c'] and records the groups in each train and test index set. Run nested\_\allowbreak{}logo\_\allowbreak{}search with space \{c: [1, 2, 3]\}, 5 trials, seed 0, baseline \{c: 1\}. & \multicolumn{2}{L{\dimexpr 0.36\linewidth+2\tabcolsep+1\arrayrulewidth\relax}|}{The selected params are \{'c': 3\}. For every outer fold, at least one recorded score call used train and test sets that both exclude that fold's held-out group, showing the inner searches ran without it. (Passed)} \\ \hline
\multicolumn{4}{|L{\dimexpr 0.84\linewidth+6\tabcolsep+3\arrayrulewidth\relax}|}{\textbf{Comments:} The nested-search step checks that inner searches without the held-out group took place; it does not by itself prove that no inner call ever touched the held-out group.} \\ \hline
\multicolumn{4}{|C{\dimexpr 0.84\linewidth+6\tabcolsep+3\arrayrulewidth\relax}|}{\textbf{\fbox{Passed} \quad Failed \quad Not Executed}} \\ \hline
\end{longtable}

\begin{longtable}{|L{0.21\linewidth}|L{0.27\linewidth}|L{0.17\linewidth}|L{0.19\linewidth}|}
\caption{Test case TC-39: Shared-frame multi-tracker runs and search prefetch}\label{tab:tc-39} \\ \hline
\multicolumn{4}{|C{\dimexpr 0.84\linewidth+6\tabcolsep+3\arrayrulewidth\relax}|}{\textbf{Shared-frame multi-tracker runs and search prefetch}} \\ \hline
\multicolumn{4}{|C{\dimexpr 0.84\linewidth+6\tabcolsep+3\arrayrulewidth\relax}|}{\textbf{scripts/\allowbreak{}tuning/\allowbreak{}tune\_\allowbreak{}tracker.py, src/\allowbreak{}tuning/\allowbreak{}search.py, src/\allowbreak{}tracking/\allowbreak{}tracker.py, src/\allowbreak{}tracking/\allowbreak{}tracklets.py; tests/\allowbreak{}test\_\allowbreak{}tuning.py (TestSharedFrames)}} \\ \hline
\textbf{Test Case ID:} & TC-39 & \textbf{QA Test Engineer:} & Abi \\ \hline
\textbf{Test case Version:} & 5 & \textbf{Reviewed By:} & - \\ \hline
\textbf{Test Date:} & 2026-11-18 & \textbf{Use Case Reference(s):} & - \\ \hline
\textbf{Revision History:} & \multicolumn{3}{L{\dimexpr 0.63\linewidth+4\tabcolsep+2\arrayrulewidth\relax}|}{v1 2026-11-14 (178c9bc); v2 2026-11-15 (a497de6); v3 2026-11-15 (be0a468); v4 2026-11-16 (d3ab703); v5 2026-11-18 (e664d52)} \\ \hline
\textbf{Objective:} & \multicolumn{3}{L{\dimexpr 0.63\linewidth+4\tabcolsep+2\arrayrulewidth\relax}|}{Check that running several trackers over one pass of a session's frames (run\_\allowbreak{}trackers) gives each tracker exactly the tracks it gives when run alone (run\_\allowbreak{}tracker), and that the prefetch hook of beam bisection search does not change the search.} \\ \hline
\textbf{Product/Ver/ Module:} & \multicolumn{3}{L{\dimexpr 0.63\linewidth+4\tabcolsep+2\arrayrulewidth\relax}|}{FYP-1 attention pipeline, commit e743616. Tracker tuning (scripts/\allowbreak{}parallel\_\allowbreak{}runs.py): decodes each frame once and feeds it to many tracker configurations, and computes a search level's probe points together.} \\ \hline
\textbf{Environment:} & \multicolumn{3}{L{\dimexpr 0.63\linewidth+4\tabcolsep+2\arrayrulewidth\relax}|}{Windows 11 (10.0.26200); conda environment fyp-local: Python 3.13.18, NumPy 2.5.4, SciPy 1.18.1, scikit-learn 1.9.1, pandas 3.1.1, OpenCV 4.11.0, PyTorch 2.9.0+cu126. CPU only; no GPU, dataset or model weights needed unless stated. Run from the repository root: python -m unittest discover -s tests -t .} \\ \hline
\textbf{Assumptions:} & \multicolumn{3}{L{\dimexpr 0.63\linewidth+4\tabcolsep+2\arrayrulewidth\relax}|}{Synthetic detections stand in for cached detector output: 60 frames 0.2 s apart, four people (60 x 160 px boxes) walking right at 3 px per frame, Gaussian box jitter (sd 1.5 px) and confidences uniform in [0.2, 0.9], numpy seed 0. Frame size 480 x 640. All tracker specs turn CMC and ReID embeddings off, so blank frames are used and the ZIP path 'none.zip' is never opened.} \\ \hline
\textbf{Pre-Requisite:} & \multicolumn{3}{L{\dimexpr 0.63\linewidth+4\tabcolsep+2\arrayrulewidth\relax}|}{BoxMOT installed (BoT-SORT tracker). No video, GPU or detection cache needed.} \\ \hline
\textbf{Step No.} & \textbf{Execution description} & \multicolumn{2}{L{\dimexpr 0.36\linewidth+2\tabcolsep+1\arrayrulewidth\relax}|}{\textbf{Procedure result}} \\ \hline
1 & \texttt{test\_\allowbreak{}several\_\allowbreak{}trackers\_\allowbreak{}on\_\allowbreak{}one\_\allowbreak{}pass\_\allowbreak{}match\_\allowbreak{}each\_\allowbreak{}run\_\allowbreak{}alone}: Three BoT-SORT specs (use\_\allowbreak{}cmc False, use\_\allowbreak{}embeddings False): track\_\allowbreak{}high\_\allowbreak{}thresh 0.5 with det\_\allowbreak{}conf 0.3; match\_\allowbreak{}thresh 0.6 with det\_\allowbreak{}conf 0.5; max\_\allowbreak{}lost\_\allowbreak{}frames 5 with det\_\allowbreak{}conf 0.2. Call run\_\allowbreak{}trackers once with all three, then run\_\allowbreak{}tracker separately for each spec. & \multicolumn{2}{L{\dimexpr 0.36\linewidth+2\tabcolsep+1\arrayrulewidth\relax}|}{For each spec, the list of (tid, timestamps, boxes) from the shared run equals the list from the run alone. (Passed)} \\ \hline
2 & \texttt{test\_\allowbreak{}prefetch\_\allowbreak{}does\_\allowbreak{}not\_\allowbreak{}change\_\allowbreak{}the\_\allowbreak{}search}: Space: a and b in [0, 1]. Objective -(a - 0.3)\textasciicircum{}2 - (b - 0.8)\textasciicircum{}2. Run beam\_\allowbreak{}bisection\_\allowbreak{}search with a budget of 20 and baseline \{a: 0.5, b: 0.5\}, once plainly and once with prefetch collecting the requested points and the objective recording every evaluated point. & \multicolumn{2}{L{\dimexpr 0.36\linewidth+2\tabcolsep+1\arrayrulewidth\relax}|}{Both runs return the same list of params. The prefetch hook was called with at least one point, and every prefetched point was later evaluated (no prefetched work is wasted). (Passed)} \\ \hline
\multicolumn{4}{|L{\dimexpr 0.84\linewidth+6\tabcolsep+3\arrayrulewidth\relax}|}{\textbf{Comments:} This class is defined after the module's 'if \_\allowbreak{}\_\allowbreak{}name\_\allowbreak{}\_\allowbreak{} == "\_\allowbreak{}\_\allowbreak{}main\_\allowbreak{}\_\allowbreak{}": unittest.main()' block. It runs under unittest or pytest discovery, but not when the file is executed directly as a script, because unittest.main() runs before the class is defined.} \\ \hline
\multicolumn{4}{|C{\dimexpr 0.84\linewidth+6\tabcolsep+3\arrayrulewidth\relax}|}{\textbf{\fbox{Passed} \quad Failed \quad Not Executed}} \\ \hline
\end{longtable}
